\documentclass[11pt,a4paper]{article}
\usepackage[T1]{fontenc}
\usepackage[utf8]{inputenc}
\usepackage{lmodern}
\usepackage[margin=26mm,headheight=15pt]{geometry}
\usepackage{amsmath,amssymb,amsthm,mathtools,booktabs,longtable,array}
\usepackage{microtype,enumitem,fancyhdr,listings,xcolor}
\usepackage[colorlinks=true,linkcolor=blue!40!black,citecolor=blue!40!black,urlcolor=blue!40!black]{hyperref}
\pagestyle{fancy}\fancyhf{}
\fancyhead[L]{\small Keio / Bessel periods and exact orders}
\fancyhead[R]{\small Research manuscript}
\fancyfoot[C]{\thepage}
\setlength{\emergencystretch}{3em}
\setcounter{tocdepth}{1}
\setlist{itemsep=3pt,topsep=5pt}
\newtheorem{theorem}{Theorem}[section]
\newtheorem{proposition}[theorem]{Proposition}
\newtheorem{lemma}[theorem]{Lemma}
\newtheorem{corollary}[theorem]{Corollary}
\newtheorem{conjecture}[theorem]{Conjecture}
\newtheorem{criterion}[theorem]{Criterion}
\newtheorem*{theoremA}{Theorem A (banana order)}
\newtheorem*{corollaryB}{Corollary A.1}
\newtheorem*{corollaryC}{Corollary A.2}
\newtheorem*{theoremB}{Theorem B (singular locus)}
\newtheorem*{theoremC}{Theorem C (cut Galois group)}
\newtheorem*{corollaryD}{Corollary D}
\newtheorem*{corollaryE}{Corollary E}
\newtheorem*{theoremF}{Theorem F (full-integral order)}
\newtheorem*{theoremG}{Theorem G (affine extension)}
\theoremstyle{definition}
\newtheorem{definition}[theorem]{Definition}
\newtheorem{example}[theorem]{Example}
\theoremstyle{remark}
\newtheorem{remark}[theorem]{Remark}
\numberwithin{equation}{section}
\newcommand{\C}{\mathbb{C}}
\newcommand{\Q}{\mathbb{Q}}
\newcommand{\Z}{\mathbb{Z}}
\newcommand{\R}{\mathbb{R}}
\newcommand{\fl}[1]{\lfloor #1\rfloor}
\newcommand{\ff}[2]{#1^{\underline{#2}}}
\DeclareMathOperator{\ord}{ord}
\DeclareMathOperator{\Res}{Res}
\DeclareMathOperator{\Sol}{Sol}
\DeclareMathOperator{\sgn}{sgn}
\DeclareMathOperator{\SL}{SL}
\DeclareMathOperator{\Gal}{Gal}
\DeclareMathOperator{\spn}{span}
\newcommand{\N}{\mathbb{N}}
\newcommand{\Gm}{\mathbb{G}_{\mathrm{m}}}
\newcommand{\Kt}{\mathbb{K}}
\newcommand{\A}{\mathbb{A}}
\newcommand{\cD}{\mathcal{D}}
\newcommand{\cL}{\mathcal{L}}
\newcommand{\cN}{\mathcal{N}}
\newcommand{\cM}{\mathcal{M}}
\newcommand{\cO}{\mathcal{O}}
\newcommand{\Four}{\mathrm{F}}
\newcommand{\sltwo}{\mathfrak{sl}_2}
\newcommand{\varpibar}{\varpi}
\DeclareMathOperator{\rk}{rk}
\DeclareMathOperator{\Sym}{Sym}
\DeclareMathOperator{\im}{im}
\DeclareMathOperator{\irr}{irr}
\newcommand{\PP}{\mathbb{P}}
\newcommand{\cK}{\mathcal{K}}
\newcommand{\Mhat}{\widehat{\cM}}
\DeclareMathOperator{\GL}{GL}
\DeclareMathOperator{\SO}{SO}
\DeclareMathOperator{\Ort}{O}
\DeclareMathOperator{\Sp}{Sp}
\DeclareMathOperator{\Cat}{Cat}
\DeclareMathOperator{\trdeg}{trdeg}
\DeclareMathOperator{\Lie}{Lie}
\providecommand{\Ga}{\mathbb G_a}
\providecommand{\CT}{\operatorname{CT}}
\providecommand{\vp}{v_p}
\lstset{basicstyle=\ttfamily\footnotesize,breaklines=true,columns=fullflexible,keepspaces=true,showstringspaces=false,frame=single,rulecolor=\color{black!20}}
\title{Minimal Differential and Recurrence Orders\\[3pt]
\large Bessel Periods, Bounded-Height Tableaux, and Banana Integrals}
\author{Keio}
\hypersetup{pdftitle={Minimal Differential and Recurrence Orders: Bessel Periods, Bounded-Height Tableaux, and Banana Integrals},pdfauthor={Keio},pdfsubject={Integrated and extended research manuscript},pdfkeywords={Bessel periods, Young tableaux, banana integrals, differential Galois groups, recurrence orders}}
\date{18 October 2026\\\small Integrated and extended research manuscript}
\begin{document}
\maketitle
\begin{abstract}
We study exact differential and recurrence orders for two families built from
Bessel functions: the exponential generating functions of standard Young
tableaux of bounded height, and the periods of two-dimensional banana graphs.
For tableaux we determine the minimal operator, its polynomial degree and its
local exponents. The existence and coefficient-shape assertion is also derived
directly from the published Jacobi integral system of Forrester and Kumar;
this comparison separates that consequence of prior work from the minimality
and singularity analysis. For banana periods, arbitrary nonzero mass
multiplicities determine the exact differential order through a tensor product
of $\mathfrak{sl}_2$ representations. Local Fourier transforms determine the
threshold monodromy, and the cut Galois groups are orthogonal or symplectic
under a stated noncollision hypothesis.

Further deductions connect these results with two methods used for Franel
sequences: a singularity bound on recurrence width and a discrete Abel
identity. We obtain exact coefficient-recurrence widths under the noncollision
hypothesis, strict order increases under splitting repeated squared masses,
and integral-solution lattices for explicit low-height and two-mass
recurrences. For all positive masses the full Feynman integral has an affine
differential Galois group over its homogeneous cut group; its first $o$
derivatives are algebraically independent over the cut Picard--Vessiot field.
These are functional statements, with no inference about irrationality of
special numerical values. All proofs and the bibliography are contained in
this single manuscript.
\end{abstract}

\noindent\textbf{Keywords.} Bessel functions; minimal differential order;
recurrence width; Young tableaux; banana integrals; differential Galois groups;
integral solution lattices.

\medskip
\noindent\textbf{Manuscript status.} This is an integrated research draft,
prepared with AI assistance and not externally peer reviewed. It revises three
supplied manuscripts and adds explicitly proved consequences. A newly included
deduction is not, by that fact alone, a claim of historical priority or the
solution of a separately published open problem. The exact scope of the
Lairez--Vanhove and rational-diagonal conjectures is stated below. Computations
are supporting checks and do not replace the analytic proofs.

\clearpage
\tableofcontents
\clearpage
\section{A common problem: exact complexity of period functions}
\label{sec:overview}

Three constructions organize the paper. A Bessel determinant becomes a
Jacobi integral for tableaux. A Bessel product becomes a banana period after
two Fourier--Laplace transforms and an inversion. A differential relation for
a generating function becomes a recurrence for its coefficients. These
constructions provide upper bounds; matching lower bounds require additional
information about solution spaces, singularities, and arithmetic.

The supplied Franel manuscript \cite{KeioFranel} suggests two particularly
useful principles. First, a recurrence with shifts spanning an interval of
length $r$ permits at most $r$ distinct nonzero finite singularities of its
ordinary generating function. Second, a discrete Abel identity can force
every integral solution of a second-order recurrence onto one rational line.
Both principles are proved again where they are used here. We do not require
the Franel paper's general minimality theorem as an input. In particular, its
definite-critical-point criterion is not applied to banana saddle points;
the latter are handled by local Fourier theory and irreducibility.

\subsection{The invariants must be kept separate}
For a function $f(t)$, differential order means the smallest $o$ such that
$\sum_{j=0}^{o}q_j(t)f^{(j)}(t)=0$, with rational coefficients and $q_o\ne0$.
For a sequence $a_k$, recurrence width means the smallest difference between
the largest and smallest shifts in a relation
$\sum_{j=0}^{r}p_j(k)a_{k+j}=0$ valid for all sufficiently large $k$.
For polynomial $p_j$, the maximum of their degrees is a third invariant.
Minimizing the width and minimizing the coefficient degree are separate
optimization problems; the same recurrence need not attain both minima.

For the banana coefficients and periods, write
\begin{gather*}
 c_k(m)=\sum_{a_1+\cdots+a_n=k}\binom{k}{a_1,\ldots,a_n}^{\!2}
             \prod_i m_i^{2a_i},\\
 F_m(z)=\sum_{k\ge0}c_k(m)z^k,\qquad
 \varpi_m(t)=t^{-1}F_m(t^{-1}).
\end{gather*}
Let the distinct squared masses have multiplicities
$\boldsymbol n=(n_1,\ldots,n_s)$, let
$D=\prod_j(n_j+1)$, and let $\delta=1$ if every $n_j$ is even and $0$ otherwise.
The nonzero signed weight sums, modulo simultaneous change of sign, give
$(D-\delta)/2$ candidate thresholds. The hypothesis (G$'$) requires their
squares to be nonzero and pairwise distinct; it is defined precisely in
Part~\ref{part:structure}.

\begin{center}
\begin{tabular}{@{}p{.27\textwidth}p{.31\textwidth}p{.29\textwidth}@{}}
\toprule
Quantity & Pairwise distinct masses under (G) & Equal nonzero masses\\
\midrule
Minimal cut differential order &
$2^n-\binom{n+1}{\lfloor(n+1)/2\rfloor}$ & $n-1$\\[4pt]
Minimal coefficient recurrence width & $2^{n-1}$ &
$\lfloor(n+1)/2\rfloor$\\[4pt]
Minimal full-integral homogeneous order, positive masses &
$2^n-\binom{n+1}{\lfloor(n+1)/2\rfloor}+1$ & $n$\\
\bottomrule
\end{tabular}
\end{center}
Here $n\ge2$ and ``distinct masses'' in the table means distinct squared
masses. The differential-order theorem itself needs no threshold-noncollision
assumption and applies to all nonzero complex masses. The recurrence-width
formula is established under (G$'$); collisions of different signed sums are
not silently included in that assertion.

\subsection{Scope, provenance, and additional deductions}
Part~\ref{part:tableaux} retains the full tableau argument and adds an explicit
comparison with the Jacobi differential system of Forrester--Kumar. Thus the
BFK existence assertion is presented with its prior-art reduction, rather than
as an unqualified claim of a first solution. The independent contribution of
the proof is its exact minimality and local-exponent mechanism. The case of
height one is treated separately.

Part~\ref{part:order} gives upper and lower bounds for the actual banana period,
including the boundary terms in the Laplace transforms. Part~\ref{part:structure}
then gives its singular locus and Galois group, proves the nonzero rational
inhomogeneity for the full integral, and develops the affine extension.
The remaining part transfers these results to coefficient recurrences and
arithmetic solution lattices. The mass-splitting result identifies the unique
extremal multiplicity partitions for $n\ge4$.

The following limitations are part of the mathematical scope. The simple-root
claim for thresholds in the leading coefficient of the \emph{minimal-order}
Picard--Fuchs operator, and the conjectured multiplicity at zero, are not proved.
An operator of larger order constructed from an optimal-width recurrence may
have simple thresholds without settling that question. Likewise, the
orthogonal/symplectic classification for diamond lattices does not prove the
conjectured minimum number of variables in every rational-diagonal
representation. No theorem here asserts the irrationality or transcendence
of the Euler--Mascheroni or Gompertz constants.

The large finite computations described in the source drafts are identified
as source-reported checks where their certificates were not supplied. The
additional exact checks made for this compilation are described in the
verification appendix. Their role is to catch normalization and index errors;
the unbounded-parameter claims rest on the proofs.

\clearpage
\part{Bounded-height tableaux and Jacobi integrals}\label{part:tableaux}


\section{Bounded-height tableaux and the BFK assertion}

For integers $k\ge 1$ and $n\ge 0$ let $y_k(n)$ be the number of standard Young tableaux with $n$ cells and at most $k$ rows. By the Robinson--Schensted correspondence, $y_k(n)$ is also the number of involutions of $\{1,\dots,n\}$ with no decreasing subsequence of length $k+1$. Put
\[
Y_k(t)=\sum_{n\ge 0}y_k(n)\,\frac{t^n}{n!}.
\]
Gessel \cite{Ges90} showed that $Y_k$ is D-finite, that is, it satisfies a linear differential equation with polynomial coefficients. Bergeron, Favreau and Krob \cite{BFK95} computed such equations for small $k$ and conjectured their general shape. We use the formulation of Bergeron and Gascon \cite{BG00}.

\begin{conjecture}[Bergeron--Favreau--Krob \cite{BFK95,BG00}]\label{bfk:conj:BFK}
Let $\ell=\ell(k)=\fl{k/2}+1$. For each $k$ there are polynomials $p_0,\dots,p_\ell$ of degree at most $\ell-1$ such that
\[
\sum_{m=0}^{\ell}p_m(t)\,Y_k^{(m)}(t)=0,
\]
with $p_\ell(t)=t^{\ell-1}$ and $p_m(t)=t^{m-1}q_m(t)$ for some polynomials $q_m$, $1\le m\le \ell$.
\end{conjecture}

Bergeron and Gascon proved that $Y_k$ and its derivatives span a vector space of dimension at most $\ell$ over $\C(t)$, so that an equation of order at most $\ell$ exists, and they verified the full conjecture for $k\le 11$ by computer algebra \cite{BG00}. The weaker statement (order at most $\fl{k/2}+1$ and degree at most $\fl{k/2}$) is Conjecture~9 in Mishna's survey \cite{Mis19}, where it is listed as open; the general-purpose bounds recorded there are exponential in $k$.

Our main theorem establishes the stated coefficient shape and identifies the
minimal operator. A precise prior-art distinction is essential. The
Forrester--Kumar system for Jacobi trace transforms \cite{FK22}, already
available as a 2020 preprint, yields the existence and coefficient-shape part
by a direct elimination argument. Section~\ref{bfk:sec:bfk-jacobi} supplies that
argument, with the exact parameter substitution. We do not claim that those
authors explicitly stated the BFK conclusion, nor claim first priority for
this consequence. The proof below additionally determines exact minimality
and local exponents and explains the absence of apparent singularities.

\begin{theorem}\label{bfk:thm:main}
Let $k\ge 1$, $r=\fl{k/2}$, $\varepsilon=k-2r\in\{0,1\}$, $n=r+1$ and $a=\varepsilon-\tfrac12$. Let
\[
L_k=\sum_{m=0}^{n'}p_m(t)\,\partial^m,\qquad \partial=\frac{d}{dt},
\]
be the linear differential operator of minimal order $n'$ annihilating $Y_k$, normalized so that the polynomials $p_m$ have no common zero. Then the following hold.
\begin{enumerate}
\item[(i)] $n'=n=\fl{k/2}+1$.
\item[(ii)] After multiplication by a nonzero constant, $p_n(t)=t^{r}$; moreover $\deg p_m\le r$ for all $m$, and $t^{m-1}$ divides $p_m$ for $1\le m\le n$. Every operator of order $n$ with polynomial coefficients annihilating $Y_k$ is a polynomial multiple of $L_k$; in particular the degree bound $r$ cannot be lowered.
\item[(iii)] The point $t=0$ is regular singular for $k\ge2$ and ordinary for $k=1$, with exponents
\[
0,\ e_1,\ \dots,\ e_r,\qquad e_b=b^2-2br-\varepsilon b .
\]
At $t=\infty$ the operator $L_k$ has a basis of formal solutions $e^{\lambda_jt}t^{-\gamma_j}\bigl(1+\sum_{i\ge1}c_{j,i}t^{-i}\bigr)$, $0\le j\le r$, where
\[
\lambda_j=\varepsilon+2r-4j,\qquad \gamma_j=(r-j)(r-j+a)+j\bigl(j+\tfrac12\bigr).
\]
\item[(iv)] $L_k$ has no singular point in $\C\setminus\{0\}$.
\end{enumerate}
\end{theorem}

\begin{corollary}\label{bfk:cor:BFK}
Conjecture~\ref{bfk:conj:BFK} holds for every $k\ge 1$. In particular Conjecture~9 of \cite{Mis19} holds.
\end{corollary}

\begin{corollary}\label{bfk:cor:rec}
For every $k\ge1$ there are polynomials $c_0,\dots,c_{r+1}\in\Q[N]$ of degree at most $r=\fl{k/2}$, with $\deg c_0=r$, such that
\[
\sum_{i=0}^{r+1}c_i(N)\,y_k(N-i)=0\qquad (N\ge 1),
\]
where $y_k(m)=0$ for $m<0$.
\end{corollary}

\subsection*{The case $k=1$}
Here $Y_1(t)=e^t$ and the primitive minimal operator is $L_1=\partial-1$.
All assertions of Theorem~\ref{bfk:thm:main} and the recurrence corollary follow
directly, with $r=0$, exponent $0$ at the ordinary point $0$, and formal
solution $e^t$ at infinity. In the Galois-orbit and asymptotic proof below we
therefore assume $k\ge2$, so that $r\ge1$; the polarization identity used there
is only invoked for positive $r$.

\subsection*{Strategy}
The Bessel determinants of Gordon and Gessel (Theorem~\ref{bfk:thm:det}) can be rewritten, via Andreief's identity, as a Laplace transform
\[
Y_k(t)=\frac{e^{\varepsilon t}}{r!\,(2\pi)^r}\int_{[-2,2]^r}e^{t(u_1+\dots+u_r)}\,\Delta(u)^2\prod_{l=1}^r w(u_l)\,du ,
\]
a $\beta=2$ Jacobi ensemble with weight $w(u)=(2-u)^{a}(2+u)^{1/2}$ (Proposition~\ref{bfk:prop:int}). Replacing $b$ of the $r$ segments $[-2,2]$ by the ray $(-\infty,-2]$ produces integrals $\Psi_b$, $0\le b\le r$ (with $\Psi_0$ the original one up to the factor $e^{\varepsilon t}$), and the differential Galois group $\SL_2$ of Bessel's equation shows that they span the solution space of the minimal operator (Proposition~\ref{bfk:prop:V}). Laplace's method at the vertices of the cube gives the exponential types and the exponents at infinity (Lemma~\ref{bfk:lem:vertex}), and a scaling argument gives the exponents at the origin (Lemma~\ref{bfk:lem:origin}). The two lists of exponents satisfy
\[
\sum_{j=0}^{r}\gamma_j+\sum_{b=0}^{r}e_b=\binom{n}{2},
\]
and Abel's identity for the Wronskian turns this into the statement that the minimal operator has no singular point other than $0$ (Section~\ref{bfk:sec:proof}). The shape of the operator then follows from the regularity at $0$, the exponents, and a top-degree argument at infinity.

The order $n=\fl{k/2}+1$ appears as the number of $S_r$-orbits of vertices of $[-2,2]^r$, that is, the number of distinct values $2r-4j$ $(0\le j\le r)$ of $u_1+\dots+u_r$ at the vertices.

\section{Three facts about linear differential operators}\label{bfk:sec:ode}

Throughout, $K=\C(t)$, $\partial=d/dt$, $\theta=t\partial$, and $\ff{x}{m}=x(x-1)\cdots(x-m+1)$. An operator $L=\sum_{m=0}^{n}p_m\partial^m$ with $p_m\in\C[t]$ and $p_n\ne0$ is called primitive if the $p_m$ have no common zero. Note that $t^m\partial^m=\ff{\theta}{m}$.

Recall Fuchs' criterion: $\tau\in\C$ is a regular singular point (or an ordinary point) of $L$ if and only if $\ord_\tau p_m\ge \ord_\tau p_n-(n-m)$ for all $m$; equivalently, all solutions have moderate growth near $\tau$ \cite{Inc56,Was65}. In that case, with $A=\ord_\tau p_n$,
\begin{equation}\label{bfk:eq:thetaform}
(t-\tau)^{n-A}L=\sum_{m=0}^{n}s_m(t)\,(t-\tau)^m\partial^m,\qquad s_m\in\C[t],\ s_n(\tau)\ne0,
\end{equation}
and the indicial polynomial is $I_\tau(\rho)=\sum_{m=0}^{n}s_m(\tau)\ff{\rho}{m}$. It has degree $n$; its roots, with multiplicity, are the exponents of $L$ at $\tau$.

\begin{lemma}\label{bfk:lem:indicial}
Suppose $0$ is a regular singular point of $L$, and let $y=\sum_{\nu\in\nu_0+\Z_{\ge0}}t^{\nu}P_\nu(\log t)$ be a solution, given by a convergent expansion on a slit disc, where the $P_\nu$ are polynomials of bounded degree and $P_{\nu_0}\ne 0$. Then $I_0(\nu_0)=0$.
\end{lemma}

\begin{proof}
Write $D=d/d\ell$ acting on polynomials in $\ell=\log t$. Then $\theta\bigl(t^\nu P(\ell)\bigr)=t^\nu(\nu+D)P(\ell)$, hence $t^m\partial^m\bigl(t^\nu P\bigr)=\ff{\theta}{m}\bigl(t^\nu P\bigr)=t^\nu\,\ff{(\nu+D)}{m}P$. Apply \eqref{bfk:eq:thetaform} with $\tau=0$ to $y$. Since $s_m(t)=s_m(0)+O(t)$, the coefficient of $t^{\nu_0}$ in $t^{n-A}Ly=0$ is $I_0(\nu_0+D)P_{\nu_0}=0$. The operator $D$ is nilpotent on polynomials of bounded degree, so $I_0(\nu_0+D)=I_0(\nu_0)+(\text{nilpotent})$. If $I_0(\nu_0)\ne0$ this operator is invertible, which forces $P_{\nu_0}=0$, a contradiction.
\end{proof}

\begin{lemma}\label{bfk:lem:abel}
Let $L$ be primitive of order $n$, and let $W$ be the Wronskian of a basis of solutions on a simply connected domain. Then $W'/W=-p_{n-1}/p_n$, and the following hold.
\begin{enumerate}
\item[(a)] If $\tau$ is a regular singular point with exponents $\rho_1,\dots,\rho_n$, then the residue of $W'/W$ at $\tau$ is $\sum_i\rho_i-\binom n2$.
\item[(b)] If $p_n(\tau)=0$ and all solutions are holomorphic at $\tau$, then $\Res_{t=\tau}(W'/W)$ is a positive integer.
\end{enumerate}
\end{lemma}

\begin{proof}
The formula $W'/W=-p_{n-1}/p_n$ is Abel's identity. For (a), \eqref{bfk:eq:thetaform} gives $p_{n-1}/p_n=s_{n-1}(t)/\bigl((t-\tau)s_n(t)\bigr)$, so the residue of $-p_{n-1}/p_n$ at $\tau$ is $-s_{n-1}(\tau)/s_n(\tau)$. On the other hand $\ff{\rho}{n}=\rho^n-\binom n2\rho^{n-1}+\cdots$ and $\ff{\rho}{n-1}=\rho^{n-1}+\cdots$, so
\[
\frac{I_\tau(\rho)}{s_n(\tau)}=\rho^n-\Bigl(\binom n2-\frac{s_{n-1}(\tau)}{s_n(\tau)}\Bigr)\rho^{n-1}+\cdots,
\]
and $\sum_i\rho_i=\binom n2-s_{n-1}(\tau)/s_n(\tau)$.

For (b), $W$ is holomorphic at $\tau$, so the residue equals $\ord_\tau W\ge0$. If $W(\tau)\ne0$, Cramer's rule applied to the linear system satisfied by the coefficients shows that every quotient $p_m/p_n$ is holomorphic at $\tau$. Then $(t-\tau)$ divides every $p_m$, contradicting primitivity. Hence ${\ord_\tau W\ge1}$.
\end{proof}

\begin{lemma}\label{bfk:lem:top}
Let $L=\sum_{m=0}^np_m\partial^m$, let $d=\max_m\deg p_m$, and put $\chi(\lambda)=\sum_{m}[t^d]p_m\,\lambda^m$. Suppose there are solutions $y_1,\dots,y_n$, pairwise distinct numbers $\lambda_1,\dots,\lambda_n$, constants $C_j\ne0$ and $\gamma_j$ such that
\[
y_j^{(m)}(t)=C_je^{\lambda_jt}t^{-\gamma_j}\bigl(\lambda_j^m+O(t^{-1})\bigr)\qquad(t\to+\infty,\ 0\le m\le n).
\]
Then $\chi(\lambda_j)=0$ for all $j$. Consequently $\deg p_n=d$ and $\chi(\lambda)=[t^d]p_n\prod_j(\lambda-\lambda_j)$.
\end{lemma}

\begin{proof}
We have $0=e^{-\lambda_jt}t^{\gamma_j}C_j^{-1}Ly_j=\sum_mp_m(t)\bigl(\lambda_j^m+O(t^{-1})\bigr)=t^d\chi(\lambda_j)+O(t^{d-1})$. Hence $\chi$ is a nonzero polynomial of degree at most $n$ with the $n$ distinct roots $\lambda_j$. Therefore its degree is $n$, so $[t^d]p_n\ne0$.
\end{proof}

\section{An integral representation}\label{bfk:sec:int}

For $m\in\Z$ let $b_m(t)=I_m(2t)=\sum_{j\ge0}t^{2j+|m|}/\bigl(j!\,(j+|m|)!\bigr)$, where $I_m$ is the modified Bessel function. Recall $b_m(t)=\frac1\pi\int_0^\pi e^{2t\cos\theta}\cos(m\theta)\,d\theta$ \cite[\S10.32]{DLMF}.

\begin{theorem}[Gordon; Gessel]\label{bfk:thm:det}
For $r\ge0$,
\[
Y_{2r}(t)=\det\bigl(b_{i-j}+b_{i+j-1}\bigr)_{1\le i,j\le r},\qquad
Y_{2r+1}(t)=e^{t}\det\bigl(b_{i-j}-b_{i+j}\bigr)_{1\le i,j\le r},
\]
the empty determinant being $1$.
\end{theorem}

Both formulas are given by Gessel \cite{Ges90}, who deduced them from a result of Gordon; see also \cite{BG00} and, for the history, \cite{Mis19}.

Let $\Delta(u)=\prod_{1\le l<l'\le r}(u_{l'}-u_l)$ and, for $\varepsilon\in\{0,1\}$ and $a=\varepsilon-\frac12$,
\[
w(u)=(2-u)^{a}(2+u)^{1/2},\qquad -2<u<2;
\]
thus $w(u)=\sqrt{(2+u)/(2-u)}$ if $\varepsilon=0$ and $w(u)=\sqrt{4-u^2}$ if $\varepsilon=1$.

\begin{proposition}\label{bfk:prop:int}
Let $k=2r+\varepsilon$ with $\varepsilon\in\{0,1\}$. Then
\[
Y_k(t)=\frac{e^{\varepsilon t}}{r!\,(2\pi)^r}\int_{[-2,2]^r}e^{t(u_1+\cdots+u_r)}\,\Delta(u)^2\prod_{l=1}^rw(u_l)\,du .
\]
\end{proposition}

\begin{proof}
Let $\varepsilon=0$. From the integral for $b_m$ and $\cos((i-j)\theta)+\cos((i+j-1)\theta)=2\cos((i-\frac12)\theta)\cos((j-\frac12)\theta)$ we get
\[
b_{i-j}+b_{i+j-1}=\frac2\pi\int_0^\pi e^{2t\cos\theta}\cos\bigl((i-\tfrac12)\theta\bigr)\cos\bigl((j-\tfrac12)\theta\bigr)\,d\theta .
\]
Now $\cos((i-\frac12)\theta)=\cos(\theta/2)\,P_i(2\cos\theta)$, where $P_i$ is a monic polynomial of degree $i-1$ (a rescaled Chebyshev polynomial of the third kind). Substituting $u=2\cos\theta$, $d\theta=-du/\sqrt{4-u^2}$ and $\cos^2(\theta/2)=(2+u)/4$ gives
\[
b_{i-j}+b_{i+j-1}=\frac1{2\pi}\int_{-2}^2e^{tu}P_i(u)P_j(u)\,w(u)\,du .
\]
Andreief's identity (see \cite{For10}), $\det\bigl(\int f_if_j\,d\mu\bigr)_{i,j=1}^r=\frac1{r!}\int\det(f_i(u_l))^2\prod_ld\mu(u_l)$, together with $\det(P_i(u_l))_{i,l}=\Delta(u)$, gives the claim by Theorem~\ref{bfk:thm:det}. For $\varepsilon=1$ use instead $\cos((i-j)\theta)-\cos((i+j)\theta)=2\sin(i\theta)\sin(j\theta)$, $\sin(i\theta)=\sin\theta\,P_i(2\cos\theta)$ with $P_i$ monic of degree $i-1$ (Chebyshev polynomials of the second kind), and $\sin^2\theta=(4-u^2)/4$.
\end{proof}

\section{The one-variable equation and the solution space}\label{bfk:sec:V}

Fix $k=2r+\varepsilon\ge2$ and $a=\varepsilon-\frac12$, so $r\ge1$. Put $\tilde w(u)=(2-u)^{a}|2+u|^{1/2}$ for $u<-2$, and define, for $\Re t>0$,
\[
\varphi_I(t)=\int_{-2}^{2}e^{tu}w(u)\,du,\qquad \varphi_K(t)=\int_{-\infty}^{-2}e^{tu}\,\tilde w(u)\,du .
\]
The function $\varphi_I$ is entire, and $\varphi_K$ extends analytically to the universal cover of $\C^*$ (rotate the ray of integration).

\begin{lemma}\label{bfk:lem:B}
Both $\varphi_I$ and $\varphi_K$ satisfy
\begin{equation}\label{bfk:eq:B}
B_\varepsilon f:=tf''+(2+\varepsilon)f'-(4t+2-2\varepsilon)f=0 .
\end{equation}
This equation has singular points only at $0$ and $\infty$, and its exponents at $0$ are $0$ and $-1-\varepsilon$. The functions $\varphi_I,\varphi_K$ are linearly independent, so $\Sol(B_\varepsilon)=\C\varphi_I\oplus\C\varphi_K$.
\end{lemma}

\begin{proof}
Let $g(u)=(2-u)^{a+1}(2+u)^{3/2}=(4-u^2)w(u)$ on $(-2,2)$. Then
\[
g'(u)=w(u)\bigl[(1-2a)-(a+\tfrac52)u\bigr]=w(u)\bigl[(2-2\varepsilon)-(2+\varepsilon)u\bigr].
\]
Since $g(\pm2)=0$, we get $0=\int_{-2}^2\bigl(e^{tu}g\bigr)'du=\int_{-2}^2e^{tu}\bigl[t(4-u^2)+(2-2\varepsilon)-(2+\varepsilon)u\bigr]w\,du$. Since $\int e^{tu}u^mw\,du=\varphi_I^{(m)}$, this is \eqref{bfk:eq:B}. On $(-\infty,-2)$ the same computation applies to $\tilde g=(2-u)^{a+1}|2+u|^{3/2}=-(4-u^2)\tilde w$, which vanishes at $-2$; the boundary term at $-\infty$ vanishes for $\Re t>0$. The indicial polynomial at $0$ is $\rho(\rho-1)+(2+\varepsilon)\rho=\rho(\rho+1+\varepsilon)$. Finally, $\varphi_I(0)=\int w>0$, while $\varphi_K(t)\to+\infty$ as $t\to0^+$ because $\tilde w(u)\sim|u|^{\varepsilon}$ is not integrable at $-\infty$. So the two functions are independent.
\end{proof}

\begin{lemma}\label{bfk:lem:galois}
Let $E$ be the Picard--Vessiot extension of $K=\C(t)$ for \eqref{bfk:eq:B} and $G=\Gal(E/K)$. Then $G$ acts on $\Sol(B_\varepsilon)$ as the full group $\SL(\Sol(B_\varepsilon))\cong\SL_2(\C)$. In particular $G$ acts transitively on $\Sol(B_\varepsilon)\setminus\{0\}$.
\end{lemma}

\begin{proof}
Let $T=t\partial^2+\partial-4t$; the function $f_0=I_0(2t)$ satisfies $Tf_0=0$, the modified Bessel equation of order $0$ in the variable $2t$. Put $G_0=1+\frac12\partial$ and $G_1=(2t)^{-1}\partial$. A direct computation gives the operator identities
\[
B_0\circ G_0=\bigl(\tfrac12\partial+1\bigr)\circ T,\qquad B_1\circ G_1=\bigl(\partial\circ(2t)^{-1}\bigr)\circ T .
\]
Hence $G_\varepsilon$ maps $\Sol(T)$ to $\Sol(B_\varepsilon)$. It is injective there: $G_0f=0$ forces $f=Ce^{-2t}$, $G_1f=0$ forces $f$ to be constant, and neither is a nonzero solution of $T$. So $G_\varepsilon$ is a $K$-gauge equivalence, and the two equations have the same Picard--Vessiot extension. (Concretely, $\varphi_I=2\pi G_\varepsilon f_0$, that is, $\varphi_I=2\pi(I_0(2t)+I_1(2t))$ for $\varepsilon=0$ and $\varphi_I=2\pi I_1(2t)/t$ for $\varepsilon=1$.) The differential Galois group of the Bessel equation of integer order is $\SL_2(\C)$ \cite{Kol68}; see also \cite{vdPS03}. The modified equation is obtained by the linear substitution $t\mapsto it$, which does not change the Galois group. The Wronskian of \eqref{bfk:eq:B} is $t^{-2-\varepsilon}\in K$, consistent with $G\subseteq\SL$.
\end{proof}

For $g_1,\dots,g_r\in\Sol(B_\varepsilon)$ define
\begin{equation}\label{bfk:eq:F}
F(g_1,\dots,g_r)=\sum_{\sigma,\tau\in S_r}\sgn(\sigma)\sgn(\tau)\prod_{l=1}^rg_l^{(\sigma(l)+\tau(l)-2)} .
\end{equation}

\begin{lemma}\label{bfk:lem:F}
\begin{enumerate}
\item[(a)] $F$ is $\C$-multilinear and symmetric in its arguments.
\item[(b)] If $g_l=\int_{\gamma_l}e^{tu}\omega_l(u)\,du$ with $(\gamma_l,\omega_l)\in\{([-2,2],w),((-\infty,-2],\tilde w)\}$, then
\[
F(g_1,\dots,g_r)=\int_{\gamma_1\times\cdots\times\gamma_r}e^{t(u_1+\cdots+u_r)}\Delta(u)^2\prod_{l=1}^r\omega_l(u_l)\,du .
\]
\item[(c)] $Z:=F(\varphi_I,\dots,\varphi_I)=r!\,(2\pi)^re^{-\varepsilon t}Y_k(t)$.
\end{enumerate}
\end{lemma}

\begin{proof}
Multilinearity is clear. For symmetry, permute the arguments and compose $\sigma$ and $\tau$ with the same permutation; the two signs cancel. For (b), expand
\[
\Delta(u)^2=\det\bigl(u_l^{\,i-1}\bigr)^2=\sum_{\sigma,\tau}\sgn(\sigma\tau)\prod_lu_l^{\sigma(l)+\tau(l)-2},
\]
use Fubini, and use $\int_\gamma e^{tu}u^m\omega(u)\,du=g^{(m)}(t)$. Part (c) follows from (b) and Proposition~\ref{bfk:prop:int}.
\end{proof}

For $0\le b\le r$ set
\[
\Psi_b=F(\underbrace{\varphi_K,\dots,\varphi_K}_{b},\underbrace{\varphi_I,\dots,\varphi_I}_{r-b})
=\int_{(-\infty,-2]^b\times[-2,2]^{r-b}}e^{t\sum_lu_l}\Delta(u)^2\prod_{l\le b}\tilde w(u_l)\prod_{l>b}w(u_l)\,du ,
\]
and let $V=\spn_\C\{F(g_1,\dots,g_r):g_l\in\Sol(B_\varepsilon)\}$. By Lemma~\ref{bfk:lem:F}(a), expanding each $g_l$ in the basis $\varphi_I,\varphi_K$ gives
\begin{equation}\label{bfk:eq:Vspan}
V=\spn_\C\{\Psi_0,\dots,\Psi_r\},\qquad \dim V\le r+1 .
\end{equation}

\begin{proposition}\label{bfk:prop:V}
Let $L_Z$ be the monic operator over $K$ of minimal order annihilating $Z$. Then $\Sol(L_Z)=V$. Consequently the minimal order of $Y_k$ equals $\dim V$, and the solution space of the minimal operator of $Y_k$ is $e^{\varepsilon t}V$.
\end{proposition}

\begin{proof}
Since $F$ has integer coefficients and involves only derivatives, $Z\in E$ and $\sigma(Z)=F(\sigma\varphi_I,\dots,\sigma\varphi_I)$ for $\sigma\in G$. Let $U=\spn_\C\{\sigma(Z):\sigma\in G\}$. By Lemma~\ref{bfk:lem:galois}, $\{\sigma\varphi_I\}=\Sol(B_\varepsilon)\setminus\{0\}$, so $U$ contains every $F(g,\dots,g)$ with $g\in\Sol(B_\varepsilon)$. Since $F$ is symmetric and multilinear, the polarization identity
\[
F(g_1,\dots,g_r)=\frac1{r!}\sum_{\emptyset\ne S\subseteq\{1,\dots,r\}}(-1)^{r-|S|}F(g_S,\dots,g_S),\qquad g_S=\sum_{l\in S}g_l ,
\]
shows $V\subseteq U$, and clearly $U\subseteq V$. So $U=V$.

The space $U$ is finite-dimensional and $G$-stable. The monic operator $L_U$ with solution space $U$ (built from Wronskians of a basis) is therefore fixed by $G$, so it has coefficients in $K$ by the Galois correspondence \cite{vdPS03}. Since $L_UZ=0$, we get $\ord L_Z\le\dim U$. Conversely, $L_Z$ has coefficients in $K$, so it annihilates every $\sigma(Z)$; hence $U\subseteq\Sol(L_Z)$ and $\dim U\le\ord L_Z$. Thus $\Sol(L_Z)=U=V$. The last assertion follows from Lemma~\ref{bfk:lem:F}(c), because $L_Z(\partial-\varepsilon)$ is the minimal operator of $e^{\varepsilon t}Z$.
\end{proof}

\section{Asymptotics at the vertices and at the origin}\label{bfk:sec:asym}

For an integer $m\ge0$ and $c>-1$ put
\[
S_m(c)=\int_{(0,\infty)^m}e^{-(x_1+\cdots+x_m)}\,\Delta(x)^2\prod_{i=1}^mx_i^{c}\,dx\in(0,\infty),\qquad S_0(c)=1 .
\]
These are Laguerre--Selberg integrals \cite{FW08}; we only use that they are finite and positive.

\begin{lemma}\label{bfk:lem:vertex}
Let $0\le b\le r$ and $\mu_b=2r-4b$, and let $\gamma_b$ be as in Theorem~\ref{bfk:thm:main}(iii). Then, as $t\to+\infty$, for every $m\ge0$,
\[
\Psi_b^{(m)}(t)=C_b\,e^{\mu_bt}t^{-\gamma_b}\bigl(\mu_b^m+O(t^{-1})\bigr),\qquad
C_b=4^{2b(r-b)+(r-b)/2+ab}\,S_{r-b}(a)\,S_b(\tfrac12)>0 .
\]
Moreover $\Psi_b$ has an asymptotic expansion $e^{\mu_bt}t^{-\gamma_b}\sum_{i\ge0}c_{b,i}t^{-i}$ with $c_{b,0}=C_b$, and the asymptotic expansions of the derivatives $\Psi_b^{(m)}$ are obtained by termwise differentiation.
\end{lemma}

\begin{proof}
We have $\Psi_b^{(m)}=\int(\textstyle\sum_lu_l)^me^{t\sum u_l}\Delta^2\prod\tilde w\prod w\,du$ over $(-\infty,-2]^b\times[-2,2]^{r-b}$. On this domain $\sum_lu_l$ attains its maximum $\mu_b$ only at the vertex $v^*$ whose first $b$ coordinates are $-2$ and whose last $r-b$ coordinates are $2$. Outside a neighbourhood of $v^*$ we have $\sum u_l\le\mu_b-\delta$ for some $\delta>0$. For $t\ge1$ write $e^{t\sum u}=e^{\sum u}e^{(t-1)\sum u}$; since $e^{\sum u}(\sum u)^m\Delta^2\prod\tilde w\prod w$ is integrable, the contribution of that region is $O(e^{(\mu_b-\delta)t})$.

Near $v^*$ use local coordinates $y_l=-2-u_l\ge0$ for $l\le b$ and $x_l=2-u_l\ge0$ for $l>b$. Then $\sum u_l=\mu_b-\sum x-\sum y$ and the integrand becomes
\[
e^{\mu_bt}e^{-t(\sum x+\sum y)}\,\Phi(x,y)\,H(x,y)\,\bigl(\mu_b-\textstyle\sum x-\sum y\bigr)^m,
\qquad \Phi=\Delta(x)^2\Delta(y)^2\prod x_l^{a}\prod y_l^{1/2},
\]
where $H(x,y)=\prod_{l\le b<l'}(u_{l'}-u_l)^2\prod_{l>b}(2+u_l)^{1/2}\prod_{l\le b}(2-u_l)^{a}$ is real-analytic near $0$ with $H(0,0)=4^{2b(r-b)}4^{(r-b)/2}4^{ab}$. The function $\Phi$ is positively homogeneous of degree
\[
(r-b)(r-b-1)+b(b-1)+a(r-b)+\tfrac b2=\gamma_b-r .
\]
By the multidimensional Watson lemma (Taylor-expand $H\cdot(\mu_b-\sum x-\sum y)^m$ at $0$ and integrate term by term against $e^{-t(\sum x+\sum y)}\Phi$; see \cite{Won89}), the local contribution has an expansion $e^{\mu_bt}\sum_{i\ge0}d_it^{-\gamma_b-i}$. Its leading coefficient is $H(0,0)\mu_b^m\int_{(0,\infty)^r}e^{-\sum x-\sum y}\Phi\,dx\,dy=\mu_b^mC_b$. The integrand $\Phi$ is positive, so $C_b>0$. When $\mu_b=0$ and $m\ge1$ the leading coefficient vanishes, in accordance with the statement. The expansion of $\Psi_b^{(m)}$ obtained this way is the termwise derivative of that of $\Psi_b$, since differentiating in $t$ under the integral sign multiplies the integrand by $\sum u_l$.
\end{proof}

\begin{lemma}\label{bfk:lem:origin}
Every element of $V$ has, on a slit neighbourhood of $0$, a convergent expansion $\sum_{\nu\ge\nu_0}t^\nu P_\nu(\log t)$ with $\nu\in\Z$ and polynomials $P_\nu$ of bounded degree. For $0\le b\le r$,
\[
\lim_{t\to0^+}t^{-e_b}\Psi_b(t)=C'_b:=S_b(2r-2b+\varepsilon)\int_{[-2,2]^{r-b}}\Delta(u)^2\prod_lw(u_l)\,du>0,
\]
where $e_b=b^2-2br-\varepsilon b$ (so $e_0=0$). Consequently the expansion of $\Psi_b$ starts with $C'_bt^{e_b}$, without a logarithmic factor.
\end{lemma}

\begin{proof}
Elements of $\Sol(B_\varepsilon)$ have convergent Frobenius expansions at the regular singular point $0$, with integer exponents $0$ and $-1-\varepsilon$ and at most one power of $\log t$. Elements of $V$ are polynomials in such functions and their derivatives, which gives the first claim.

For the limit, substitute $u_l=-s_l/t$ ($s_l\ge2t$) for the $b$ coordinates on rays and keep the other $r-b$ coordinates $u'\in[-2,2]^{r-b}$. Then $e^{tu_l}=e^{-s_l}$, $du_l=t^{-1}ds_l$, and $\tilde w(u_l)=t^{-a-1/2}(s_l+2t)^{a}(s_l-2t)^{1/2}$ with $a+\frac12=\varepsilon$. Moreover $(u_l-u_{l'})^2=t^{-2}(s_l-s_{l'})^2$ for two ray coordinates, and $(u_l-u'_i)^2=t^{-2}(s_l+tu'_i)^2$ for a ray and a segment coordinate. Collecting powers of $t$,
\[
t^{-e_b}\Psi_b(t)=\int\!\!\int e^{t\sum u'_i}e^{-\sum s_l}\Delta(s)^2\prod_{l,i}(s_l+tu'_i)^2\prod_l(s_l+2t)^a(s_l-2t)^{1/2}\,\Delta(u')^2\prod_iw(u'_i)\,ds\,du',
\]
because $b+\varepsilon b+b(b-1)+2b(r-b)=-e_b$. For $0<t\le1$ the integrand is dominated by $e^{2(r-b)}e^{-\sum s}\Delta(s)^2\prod_l(s_l+2)^{2(r-b)+1}s_l^{\min(a,0)}\,\Delta(u')^2\prod w(u'_i)$, which is integrable. Dominated convergence gives the limit $S_b(2(r-b)+a+\frac12)\int\Delta(u')^2\prod w$, which equals $C'_b$. It is positive and finite, so the lowest term of the expansion of $\Psi_b$ is $C'_bt^{e_b}$.
\end{proof}

\section{Proof of \texorpdfstring{Theorem~\ref{bfk:thm:main}}{the main theorem}}\label{bfk:sec:proof}

\subsection{The order}
By Lemma~\ref{bfk:lem:vertex} the functions $\Psi_0,\dots,\Psi_r$ have pairwise distinct exponential rates $\mu_b$ along $\R_{>0}$, hence they are linearly independent. By \eqref{bfk:eq:Vspan}, $\dim V=r+1$, and Proposition~\ref{bfk:prop:V} gives (i).

From now on let $L=L_k=\sum_{m=0}^np_m\partial^m$ be the primitive minimal operator of $Y_k$, $n=r+1$. By Proposition~\ref{bfk:prop:V} its solution space has the basis $y_b=e^{\varepsilon t}\Psi_b$, $0\le b\le r$. Since $Y_k$ has rational Taylor coefficients, uniqueness of the minimal operator allows us to take $p_m\in\Q[t]$.

\subsection{Singular points}
Every $y_b$ extends analytically to the universal cover of $\C^*$, because $\varphi_I$ and $\varphi_K$ do. Hence at a zero $\tau\ne0$ of $p_n$ all solutions are holomorphic, and $\tau$ is an apparent singular point. By Lemma~\ref{bfk:lem:origin} all solutions have moderate growth at $0$, so $0$ is a regular singular point (Fuchs' theorem).

\subsection{Exponents at \texorpdfstring{$0$}{0}}
By Lemma~\ref{bfk:lem:origin} the expansion of $y_b=e^{\varepsilon t}\Psi_b$ starts with $C'_bt^{e_b}$. By Lemma~\ref{bfk:lem:indicial}, each $e_b$ is a root of the indicial polynomial $I_0$. Since $e_b-e_{b+1}=2r+\varepsilon-2b-1>0$ for $0\le b<r$, these are $n$ distinct numbers $0=e_0>e_1>\cdots>e_r$. As $\deg I_0=n$, they are exactly the exponents of $L$ at $0$. This proves the first half of (iii).

\subsection{Behaviour at \texorpdfstring{$\infty$}{infinity}}
Put $\lambda_b=\varepsilon+\mu_b=\varepsilon+2r-4b$. Since $y_b^{(m)}=e^{\varepsilon t}\sum_i\binom mi\varepsilon^{m-i}\Psi_b^{(i)}$, Lemma~\ref{bfk:lem:vertex} gives
\begin{equation}\label{bfk:eq:yinf}
y_b^{(m)}(t)=C_be^{\lambda_bt}t^{-\gamma_b}\bigl(\lambda_b^m+O(t^{-1})\bigr)\qquad(t\to+\infty,\ m\ge0).
\end{equation}
Applying $L$ termwise to the asymptotic expansion of $y_b$, which is legitimate by Lemma~\ref{bfk:lem:vertex}, gives the asymptotic expansion of $Ly_b=0$. Hence $e^{\lambda_bt}t^{-\gamma_b}\sum_ic_{b,i}t^{-i}$ is a formal solution of $L$ at $\infty$. The $\lambda_b$ are distinct, so these $n$ formal solutions form a basis. This proves the second half of (iii).

\subsection{No apparent singularities}
Let $W$ be the Wronskian of $y_0,\dots,y_r$ on $(0,\infty)$ and $R=W'/W=-p_{n-1}/p_n\in\C(t)$. The poles of $R$ lie among the zeros of $p_n$. At $0$ and at each apparent point the pole is simple, because $p_{n-1}/p_n$ has at most a simple pole at a regular singular point and $W$ is holomorphic at an apparent point. So
\[
R(t)=Q(t)+\sum_{\tau}\frac{\kappa_\tau}{t-\tau},\qquad Q\in\C[t],\quad \kappa_\tau=\Res_{t=\tau}R .
\]
By Lemma~\ref{bfk:lem:abel}, $\kappa_0=\sum_be_b-\binom n2$, and $\kappa_\tau\ge1$ for every zero $\tau\ne0$ of $p_n$. Integrating $R$ along $(t_0,\infty)$ with $t_0>\max_\tau|\tau|$ gives
\[
W(t)=c\,\exp\Bigl(\int^tQ\Bigr)\prod_\tau(t-\tau)^{\kappa_\tau},\qquad c\ne0 .
\]
On the other hand, \eqref{bfk:eq:yinf} gives
\[
W(t)=\Bigl(\prod_bC_b\Bigr)e^{(\sum_b\lambda_b)t}t^{-\sum_b\gamma_b}\bigl(V(\lambda)+O(t^{-1})\bigr),\qquad V(\lambda)=\prod_{b<b'}(\lambda_{b'}-\lambda_b)\ne0 .
\]
Comparing the two expressions as $t\to+\infty$ shows that $Q$ is the constant $\sum_b\lambda_b$ and that $\sum_\tau\kappa_\tau=-\sum_b\gamma_b$. Therefore
\[
\sum_{\tau\ne0}\kappa_\tau=-\sum_{b=0}^r\gamma_b-\sum_{b=0}^re_b+\binom n2 .
\]
Now, using $a+\frac12=\varepsilon$,
\begin{align*}
\sum_{b=0}^r\gamma_b&=\sum_{b=0}^r\Bigl[(r-b)^2+a(r-b)+b^2+\tfrac b2\Bigr]=\frac{r(r+1)(2r+1)}3+\varepsilon\,\frac{r(r+1)}2,\\
\sum_{b=0}^re_b&=\sum_{b=1}^r\bigl(b^2-2br-\varepsilon b\bigr)=\frac{r(r+1)(2r+1)}6-r^2(r+1)-\varepsilon\,\frac{r(r+1)}2 .
\end{align*}
Hence
\[
\sum_b\gamma_b+\sum_be_b=\frac{r(r+1)(2r+1)}2-r^2(r+1)=\frac{r(r+1)}2=\binom n2,
\]
and $\sum_{\tau\ne0}\kappa_\tau=0$. Since each $\kappa_\tau$ with $\tau\ne0$ is at least $1$, the polynomial $p_n$ has no zero other than $0$. This proves (iv): $p_n=c\,t^{A}$ for some $c\ne0$ and $A\ge0$.

\subsection{The shape of \texorpdfstring{$L_k$}{L\_k}}
By \eqref{bfk:eq:yinf} and Lemma~\ref{bfk:lem:top}, $\deg p_m\le\deg p_n=A$ for all $m$. Since $0$ is a regular singular point, Fuchs' criterion gives $\ord_0p_m\ge A-n+m$. The indicial polynomial is $I_0(\rho)=\sum_{m=0}^n[t^{A-n+m}]p_m\,\ff\rho m$, with the convention that coefficients of negative powers vanish.

If $A\le n-2$, the terms with $m<n-A$ vanish, so $I_0$ is divisible by $\ff\rho{\,n-A}$ and has the root $1$. This contradicts the fact that the exponents $e_b$ are all $\le0$. Hence $A\ge n-1$.

Since $0$ is an exponent, $I_0(0)=[t^{A-n}]p_0=0$, so $\ord_0p_m\ge A-n+1$ for every $m$. Primitivity forces $A-n+1\le0$. Hence $A=n-1=r$.

Consequently $p_n=c\,t^{r}$, $\deg p_m\le r$, and $\ord_0p_m\ge m-1$, which is (ii) after dividing by $c$. Finally, every operator of order $n$ annihilating $Y_k$ is $f\cdot L_k$ with $f\in\C(t)$ by minimality. If its coefficients are polynomials, then $f$ is a polynomial because $L_k$ is primitive, so its degree is at least $r$. \qed

\begin{proof}[Proof of Corollary~\ref{bfk:cor:BFK}]
Theorem~\ref{bfk:thm:main}(ii) provides the operator with $\ell=n$, $p_\ell=t^{\ell-1}$, $\deg p_m\le\ell-1$ and $t^{m-1}\mid p_m$.
\end{proof}


\section{A general Jacobi construction and its provenance}\label{bfk:sec:bfk-jacobi}
The following result is a scalar consequence of a published differential
system, rather than a claim of a new first solution. It makes the overlap
with earlier work precise and supplies an independent construction of the
coefficient shape in Conjecture~\ref{bfk:conj:BFK}.

Fix an integer $N\ge1$ and real parameters $\alpha>-1$, $b>0$ and $\beta>0$.
For $0\le p\le N$ define the entire functions
\begin{equation}\label{bfk:eq:bfk-Jacobi-H}
 H_p(x)=\binom Np^{-1}\int_{[0,1]^N}
 e^{-x\sum u_i}\prod_{i=1}^N u_i^\alpha(1-u_i)^{b-1}
 \prod_{i<j}|u_i-u_j|^\beta
 e_p(1-u_1,\ldots,1-u_N)\,du,
\end{equation}
where $e_p$ is the elementary symmetric polynomial and $e_0=1$.
These are exactly the functions of \cite[Corollary~2.5]{FK22}, with its
parameter $a$ renamed $\alpha$. Put
\[
 B_p=p\left(\alpha+b+1+\frac\beta2(2N-p-1)\right),\qquad
 D_p=p\left(b+\frac\beta2(N-p)\right).
\]
The cited result, equation~(2.20) of its arXiv version, is
\begin{equation}\label{bfk:eq:bfk-Jacobi-system}
 (N-p)xH_{p+1}=\bigl((N-p)x+B_p\bigr)H_p+xH_p'-D_pH_{p-1},
 \qquad 0\le p\le N,
\end{equation}
where the terms with an out-of-range index have zero coefficient.
In particular, its $p=0$ case gives $H_1=H_0+H_0'/N$.

\begin{proposition}[Scalar bounds from the Jacobi system]\label{bfk:prop:bfk-jacobi-shape}
There is an operator $\mathcal A=\sum_{m=0}^{N+1}a_m(x)\partial_x^m$
annihilating $H_0$, with
\[
 a_m\in\R[x],\qquad a_{N+1}=x^N,\qquad
 \deg a_m\le N,\qquad x^{m-1}\mid a_m\quad(1\le m\le N+1).
\]
If $\alpha,b,\beta$ are rational, the coefficients are rational.
The degree, leading-coefficient, and divisibility assertions also hold,
after constant normalization, for $e^{ct}H_0(dt)$ whenever $c\in\R$ and
$d\in\R\setminus\{0\}$; its coefficients are rational if
$\alpha,b,\beta,c,d$ are all rational.
This is an existence statement and does not assert
minimality for every choice of Jacobi parameters.
\end{proposition}
\begin{proof}
All operator products below are compositions, with coefficients written to
the left of derivatives. Write $D=\partial_x$ and define
\[
 P_0=1,\qquad P_1=1+D/N,
\]
and, for $1\le p<N$, define
\begin{equation}\label{bfk:eq:bfk-P-recursion}
 (N-p)P_{p+1}
 =\left(D+N-p+\frac{B_p}{x}\right)P_p
       -\frac{D_p}{x}P_{p-1}.
\end{equation}
Equation~\eqref{bfk:eq:bfk-Jacobi-system} shows inductively that $H_p=P_pH_0$.
Write $P_p=\sum_{m=0}^p q_{p,m}(x)D^m$, with missing coefficients zero.
For every $p\ge1$, the following three properties hold:
\begin{enumerate}
\item $q_{p,p}=1/\ff Np$;
\item each $q_{p,m}$ is a Laurent polynomial whose exponents lie between
$-(p-1)$ and $0$;
\item $\ord_0 q_{p,m}\ge m-p$, with $\ord_0 0=+\infty$.
\end{enumerate}
They hold for $P_1$. To propagate them, use
$D\circ(qD^m)=qD^{m+1}+q'D^m$ in
\eqref{bfk:eq:bfk-P-recursion}. Differentiating a Laurent coefficient lowers its
smallest exponent by at most one and produces no positive exponent.
Division by $x$ has the same lower-bound effect. The coefficient of $D^m$
in $D\circ P_p$ is $q_{p,m}'+q_{p,m-1}$, whose valuation is at least
$m-p-1$. The remaining terms satisfy the same bound. The highest derivative
comes only from $D\circ P_p$, giving the stated leading coefficient.
For the global Laurent bound, $P_{p-1}$ has smallest exponent at least
$-(p-2)$ when $p\ge2$; for $p=1$, $P_0=1$ gives the same required bound
directly. This proves all three assertions by induction.

The $p=N$ equation is
\[
 xH_N'+B_NH_N-D_NH_{N-1}=0.
\]
Consequently the operator
\begin{equation}\label{bfk:eq:bfk-A-construction}
 \mathcal A=N!x^{N-1}\bigl(xD\circ P_N+B_NP_N-D_NP_{N-1}\bigr)
\end{equation}
annihilates $H_0$. Its bracket has Laurent exponents between $-(N-1)$
and $1$: in $xD\circ P_N$, differentiation of a coefficient is compensated
by multiplication by $x$. Thus every coefficient of $\mathcal A$ is a
polynomial of degree at most $N$. Its leading coefficient is $x^N$.
The coefficient of $D^m$ in the bracket has valuation at least $m-N$,
by property~(3); after multiplication by $x^{N-1}$ this is at least $m-1$.
This proves divisibility. All recursions use only field operations in the
parameters and integer differentiation, proving the rationality assertion.

For the transformed function, substitute $x=dt$ and
$D=d^{-1}(\partial_t-c)$. The coefficient of $\partial_t^m$ is a constant
linear combination of the old coefficients with derivative indices $j\ge m$,
each divisible by $t^{j-1}$ when $m\ge1$ and of degree at most $N$.
The leading coefficient becomes $d^{-1}t^N$, so multiplication by $d$
normalizes it to $t^N$. The stated properties are preserved.
\end{proof}

\begin{corollary}[Exact BFK parameter specialization]\label{bfk:cor:bfk-FK-specialization}
The existence and coefficient-shape assertions of
Conjecture~\ref{bfk:conj:BFK} follow from the published system
\eqref{bfk:eq:bfk-Jacobi-system} and Proposition~\ref{bfk:prop:bfk-jacobi-shape}.
\end{corollary}
\begin{proof}
For $r\ge1$, substitute $u_i=4v_i-2$ in
Proposition~\ref{bfk:prop:int}. The resulting Jacobi parameters are
\[
 N=r,\qquad \alpha=\tfrac12,\qquad b=\varepsilon+\tfrac12,
 \qquad \beta=2.
\]
All normalization factors are constant in $t$, and specifically
\[
 Y_k(t)=\frac{4^{r(r+\varepsilon)}}{r!(2\pi)^r}
        e^{(\varepsilon-2r)t}H_0(-4t).
\]
Indeed the Vandermonde, weights and Jacobian contribute, respectively,
$4^{r(r-1)}$, $4^{\varepsilon r}$ and $4^r$.
Apply Proposition~\ref{bfk:prop:bfk-jacobi-shape} with $c=\varepsilon-2r$ and
$d=-4$. The case $k=1$ was treated separately.
\end{proof}

\begin{remark}[What this comparison does and does not establish]
The preceding corollary is a derivation made explicit in this compilation.
It is not evidence that \cite{FK22} itself announced the combinatorial
conjecture as solved. It does establish that existence and the full
coefficient shape are consequences of an earlier general system. The exact
minimality and exponent analysis in Sections~\ref{bfk:sec:V}--\ref{bfk:sec:proof}
are separate assertions. No exhaustive historical priority claim is made
for either the consequence or its additional refinements.
\end{remark}

\section{The recurrence, an example, and computations}\label{bfk:sec:rec}

\begin{proof}[Proof of Corollary~\ref{bfk:cor:rec}]
Write $p_m=t^{m-1}q_m$ for $m\ge1$ and put $q_0=tp_0$. Then $tL_k=\sum_{m=0}^nq_m(t)\ff\theta m$, where $\deg q_m\le r+1-m$ for $m\ge1$, $q_n=1$, $\deg q_0\le r+1$ and $q_0(0)=0$. Write $q_m=\sum_iq_{m,i}t^i$ and $a_N=y_k(N)/N!$. Since $\ff\theta m\,t^N=\ff Nm\,t^N$, taking the coefficient of $t^N$ in $tL_kY_k=0$ gives $\sum_i\bigl(\sum_mq_{m,i}\ff{(N-i)}m\bigr)a_{N-i}=0$. Multiply by $N!$ and use $N!\,a_{N-i}=\ff Ni\,y_k(N-i)$:
\[
\sum_{i=0}^{r+1}C_i(N)\,y_k(N-i)=0,\qquad C_i(N)=\ff Ni\sum_{m}q_{m,i}\ff{(N-i)}m .
\]
Since $q_{m,i}\ne0$ forces $m\le r+1-i$, we have $\deg C_i\le r+1$. Each $C_i$ is divisible by $N$: for $i\ge1$ because of the factor $\ff Ni$, and for $i=0$ because $q_{0,0}=0$. Put $c_i=C_i/N$. Then $\deg c_i\le r$, and $c_0$ has degree exactly $r$ because $q_{n,0}=1$.
\end{proof}

For example, for $k=2$ this gives $(N+1)y_2(N)-2y_2(N-1)-4(N-1)y_2(N-2)=0$, in agreement with \cite{BG00}. In general $c_{r+1}(N)=\ff{(N-1)}{r}\,[t^r]p_0$, and by Lemma~\ref{bfk:lem:top} the number $[t^r]p_0=\chi(0)$ vanishes exactly when some $\lambda_j$ is $0$, that is, when $\varepsilon=0$ and $r$ is even. In that case the recurrence has order at most $r$; for $k=4$ one obtains the second-order recurrence $(N+3)(N+4)y_4(N)-(8N+12)y_4(N-1)-16N(N-1)y_4(N-2)=0$ of \cite{BG00}.

\begin{example}[$k=12$]\label{bfk:ex:k12}
The case $k=12$ is the first one outside the range verified in \cite{BG00}. The primitive minimal operator is $L_{12}=\sum_{m=0}^7p_m\partial^m$ with
\begin{align*}
p_7&=t^6,\qquad p_6=182\,t^5,\qquad p_5=-14\,t^4\,(16t^2+8t-937),\\
p_4&=-8\,t^3\,(3304t^2+1743t-59554),\\
p_3&=4\,t^2\,(3136t^4+3136t^3-282848t^2-159136t+2280623),\\
p_2&=8\,t\,(98496t^4+105072t^3-2663704t^2-1626026t+10791605),\\
p_1&=-16\,(9216t^6+13824t^5-909792t^4-1053472t^3+10407753t^2+7124249t-19384596),\\
p_0&=-576\,(4608t^5+7488t^4-130240t^3-169050t^2+640332t+538461).
\end{align*}
Its indicial polynomial at $0$ is $\rho(\rho+11)(\rho+20)(\rho+27)(\rho+32)(\rho+35)(\rho+36)$, and the top-degree polynomial of Lemma~\ref{bfk:lem:top} is $\lambda(\lambda^2-16)(\lambda^2-64)(\lambda^2-144)$. Both agree with Theorem~\ref{bfk:thm:main}(iii). The source reports a Taylor check through $t^{384}$ and a separate lattice-walk check through $t^{39}$; those historical runs are not used as proof inputs here. For this compilation, all eight displayed coefficients were independently
verified in exact rational arithmetic using the finite Laurent-operator
recursion of Section~\ref{bfk:sec:bfk-jacobi}, followed by $x=-4t$ and the
exponential gauge. This is a finite algebraic identity check, separate from
the source-reported Taylor tests.
\end{example}

\begin{table}[ht]
\centering
\caption{Exponential rates $\lambda_j$, exponents $\gamma_j$ at $\infty$, and exponents $e_b$ at $0$ (Theorem~\ref{bfk:thm:main}(iii)).}\label{bfk:tab:exp}
\begin{tabular}{@{}rlll@{}}
\toprule
$k$ & $\lambda_0,\dots,\lambda_r$ & $\gamma_0,\dots,\gamma_r$ & $e_0,\dots,e_r$\\
\midrule
2 & $2,-2$ & $\tfrac12,\tfrac32$ & $0,-1$\\
3 & $3,-1$ & $\tfrac32,\tfrac32$ & $0,-2$\\
4 & $4,0,-4$ & $3,2,5$ & $0,-3,-4$\\
5 & $5,1,-3$ & $5,3,5$ & $0,-4,-6$\\
6 & $6,2,-2,-6$ & $\tfrac{15}2,\tfrac92,\tfrac{11}2,\tfrac{21}2$ & $0,-5,-8,-9$\\
7 & $7,3,-1,-5$ & $\tfrac{21}2,\tfrac{13}2,\tfrac{13}2,\tfrac{21}2$ & $0,-6,-10,-12$\\
8 & $8,4,0,-4,-8$ & $14,9,8,11,18$ & $0,-7,-12,-15,-16$\\
9 & $9,5,1,-3,-7$ & $18,12,10,12,18$ & $0,-8,-14,-18,-20$\\
10 & $10,6,2,-2,-6,-10$ & $\tfrac{45}2,\tfrac{31}2,\tfrac{25}2,\tfrac{27}2,\tfrac{37}2,\tfrac{55}2$ & $0,-9,-16,-21,-24,-25$\\
11 & $11,7,3,-1,-5,-9$ & $\tfrac{55}2,\tfrac{39}2,\tfrac{31}2,\tfrac{31}2,\tfrac{39}2,\tfrac{55}2$ & $0,-10,-18,-24,-28,-30$\\
12 & $12,8,4,0,-4,-8,-12$ & $33,24,19,18,21,28,39$ & $0,-11,-20,-27,-32,-35,-36$\\
\bottomrule
\end{tabular}
\end{table}

\subsection*{Computations}
The source manuscript reported additional finite computations. They are not assumptions of the proofs above and have not been reproduced in their entirety for this compilation. The following historical reports are retained only to identify what the source claimed:
\begin{enumerate}
\item The source reports that for $2\le k\le13$ it computed, in exact rational arithmetic, the space of operators of order $\fl{k/2}+1$ with coefficients of degree at most $\fl{k/2}$ annihilating the Taylor series of $Y_k$ (obtained from Theorem~\ref{bfk:thm:det}). In each case it is one-dimensional, and the operator has the shape of Theorem~\ref{bfk:thm:main}(ii). For $k\le11$ these operators coincide with those listed in \cite{BG00}, and for $k\le5$ the recurrences of Corollary~\ref{bfk:cor:rec} obtained from them coincide with the recurrences listed there.
\item The source reports that for $2\le k\le45$ it repeated the computation modulo the prime $p=67108859$, and for $k\in\{20,33,45\}$ also modulo $67108837$. In every case the space is one-dimensional, no operator with coefficients of degree at most $\fl{k/2}-1$ exists, and the indicial roots at $0$, the top-degree roots $\lambda_j$, and the exponents $\gamma_j=\psi(\lambda_j)/\chi'(\lambda_j)$ agree with Theorem~\ref{bfk:thm:main}(iii). Here $\psi(\lambda)=\sum_m[t^{r-1}]p_m\lambda^m$.
\item The source reports that a second implementation confirmed these statements for $k\le9$ in exact arithmetic, starting from a lattice-path count of $y_k(n)$ for $n\le110$, and for $10\le k\le20$ modulo $2^{61}-1$, starting from its own implementation of Theorem~\ref{bfk:thm:det}.
\end{enumerate}


\section{Integral directions in three tableau recurrences}\label{bfk:sec:bfk-lattices}
The following applications are inspired by the discrete Abel argument for
Franel recurrences. They concern a \emph{fixed recurrence with two free
initial values}, not the minimal order of an arbitrary recurrence for the
counting sequence. The height-$3$ integral-lattice assertion is the theorem
of Klazar--Luca \cite{KL05}, and its bounded-denominator formulation appears
in \cite[Theorem~3.4]{Mat25}. Our proof below gives a unified Casoratian
argument for heights $2,3,4$, together with explicit bounds at prime-power
indices. These are newly included deductions in this compilation; no first
priority or new solution of an external conjecture is claimed.

Let $U^{(k)}_n=y_k(n)$ for $k=2,3,4$. The equations already obtained above,
or directly listed in \cite{BG00}, give the following recurrences:
\begin{align}
 (n+2)f_{n+1}&=2f_n+4nf_{n-1}, && k=2,\label{bfk:eq:bfk-height2}\\
 (n+3)f_{n+1}&=(2n+3)f_n+3nf_{n-1}, && k=3,\label{bfk:eq:bfk-height3}\\
 (n+4)(n+5)f_{n+1}&=(8n+20)f_n+16n(n+1)f_{n-1},
 && k=4.\label{bfk:eq:bfk-height4}
\end{align}
In this section each equation is imposed only for $n\ge1$.
It therefore has a two-dimensional rational solution space, parametrized by
$f_0,f_1$. The counting solution is integral and satisfies
$U^{(k)}_0=U^{(k)}_1=1$. We do not impose the $n=0$ equation, which would
already force $f_1=f_0$.

\begin{proposition}[Explicit Casoratians]\label{bfk:prop:bfk-casoratians}
For a solution $f$ of the height-$k$ recurrence define
$W_n=U^{(k)}_nf_{n+1}-U^{(k)}_{n+1}f_n$ and $c=f_1-f_0$.
For every $n\ge0$,
\begin{align}
 W_n&=\frac{2(-4)^n}{(n+1)(n+2)}c, && k=2,\label{bfk:eq:bfk-W2}\\
 W_n&=\frac{6(-3)^n}{(n+1)(n+2)(n+3)}c, && k=3,\label{bfk:eq:bfk-W3}\\
 W_n&=\frac{2880(-16)^n}{(n+1)(n+2)^2(n+3)^2(n+4)^2(n+5)}c,
 && k=4.\label{bfk:eq:bfk-W4}
\end{align}
\end{proposition}
\begin{proof}
For $a_nf_{n+1}=b_nf_n+d_nf_{n-1}$, direct substitution for both solutions
shows $a_nW_n=-d_nW_{n-1}$. Hence
$W_n=(-1)^n\prod_{j=1}^n(d_j/a_j)\,c$.
For heights $2$ and $3$ the products of $j/(j+2)$ and $j/(j+3)$
are, respectively, $2/((n+1)(n+2))$ and
$6/((n+1)(n+2)(n+3))$.
For height $4$ the product is
\[
 \prod_{j=1}^n\frac{j(j+1)}{(j+4)(j+5)}
 =\frac{4!5!\,n!(n+1)!}{(n+4)!(n+5)!},
\]
which is exactly the rational factor in \eqref{bfk:eq:bfk-W4}.
All formulas also give $W_0=c$.
\end{proof}

For a prime $p$, write $v_p$ for the usual additive valuation on
$\Q$, with $v_p(0)=+\infty$.

\begin{theorem}[Integral lattices and local denominator growth]\label{bfk:thm:bfk-integral-lattices}
For each $k\in\{2,3,4\}$, the integral solutions of its recurrence
\eqref{bfk:eq:bfk-height2}--\eqref{bfk:eq:bfk-height4} are exactly
$\Z U^{(k)}$. The rational solutions with a common denominator are exactly
$\Q U^{(k)}$.
More precisely, if $c=f_1-f_0\ne0$, the following bounds hold:
\begin{enumerate}
\item For height $2$, every odd prime $p$ and every integer $a\ge1$ satisfy
\[
 \min\{v_p(f_{p^a-1}),v_p(f_{p^a})\}\le v_p(c)-a.
\]
\item For height $3$, the same inequality holds for every prime
$p\ge5$ and every $a\ge1$, and also for $p=2$ and every $a\ge2$.
\item For height $4$, every odd prime $p$ and every $a\ge1$ satisfy
\[
 \min\{v_p(f_{p^a-2}),v_p(f_{p^a-1})\}
       \le v_p(2880c)-2a.
\]
\end{enumerate}
Thus every solution outside the counting line has unbounded denominators
at each of the indicated primes.
\end{theorem}
\begin{proof}
Integrality of $U^{(k)}$ gives the general bound
\begin{equation}\label{bfk:eq:bfk-valuation-W}
 v_p(W_n)\ge\min\{v_p(f_n),v_p(f_{n+1})\}.
\end{equation}
For height $2$, take $n=p^a-1$ in \eqref{bfk:eq:bfk-W2}.
For odd $p$, its valuation is exactly $v_p(c)-a$.
For height $3$ and $p\ge5$, the same substitution in
\eqref{bfk:eq:bfk-W3} again gives exactly $v_p(c)-a$.
If $p=2$ and $a\ge2$, the three denominator factors have valuations
$a,0,1$, while $v_2(6)=1$ and $v_2(3)=0$; the result is again
$v_2(c)-a$.
For height $4$, take $n=p^a-2$ in \eqref{bfk:eq:bfk-W4}.
The factor $(n+2)^2$ has valuation $2a$, all other denominator factors
are nonzero integers, and $p\nmid16$. Hence
$v_p(W_n)\le v_p(2880c)-2a$.
Combining each calculation with \eqref{bfk:eq:bfk-valuation-W} proves the bounds.

An integral solution cannot satisfy these bounds with $c\ne0$, since the
right sides tend to $-\infty$. Thus $f_1=f_0$. Each forward coefficient
in the recurrences is nonzero for $n\ge1$, so uniqueness from the two
initial values gives $f=f_0U^{(k)}$. Its scalar $f_0$ is integral.
Conversely, integer multiples of the counting solution are integral.
If a rational solution has common denominator $D$, apply this result to
$Df$ to obtain $f\in\Q U^{(k)}$. The converse is immediate.
\end{proof}

\begin{corollary}[One-prime and localization criteria]\label{bfk:cor:bfk-localization}
Let $\mathcal P_2=\mathcal P_4$ be the set of odd primes and let
$\mathcal P_3$ be the set of primes different from $3$.
A rational solution of the height-$k$ recurrence belongs to
$\Q U^{(k)}$ if and only if it is bounded in $\Q_p$ for at least one
$p\in\mathcal P_k$; this is then true for every such prime.
For every finite set of primes $S$, the solutions with all terms in
$\Z[S^{-1}]$ are exactly $\Z[S^{-1}]U^{(k)}$.
\end{corollary}
\begin{proof}
For $f=qU^{(k)}$, integrality gives $v_p(f_n)\ge v_p(q)$ for every $n$.
The preceding theorem excludes boundedness at any indicated prime when
$c\ne0$. For the localization statement choose
$p\in\mathcal P_k\setminus S$. A solution over $\Z[S^{-1}]$ is integral
at $p$, hence belongs to the counting line; its scalar is $f_0$ and is
therefore in $\Z[S^{-1}]$. Conversely such a scalar times the integral
counting sequence lies in the localization.
\end{proof}

\begin{remark}
These lattice statements do not assert that the recurrences at larger
heights have only one integral direction. They also give no irrationality
statement for an associated limit without an independent approximation
estimate. The analogy with the Franel argument is the use of the exact
Casoratian to turn prime-power denominators into rigidity of initial values.
\end{remark}

\section{Remarks}\label{bfk:sec:rem}

\begin{remark}
The argument shows that the order $\fl{k/2}+1$ is governed by the $S_r$-orbits of vertices of the cube $[-2,2]^r$: the orbit with $j$ coordinates equal to $-2$ produces the exponential type $e^{\lambda_jt}$. The absence of apparent singularities is equivalent to the identity $\sum_j\gamma_j+\sum_be_b=\binom n2$. Questions about minimal orders of recurrences for related binomial sums, such as the generalized Franel numbers, are discussed in \cite{SZ23}. Arguments based on local exponents and the Fuchs relation are classical in the study of apparent singularities and degrees of minimal operators; see, for example, \cite{BCLSS07,BRS}.
\end{remark}

\begin{remark}
Gessel's formula $\sum_nu_k(n)x^n/(n!)^2=\det\bigl(I_{i-j}(2\sqrt x)\bigr)_{i,j=1}^k$ for the number $u_k(n)$ of permutations of $\{1,\dots,n\}$ with no increasing subsequence of length $k+1$ \cite{Ges90} leads to integrals over the unitary group $U(k)$. Forrester and Wei \cite{FW26} derived scalar differential equations of order $k+1$ for this family. In the computations reported by the source manuscript for $2\le k\le16$ (modulo the prime $67108859$), the operators in $x$ of order $k+1$ with coefficients of degree at most $k$ annihilating this series form a one-dimensional space, and there is no such operator with coefficients of degree at most $k-1$. The operator has leading coefficient $x^k$, and the coefficient of its $m$-th derivative is divisible by $x^{m-1}$ and has degree exactly $k-\lceil(k+1-m)/2\rceil$. These reported computations are not rerun here. A proof of all the stated degree refinements is not supplied, and they are not included among the proved results of this compilation.
\end{remark}


\clearpage
\part{Banana periods and minimal differential order}\label{part:order}


\section{Banana periods and exact differential orders}

\subsection{Banana integrals and the conjecture}
The $(n-1)$-loop banana (or sunset) graph has two vertices joined by $n$ edges. In two space-time dimensions the associated Feynman integral, with internal masses $m_1,\dots,m_n$ and external momentum $p$, is a relative period of a family of Calabi--Yau varieties of dimension $n-2$ parametrized by $t=p^2$ \cite{BKV15,Van14,BFKNS21,LV23}. Its maximal cut, which is the period over the torus $|x_1|=\dots=|x_n|=1$ in Feynman-parametric form, equals (up to normalization) the series
\begin{equation}\label{bo:eq:varpi}
\varpi_m(t)=\sum_{k\ge0}c_k(m)\,t^{-k-1},\qquad
c_k(m)=\sum_{\substack{a\in\N^n\\ |a|=k}}\binom{k}{a_1,\dots,a_n}^{2}\prod_{i=1}^n m_i^{2a_i},
\end{equation}
which converges for $|t|>R^2$, $R=\sum_i|m_i|$ (see \cite[\S4]{LV23}). Indeed, $c_k(m)$ is the constant term of $\bigl(\sum_i x_i^{-1}\bigr)^k\bigl(\sum_j m_j^2x_j\bigr)^k$. For positive masses, $c_k(m)$ is the $k$-th moment of $|\sum_i m_ie^{i\theta_i}|^2$ for independent uniform angles, so $\varpi_m(t)$ is the Stieltjes transform of the distribution of the squared distance travelled by a planar random walk with step lengths $m_1,\dots,m_n$. For equal masses this is the setting of \cite{BSWZ12}.

The function $\varpi_m$ is D-finite in $t$. Lairez and Vanhove computed minimal Picard--Fuchs operators for many Feynman integrals \cite{LV23}, and for the banana family they made the following conjecture.

\begin{conjecture}[{Lairez--Vanhove, order part of \cite[Conj.~4.0.1]{LV23}}]\label{bo:conj:LV}
For $n\ge2$ and generic masses, the minimal differential operator annihilating $\varpi_m$ has order
\[
o_n=2^n-\binom{n+1}{\lfloor (n+1)/2\rfloor}.
\]
\end{conjecture}

The values are $1,2,6,12,29,58,130$ for $n=2,\dots,8$. In \cite[\S\S4.1--4.5]{LV23} the cases $n\le6$ are established. For $n=7$, \cite[\S4.6]{LV23} reports a computation at fixed random integer masses modulo a large prime, producing an order-$58$ operator. This supplies evidence compatible with the conjecture, rather than a characteristic-zero proof of generic minimality. Conjecture~4.0.1 of \cite{LV23} also describes the leading coefficient of the minimal operator: apart from apparent singularities, its singular points in $\C$ should be $t=0$ and the thresholds $\bigl(\sum_i\varepsilon_im_i\bigr)^2$. That second statement is the subject of Part~\ref{part:structure} ; here we only prove the order statement. The same count was proposed in \cite{BFKNS21} on the basis of the cases with at most four loops, and the formula is also stated in the recent work \cite[\S3.1]{Van26} on the multiloop sunset. For equal masses the order is known to be at most $n-1$; operators of this order were constructed by Vanhove \cite{Van14} (see also \cite{BS08,BSWZ12}).

\subsection{Results}
Let $V=\C^2$ be the standard representation of $\sltwo$ and $V_d=\Sym^dV$ its irreducible representation of dimension $d+1$. For $m\in(\C^\times)^n$ let $\mu_1^2,\dots,\mu_k^2$ be the distinct values among $m_1^2,\dots,m_n^2$, with multiplicities $n_1,\dots,n_k$. Decompose
\begin{equation}\label{bo:eq:W}
W_m:=V_{n_1}\otimes\cdots\otimes V_{n_k}\;\cong\;\bigoplus_{d\ge0}V_d^{\oplus a_d(m)}
\end{equation}
under the diagonal action of $\sltwo$, and put
\begin{equation}\label{bo:eq:o}
o(m)=\sum_{d\ge2}(d-1)\,a_d(m).
\end{equation}

\begin{theoremA}\label{bo:thm:main}
Let $n\ge2$ and $m\in(\C^\times)^n$. Then the following hold.
\begin{enumerate}
\item[(i)] The minimal order of a nonzero linear differential operator with coefficients in $\C(t)$ annihilating $\varpi_m$ is $o(m)$.
\item[(ii)] The $\cD$-module generated by $\varpi_m$ is irreducible over $\C(t)$. Equivalently, the minimal operator of $\varpi_m$ is irreducible in $\C(t)[\partial_t]$.
\item[(iii)] The local monodromy of the minimal operator at $t=\infty$ is unipotent. It has one Jordan block of size $d-1$ for each summand $V_d$ of \eqref{bo:eq:W} with $d\ge2$.
\end{enumerate}
\end{theoremA}

\begin{corollaryB}
Conjecture~\ref{bo:conj:LV} holds for all $n\ge2$ and all masses with pairwise distinct squares $m_1^2,\dots,m_n^2$.
\end{corollaryB}

\begin{corollaryC}\label{bo:cor:equal}
If $m_1=\dots=m_n\ne0$, the minimal order is $n-1$, and the local monodromy at $\infty$ is a single unipotent Jordan block of size $n-1$ (maximal unipotent monodromy).
\end{corollaryC}

For pairwise distinct squares, $W_m=V^{\otimes n}$ and $o(m)=o_n$ by Lemma~\ref{bo:lem:count}. For equal masses, $W_m=V_n$ and $o(m)=n-1$. Table~\ref{bo:tab:conf} lists the values of $o(m)$ for $n=4,5$, and Section~\ref{bo:sec:numerics} records numerical comparisons reported in the supplied source. For $n=4$, the dependence of the maximal cuts on the mass configuration was studied from the point of view of modular forms in \cite{Duh25}. For generic masses with $n\ge4$, Theorem~A(iii) shows that the local monodromy at $\infty$ is not a single Jordan block. For instance, for $n=4$ the blocks have sizes $3,1,1,1$, whereas for $n=3$ there is a single block of size $2$.

\subsection{Strategy}
The proof uses the Bessel representation
\begin{equation}\label{bo:eq:bessel}
\varpi_m(t)=\frac12\int_0^\infty K_0\bigl(\sqrt{ty}\bigr)\,G_m(y)\,dy,\qquad G_m(y)=\prod_{i=1}^nI_0\bigl(m_i\sqrt y\bigr)
\end{equation}
(Lemma~\ref{bo:lem:reps}). This is the variant, adapted to the variable $y=x^2$, of the familiar position-space formula for the maximal cut.

\emph{Upper bound} (Section~\ref{bo:sec:upper}). The integrand is a solution of a connection $\cN$ on $\Gm$ of rank $2\dim W_m$, a tensor product of Bessel connections. It is regular at $0$ with unipotent monodromy, and every formal component at $\infty$ has slope $1/2$. The Euler characteristic formula gives $\dim H^1_{\rm dR}(\Gm,\cN)=\dim W_m$. Removing the invariants of the local monodromy at $0$ leaves a middle part of dimension $o(m)$. The integrand defines a class $\omega$ in this middle part, and the Gauss--Manin connection in $t$ preserves the middle part. The integral \eqref{bo:eq:bessel} is a period of a genuine rapid decay cycle; this is where $n\ge2$ is used.

\emph{Lower bound} (Section~\ref{bo:sec:lower}). By Kolchin and Goursat, the Bessel equations attached to $\mu_1,\dots,\mu_k$ have differential Galois group $\SL_2^k$. Hence the $\cD$-module generated by $G_m$ is irreducible of rank $\dim W_m$. We write $\varpi_m$ as a Laplace transform, followed by an inversion and a second Laplace transform, of $G_m$. Fourier transforms and intermediate extensions preserve simplicity, and boundary effects contribute only constant or $\delta$-type composition factors. Hence the $\cD$-module of $\varpi_m$ has a simple composition factor whose rank, computed by the vanishing cycle formula and stationary phase, is again $o(m)$.

As usual for minimal orders of D-finite functions, the lower bound is the harder part; compare Franel's suspicion on sums of powers of binomial coefficients, which \cite{SZ23} settle for recurrences obtained by creative telescoping. The case $n=1$, where $\varpi_m=(t-m_1^2)^{-1}$ has order $1>o(m)=0$, is excluded precisely because the rapid decay cycle cannot be closed (Remark~\ref{bo:rem:n1}).

\begin{table}[ht]
\centering
\caption{Minimal orders $o(m)$ for mass configurations with $n=4,5$. A configuration $[n_1,\dots,n_k]$ means $k$ distinct squared masses with multiplicities $n_j$. The right column lists the Jordan block sizes at $t=\infty$ predicted by Theorem~A(iii).}\label{bo:tab:conf}
\begin{tabular}{@{}lll@{\qquad}lll@{}}
\toprule
$n=4$ & $o(m)$ & blocks at $\infty$ & $n=5$ & $o(m)$ & blocks at $\infty$\\
\midrule
$[1,1,1,1]$ & $6$ & $3,1,1,1$ & $[1,1,1,1,1]$ & $12$ & $4,2,2,2,2$\\
$[2,1,1]$ & $5$ & $3,1,1$ & $[2,1,1,1]$ & $10$ & $4,2,2,2$\\
$[2,2]$ & $4$ & $3,1$ & $[2,2,1]$ & $8$ & $4,2,2$\\
$[3,1]$ & $4$ & $3,1$ & $[3,1,1]$ & $8$ & $4,2,2$\\
$[4]$ & $3$ & $3$ & $[3,2]$ & $6$ & $4,2$\\
 & & & $[4,1]$ & $6$ & $4,2$\\
 & & & $[5]$ & $4$ & $4$\\
\bottomrule
\end{tabular}
\end{table}

\subsection*{Conventions}
A $\cD$-module on an open set $U\subseteq\A^1$ means a left module over $\cO(U)\langle\partial\rangle$. A connection is a $\cD$-module that is free of finite rank over $\cO(U)$. For a function $f$ in a differential field or in a space of analytic functions, $\cD\cdot f$ is the $\cD$-module it generates, and $\ord f$ is the order of its minimal operator. Thus $\ord f=\dim_{\C(t)}\C(t)\otimes\cD\cdot f$, the generic rank of $\cD\cdot f$. We write $\theta=y\,d/dy$.

\section{Facts about \texorpdfstring{$\sltwo$}{sl2}-modules}\label{bo:sec:rep}

Let $(e,h,f)$ be the standard basis of $\sltwo$. Fix a basis $I,K$ of $V$ with $eI=0$ and $eK=I$, so that $I$ has weight $1$ and $K$ has weight $-1$. On $V_d$ the nilpotent $e$ acts with a single Jordan block of size $d+1$; its kernel is the highest weight line, and its image is the sum of all weight spaces except the lowest one.

\begin{lemma}\label{bo:lem:jordan}
Let $T_j=\exp(\kappa e)$ on $V$ with $\kappa\ne0$, and let $T=\bigotimes_j\Sym^{n_j}T_j$ act on $W_m=\bigotimes_jV_{n_j}$. Then $T$ is unipotent, and it has one Jordan block of size $d+1$ for each summand $V_d$ in \eqref{bo:eq:W}. In particular $\dim\ker(T-1)=\sum_da_d$, and $\im(T-1)=e\,W_m$.
\end{lemma}

\begin{proof}
We have $T=\exp(\kappa\rho(e))$, where $\rho$ is the diagonal action of $\sltwo$ on $W_m$. The Jordan blocks of the nilpotent $\rho(e)$ correspond to the irreducible summands, with sizes $\dim V_d=d+1$. The same holds for $\exp(\kappa\rho(e))$, since $\exp(\kappa\rho(e))-1=\kappa\rho(e)\,u$ with $u$ invertible and commuting with $\rho(e)$.
\end{proof}

\begin{lemma}\label{bo:lem:count}
\begin{enumerate}
\item[(a)] $\sum_d(d+1)a_d=\dim W_m=\prod_j(n_j+1)$.
\item[(b)] The number of irreducible summands of $V\otimes W_m$ is $2\sum_{d\ge1}a_d+a_0$. Hence
\[
\dim W_m-\#\{\text{summands of }V\otimes W_m\}=\sum_{d\ge2}(d-1)a_d=o(m).
\]
\item[(c)] If all $n_j=1$, so that $W_m=V^{\otimes n}$, then $o(m)=2^n-\binom{n+1}{\lfloor(n+1)/2\rfloor}=o_n$.
\end{enumerate}
\end{lemma}

\begin{proof}
(a) is clear. (b) By the Clebsch--Gordan rule, $V\otimes V_0=V_1$ and $V\otimes V_d=V_{d+1}\oplus V_{d-1}$ for $d\ge1$. Then
\[
\sum_d(d+1)a_d-2\sum_{d\ge1}a_d-a_0=\sum_{d\ge1}(d-1)a_d=\sum_{d\ge2}(d-1)a_d .
\]
(c) For $W_m=V^{\otimes n}$ we have $V\otimes W_m=V^{\otimes(n+1)}$. The number of irreducible summands of $V^{\otimes N}$ equals the dimension of its weight space of weight $0$ or $1$ (according to the parity of $N$), namely $\binom{N}{\lfloor N/2\rfloor}$. With $N=n+1$, (a) and (b) give $o(m)=2^n-\binom{n+1}{\lfloor(n+1)/2\rfloor}$.
\end{proof}

\begin{lemma}\label{bo:lem:weight}
Let $w\in W_m$ be a highest weight vector of weight $n=\sum_jn_j$; it spans the summand $V_n$, which has multiplicity one. If $n\ge2$, then $K\otimes w\in e\,(V\otimes W_m)$. If $n=1$, then $K\otimes w\notin e\,(V\otimes W_m)$.
\end{lemma}

\begin{proof}
The vector $K\otimes w$ lies in $V\otimes V_n\cong V_{n+1}\oplus V_{n-1}$ and has weight $n-1$. Its component in $V_{n+1}$ is in the image of $e$, because $n-1$ is not the lowest weight $-(n+1)$. Its component in $V_{n-1}$ is a highest weight vector, which lies in the image of $e$ if and only if $n-1>-(n-1)$, that is, $n\ge2$. For $n=1$ we have $w=I$ and $K\otimes I=\tfrac12(K\otimes I+I\otimes K)+\tfrac12(K\otimes I-I\otimes K)$, where the second summand is a nonzero element of $V_0\subset V\otimes V$. Since $e$ maps $V_2$ into $V_2$ and kills $V_0$, the image of $e$ meets $V_0$ trivially, so $K\otimes I$ is not in the image of $e$.
\end{proof}

\section{Bessel connections}\label{bo:sec:bessel}

For $\mu\in\C^\times$ let $\beta_\mu=\theta^2-\tfrac{\mu^2}{4}y$ and let $B_\mu=\cD_{\Gm}/\cD_{\Gm}\beta_\mu$, a connection of rank $2$ on $\Gm=\operatorname{Spec}\C[y,y^{-1}]$. It depends only on $\mu^2$.

\begin{lemma}\label{bo:lem:bessel}
\begin{enumerate}
\item[(a)] The solutions of $\beta_\mu u=0$ are spanned by $I_0(\mu\sqrt y)$, which is entire in $y$, and $K_0(\mu\sqrt y)$.
\item[(b)] The point $y=0$ is a regular singular point with both exponents $0$. In the basis $(I_0,K_0)$ the local monodromy is $I_0\mapsto I_0$, $K_0\mapsto K_0-\pi i\,I_0$, a single unipotent Jordan block.
\item[(c)] At $y=\infty$ the formal structure is $E^{\mu\sqrt y}\oplus E^{-\mu\sqrt y}$, twisted by a regular rank one factor. It is irreducible over $\C((y^{-1}))$ and has slope $1/2$.
\item[(d)] The differential Galois group of $\beta_\mu$ over $\C(y)$ is $\SL_2(\C)$.
\end{enumerate}
\end{lemma}

\begin{proof}
With $x=\sqrt y$ we have $\theta=\tfrac12x\,d/dx$, so $\beta_\mu u=0$ is the modified Bessel equation of order $0$ in $\mu x$. This gives (a), and (b) follows from $K_0(z)=-(\log(z/2)+\gamma)I_0(z)+(\text{even power series})$. Part (c) follows from the asymptotics $I_0(z)\sim e^z/\sqrt{2\pi z}$ and $K_0(z)\sim\sqrt{\pi/(2z)}\,e^{-z}$ in their standard sectors \cite[\S10.40(i), (10.40.1)--(10.40.2)]{DLMF}. For (d), the Galois group of the Bessel equation of order $0$ over $\C(x)$ is $\SL_2$ by Kolchin \cite{Kol68}. Over the subfield $\C(y)$ of index $2$, the Galois group contains that group, and it is contained in $\SL_2$ because the Wronskian of $\beta_\mu$ lies in $\C(y)$.
\end{proof}

\begin{proposition}\label{bo:prop:galois}
Let $\mu_1^2,\dots,\mu_k^2$ be pairwise distinct and nonzero. The differential Galois group of $\bigoplus_jB_{\mu_j}$ over $\C(y)$, acting on $\bigoplus_j\Sol(\beta_{\mu_j})$, is $\SL_2^k$.
\end{proposition}

\begin{proof}
By Lemma~\ref{bo:lem:bessel}(d), the Galois group $\Gamma$ is a closed subgroup of $\SL_2^k$ mapping onto each factor. Its identity component $\Gamma^0$ also maps onto each factor, since the image is a connected subgroup of finite index in $\SL_2$. By the Goursat--Kolchin--Ribet lemma \cite[Prop.~1.8.2]{Kat90}, $\Gamma^0=\SL_2^k$ unless, for some $i\ne j$, the representations $\Sol(\beta_{\mu_i})$ and $\Sol(\beta_{\mu_j})$ of $\Gamma^0$ become isomorphic after twisting by a character. Suppose this happens, and let $F$ be the fixed field of $\Gamma^0$, a finite extension of $\C(y)$. Then $B_{\mu_i}\otimes F\cong(B_{\mu_j}\otimes F)\otimes L$ for a rank one connection $L$ over $F$. At a point of $F$ above $y=\infty$, the exponential factors of $B_{\mu_i}$ are $E^{\pm\mu_i\sqrt y}$ (Lemma~\ref{bo:lem:bessel}(c)), while those of $B_{\mu_j}\otimes L$ are $E^{\pm\mu_j\sqrt y+q}$, where $E^q$ is the exponential factor of $L$. Hence
\[
\{\mu_i\sqrt y,\,-\mu_i\sqrt y\}=\{\mu_j\sqrt y+q,\,-\mu_j\sqrt y+q\}.
\]
Comparing the sums of the two sets gives $q=0$, and then $\mu_i^2=\mu_j^2$, a contradiction. Thus $\Gamma^0=\SL_2^k$, and so $\Gamma=\SL_2^k$.
\end{proof}

\begin{lemma}\label{bo:lem:simple}
Let $F\supseteq\C$ be a field and let $E$ be a connection on $\Gm$ over $F$, that is, a free $F[y,y^{-1}]$-module of finite rank with an action of $\partial_y$ satisfying the Leibniz rule. If $E\otimes F(y)$ is irreducible, then $E$ is a simple $F[y,y^{-1}]\langle\partial_y\rangle$-module. Consequently, every nonzero $F[y,y^{-1}]\langle\partial_y\rangle$-linear map $s$ from $E$ to another module is injective, and $s(E)$ is generated by $s(\lambda)$ for any $\lambda\ne0$.
\end{lemma}

\begin{proof}
Let $E'\subseteq E$ be a nonzero submodule. Since $E$ is torsion free, $E'\otimes F(y)\ne0$, so $E'\otimes F(y)=E\otimes F(y)$. Hence $Q=E/E'$ is a finitely generated torsion $F[y,y^{-1}]$-module stable under $\partial_y$. Let $p\in F[y]$ with $p(0)\ne0$ generate its annihilator. For $q\in Q$ we have $0=\partial_y(pq)=p'q+p\,\partial_yq=p'q$, so $p'\in pF[y,y^{-1}]$. Since $p$ is prime to $y$, it divides $p'$ in $F[y]$, so $p'=0$ by comparing degrees, and $p$ is constant. Thus $Q=0$ and $E'=E$. For the second statement, $\ker s$ is a proper submodule, hence zero, and $\lambda$ generates $E$ for every $\lambda\ne0$.
\end{proof}

\begin{corollary}\label{bo:cor:M}
Let $m\in(\C^\times)^n$ with distinct squared masses $\mu_j^2$ and multiplicities $n_j$, and put $G_m(y)=\prod_iI_0(m_i\sqrt y)=\prod_jI_0(\mu_j\sqrt y)^{n_j}$.
\begin{enumerate}
\item[(a)] The minimal operator of $G_m$ over $\C(y)$ has solution space $\bigotimes_j\Sym^{n_j}\Sol(\beta_{\mu_j})\cong W_m$. Hence $M:=(\cD\cdot G_m)|_{\Gm}$ is an irreducible connection of rank $\dim W_m$.
\item[(b)] At $y=0$, $M$ is regular singular with unipotent local monodromy. Under the identification of the solution space with $W_m$, this monodromy is $\exp(-\pi i\,\rho(e))$, and the function $G_m$ corresponds to the highest weight vector $\bigotimes_jI^{n_j}$.
\item[(c)] At $y=\infty$, every formal component of $M$ has the form $E^{c\sqrt y}\otimes\mathcal R$ with $\mathcal R$ regular and $c\in\C$. In particular all slopes of $M$ at $\infty$ are $\le1/2$.
\end{enumerate}
\end{corollary}

\begin{proof}
(a) Let $M'=\bigotimes_j\Sym^{n_j}B_{\mu_j}$, a connection on $\Gm$ of rank $\dim W_m$ whose solution space is $\bigotimes_j\Sym^{n_j}\Sol(\beta_{\mu_j})$; here a power $g^{n}$ of a solution $g$ of $B_\mu$ is the solution of $\Sym^{n}B_\mu$ sending $\lambda_1\cdots\lambda_n$ to $g(\lambda_1)\cdots g(\lambda_n)$. By Proposition~\ref{bo:prop:galois}, the Galois group $\SL_2^k$ acts on this space through the irreducible representation $\bigotimes_j\Sym^{n_j}V$, so $M'\otimes\C(y)$ is irreducible. Let $b_j\in B_{\mu_j}$ be the class of $1$. The solution $\bigotimes_jI_0(\mu_j\sqrt y)^{n_j}$ of $M'$ sends $\bigotimes_jb_j^{n_j}$ to $G_m$, so by Lemma~\ref{bo:lem:simple} it maps $M'$ isomorphically onto $(\cD\cdot G_m)|_{\Gm}$. Hence $M\cong M'$ is an irreducible connection of rank $\dim W_m$. Moreover, the solutions of the minimal operator of $G_m$ are the images of $\bigotimes_jb_j^{n_j}$ under the solutions of $M'$. In other words, the multiplication map from $\bigotimes_j\Sym^{n_j}\Sol(\beta_{\mu_j})$ to functions is injective, and its image is the solution space of the minimal operator. (b) follows from Lemma~\ref{bo:lem:bessel}(b) and Lemma~\ref{bo:lem:jordan} with $\kappa=-\pi i$. (c) follows from Lemma~\ref{bo:lem:bessel}(c); here $c$ runs over the numbers $\sum_j\ell_j\mu_j$ with $\ell_j$ a weight of $V_{n_j}$, and $c=0$ may occur.
\end{proof}

\section{Integral representations}\label{bo:sec:reps}

\begin{lemma}\label{bo:lem:reps}
Let $m\in(\C^\times)^n$, $R=\sum_i|m_i|$ and $t>R^2$. Then $G_m(y)=\sum_kc_k(m)\,(y/4)^k/(k!)^2$, and
\begin{equation}\label{bo:eq:repK}
\varpi_m(t)=\frac12\int_0^\infty K_0\bigl(\sqrt{ty}\bigr)G_m(y)\,dy .
\end{equation}
Moreover, let $\widehat G(s)=\int_0^\infty e^{-sy}G_m(y)\,dy$ for $\Re s>0$ and $H(v)=v^{-1}\widehat G(1/v)$. Then $H(v)=\sum_kc_k(m)v^k/(4^kk!)$ is entire, and
\begin{equation}\label{bo:eq:repL}
\varpi_m(t)=\frac14\int_0^\infty e^{-tv/4}H(v)\,dv .
\end{equation}
\end{lemma}

\begin{proof}
Since $I_0(\mu\sqrt y)=\sum_a(\mu^2y/4)^a/(a!)^2$, the coefficient of $(y/4)^k$ in $G_m$ is
\[
\sum_{|a|=k}\prod_i\frac{m_i^{2a_i}}{(a_i!)^2}=\frac{c_k(m)}{(k!)^2}.
\]
Moreover, $|G_m(y)|\le\prod_iI_0(|m_i|\sqrt{|y|})\le e^{R\sqrt{|y|}}$. By the standard integral $\int_0^\infty K_0(\sqrt{ty})\,y^k\,dy=2\cdot4^k(k!)^2t^{-k-1}$, termwise integration (justified by dominated convergence, as $K_0(z)\le Ce^{-z}$ for $z\ge1$) gives \eqref{bo:eq:repK}. Similarly, $\widehat G(s)=\sum_kc_k(m)k!\,4^{-k}(k!)^{-2}s^{-k-1}$ converges for every $s\ne0$, so $H(v)=\sum_kc_k(m)v^k/(4^kk!)$ is entire of exponential type at most $R^2/4$. Then \eqref{bo:eq:repL} follows by termwise integration for $t>R^2$.
\end{proof}

\section{The upper bound}\label{bo:sec:upper}

Throughout this section $\Kt=\C(t)$, $\overline{\Kt}$ is an algebraic closure, $m$ and $W_m$ are as above, and $N:=\dim W_m$.

\subsection{The connection \texorpdfstring{$\cN$}{N}}
Let $\beta_t=\theta^2-\tfrac t4y$ and let $B_t=\cD/\cD\beta_t$ be the corresponding rank $2$ connection on $\Gm$ over $\Kt$; thus $B_t=B_{\sqrt t}$ in the notation of Section~\ref{bo:sec:bessel}. Put
\[
\cN=B_t\otimes\bigotimes_{j=1}^k\Sym^{n_j}B_{\mu_j}\qquad(\text{tensor products over }\Kt[y,y^{-1}]),
\]
a connection on $\Gm$ over $\Kt$ of rank $2N$. Let $b_t\in B_t$ and $b_j\in B_{\mu_j}$ be the classes of $1$. The elements
\begin{equation}\label{bo:eq:basis}
\theta^ab_t\otimes\bigotimes_jb_j^{\,n_j-r_j}(\theta b_j)^{r_j}\qquad(a\in\{0,1\},\ 0\le r_j\le n_j)
\end{equation}
form a $\Kt[y,y^{-1}]$-basis of $\cN$. Let $\cN_0$ be their $\Kt[y]$-span, and put $\eta=b_t\otimes\bigotimes_jb_j^{\,n_j}\in\cN_0$.

The connection $\cN$ is defined over $\C[t]$. For $t_0\in\C^\times$, every element of $\cN$ whose coefficients in the basis \eqref{bo:eq:basis} have no pole at $t_0$ specializes to the connection $\cN_{t_0}=B_{\sqrt{t_0}}\otimes\bigotimes_j\Sym^{n_j}B_{\mu_j}$ over $\C[y,y^{-1}]$. The local solutions of $\cN_{t_0}$ form the space $\Sol(\beta_{\sqrt{t_0}})\otimes\bigotimes_j\Sym^{n_j}\Sol(\beta_{\mu_j})\cong V\otimes W_m$, where $(I_0,K_0)\leftrightarrow(I,K)$ in each factor. The solution attached to $f\otimes\bigotimes_jg_j^{n_j}$ sends the element \eqref{bo:eq:basis} to $\theta^af\cdot\prod_jg_j^{\,n_j-r_j}(\theta g_j)^{r_j}$.

Finally, let $u(y,t)=K_0(\sqrt{ty})\,G_m(y)$, where $\sqrt{ty}=\sqrt t\sqrt y$ with principal branches. Let $s_u$ be the $\Kt[y,y^{-1}]$-linear map from $\cN$ to meromorphic functions on $\{(y,t):y,t\in\C\setminus(-\infty,0]\}$ given on the basis \eqref{bo:eq:basis} by
\[
s_u\Bigl(\theta^ab_t\otimes\bigotimes_jb_j^{\,n_j-r_j}(\theta b_j)^{r_j}\Bigr)=\theta^aK_0(\sqrt{ty})\cdot\prod_jI_0(\mu_j\sqrt y)^{n_j-r_j}\bigl(\theta I_0(\mu_j\sqrt y)\bigr)^{r_j}.
\]
If the coefficients of $\lambda\in\cN$ have no pole at $t_0\in\C\setminus(-\infty,0]$, then $s_u(\lambda)(\cdot,t_0)$ is the value at $\lambda_{t_0}$ of the solution of $\cN_{t_0}$ attached to $K_0(\sqrt{t_0y})\otimes\bigotimes_jI_0(\mu_j\sqrt y)^{n_j}$; we call this solution $s_u$ as well. In particular $s_u$ commutes with $\partial_y$, and $s_u(\eta)=u$.

\begin{lemma}\label{bo:lem:N}
The connection $\cN$ has the following properties.
\begin{enumerate}
\item[(a)] $\cN$ is irreducible, and $s_u$ is injective with image $\Kt[y,y^{-1}]\langle\partial_y\rangle\cdot u$. Under $\Sol(\cN_{t_0})\cong V\otimes W_m$, the solution $s_u$ of $\cN_{t_0}$ corresponds to $K\otimes w$, where $w$ is the highest weight vector of Lemma~\ref{bo:lem:weight}.
\item[(b)] The lattice $\cN_0$ is $\theta$-stable, and the residue $\mathrm{Res}$, that is, the endomorphism induced by $\theta$ on $\cN_0/y\cN_0\cong\Kt\otimes(V\otimes W_m)$, is $\rho(X)$ for a nonzero nilpotent $X\in\sltwo$. In particular $\mathrm{Res}$ is nilpotent and has one Jordan block of size $d+1$ for each irreducible summand $V_d$ of $V\otimes W_m$. The local monodromy of $\cN_{t_0}$ at $y=0$ is $\exp(-\pi i\rho(e))$ on $V\otimes W_m$, and $\dim\ker(T_0-1)=2\sum_{d\ge1}a_d+a_0$.
\item[(c)] At $y=\infty$, every formal component of $\cN$ is $E^{c\sqrt y}\otimes\mathcal R$ with $\mathcal R$ regular and $c=\pm\sqrt t+c'$, where $c'$ is as in Corollary~\ref{bo:cor:M}(c). In particular $c\ne0$, all slopes equal $1/2$, and $\irr_\infty(\cN)=N$.
\item[(d)] There is a unique additive map $\nabla_t:\cN\to\cN$ with $s_u\circ\nabla_t=\partial_t\circ s_u$. It commutes with $\partial_y$, satisfies $\nabla_t(q\lambda)=(\partial_tq)\lambda+q\nabla_t\lambda$ for $q\in\Kt[y,y^{-1}]$, and maps $\cN_0$ into itself.
\end{enumerate}
\end{lemma}

\begin{proof}
(a) Since $t$ is transcendental over $\C$, the squared masses $t,\mu_1^2,\dots,\mu_k^2$ are pairwise distinct and nonzero. By Proposition~\ref{bo:prop:galois}, applied over $\overline{\Kt}(y)$ with the same proof, the Galois group of $B_t\oplus\bigoplus_jB_{\mu_j}$ is $\SL_2^{k+1}$. (Lemma~\ref{bo:lem:bessel}(d) holds over $\overline{\Kt}(y)$: the substitution $y\mapsto y/t$ transforms $\beta_t$ into $\beta_1$, which is defined over $\C$, and a differential Galois group does not change when the algebraically closed field of constants is enlarged.) It acts on the solution space of $\cN$ through the irreducible representation $V\otimes\bigotimes_j\Sym^{n_j}V$, so $\cN\otimes\Kt(y)$ is irreducible. Lemma~\ref{bo:lem:simple}, applied to the nonzero map $s_u$, gives the first assertion. The second holds because $K_0\leftrightarrow K$ and $I_0\leftrightarrow I$.

(b) In $B_\mu$ we have $\theta(b_\mu)=\theta b_\mu$ and $\theta(\theta b_\mu)=\tfrac{\mu^2}4y\,b_\mu$. By the Leibniz rule, $\theta$ maps each element \eqref{bo:eq:basis} to a $\Kt[y]$-linear combination of such elements. Modulo $y$ it acts on $\Sym^{n}$ by $b^{n-r}(\theta b)^r\mapsto(n-r)\,b^{n-r-1}(\theta b)^{r+1}$. This is the action of the nilpotent $X\in\sltwo$ with $X(b)=\theta b$ and $X(\theta b)=0$ on $V=\langle b,\theta b\rangle$, extended to $V\otimes\bigotimes_j\Sym^{n_j}V$. The Jordan blocks of $\rho(X)$ correspond to the irreducible summands, as in Lemma~\ref{bo:lem:jordan}. The monodromy follows from Lemmas~\ref{bo:lem:bessel}(b) and \ref{bo:lem:jordan}, applied to $V\otimes W_m=V_1\otimes\bigotimes_jV_{n_j}$, and the dimension from Lemma~\ref{bo:lem:count}(b).

(c) Here $\pm\sqrt t+c'\ne0$ because $\sqrt t\notin\C$. A formal connection all of whose components have slope $1/2$ has irregularity equal to half its rank, which is $N$.

(d) Since $K_0(\sqrt{ty})$ depends only on $ty$, we have $\partial_tK_0(\sqrt{ty})=t^{-1}\theta K_0(\sqrt{ty})$ and $\partial_t\,\theta K_0(\sqrt{ty})=t^{-1}\theta^2K_0(\sqrt{ty})=\tfrac y4K_0(\sqrt{ty})$. The other factors do not depend on $t$. Hence the map defined on the basis \eqref{bo:eq:basis} by
\[
\nabla_t(b_t\otimes\xi)=t^{-1}\,\theta b_t\otimes\xi,\qquad\nabla_t(\theta b_t\otimes\xi)=\tfrac y4\,b_t\otimes\xi,
\]
and extended by the Leibniz rule (with $\partial_t$ acting on the coefficients), satisfies $s_u\circ\nabla_t=\partial_t\circ s_u$ and preserves $\cN_0$. Uniqueness follows from the injectivity of $s_u$. So does the commutation with $\partial_y$, since $\partial_t$ and $\partial_y$ commute on functions and $s_u$ is $\partial_y$-equivariant.
\end{proof}

\subsection{Middle cohomology}
Let $H^1(\cN)=\cN\,dy/\partial_y\cN\,dy$ be the algebraic de Rham cohomology of $\cN$ on $\Gm$ over $\Kt$, and let $H^1_{\rm loc,0}(\cN)$ be the corresponding cohomology of the formal completion $\cN\otimes \Kt((y))$. Define
\[
H^1_{\rm mid}(\cN)=\ker\bigl(H^1(\cN)\to H^1_{\rm loc,0}(\cN)\bigr).
\]
By the exact sequence in the proof below, this agrees with the image of compactly supported cohomology, that is, with the middle cohomology (compare \cite{FSY22}).

Let $\Lambda=\cN_0\otimes_{\Kt[y]}\Kt[[y]]\subseteq\cN\otimes\Kt((y))$. By Lemma~\ref{bo:lem:N}(b), $\Lambda$ is a $\theta$-stable lattice, and on $y^k\Lambda/y^{k+1}\Lambda\cong\Lambda/y\Lambda$ the operator $\theta$ acts as $k+\mathrm{Res}$. Since $\mathrm{Res}$ is nilpotent, this is invertible for $k\ne0$.

\begin{proposition}\label{bo:prop:dimmid}
We have $\dim_{\Kt}H^1(\cN)=N$ and $\dim_{\Kt}H^1_{\rm mid}(\cN)=o(m)$. The derivation $\nabla_t$ induces a connection (Gauss--Manin connection) on $H^1(\cN)$ that preserves $H^1_{\rm mid}(\cN)$.
\end{proposition}

\begin{proof}
Work over $\overline{\Kt}$; the dimensions do not change. Since $\cN$ is irreducible of rank $2N\ge2$, $H^0(\Gm,\cN)=0$, and dually $H^2_c(\Gm,\cN)\cong H^0(\Gm,\cN^\vee)^\vee=0$. The Euler characteristic formula of Deligne and Malgrange \cite{Mal74}, \cite[Ch.~IV]{Mal91} gives
\[
\chi(\Gm,\cN)=\rk\cN\cdot\chi(\Gm)-\irr_0(\cN)-\irr_\infty(\cN)=-N .
\]
Hence $\dim H^1(\cN)=N$. The local cohomology at $\infty$ vanishes, because $\cN$ has no slope-$0$ part there. At $0$, since $\partial_y=y^{-1}\theta$, the cohomology $H^1_{\rm loc,0}(\cN)$ is the cokernel of $\theta$ on $\cN\otimes\Kt((y))=\bigcup_k y^{-k}\Lambda$. As $\theta$ is bijective on $y\Lambda$ and on $y^{-k}\Lambda/\Lambda$ for $k\ge1$, this cokernel is $\operatorname{coker}\mathrm{Res}$. Its dimension is the number of Jordan blocks of $\mathrm{Res}$, which by Lemma~\ref{bo:lem:N}(b) and Lemma~\ref{bo:lem:count}(b) is $2\sum_{d\ge1}a_d+a_0$. The long exact sequence
\[
H^1_c(\Gm,\cN)\to H^1(\Gm,\cN)\to H^1_{\rm loc,0}(\cN)\oplus H^1_{\rm loc,\infty}(\cN)\to H^2_c(\Gm,\cN)=0
\]
shows that $H^1(\cN)\to H^1_{\rm loc,0}(\cN)$ is surjective. Hence, by Lemma~\ref{bo:lem:count}(b),
\[
\dim H^1_{\rm mid}(\cN)=N-\Bigl(2\sum_{d\ge1}a_d+a_0\Bigr)=o(m).
\]
The map $\nabla_t$ of Lemma~\ref{bo:lem:N}(d) commutes with $\partial_y$. Extended to $\cN\otimes \Kt((y))$ by letting $\partial_t$ act on coefficients, it still commutes with $\partial_y$. So it acts on $H^1(\cN)$ and on $H^1_{\rm loc,0}(\cN)$, by $[\lambda\,dy]\mapsto[\nabla_t\lambda\,dy]$, and the restriction map commutes with it. The kernel $H^1_{\rm mid}(\cN)$ of the restriction map is therefore $\nabla_t$-stable.
\end{proof}

\begin{lemma}\label{bo:lem:localexact}
For every $\lambda\in\cN_0$ we have $[\lambda\,dy]\in H^1_{\rm mid}(\cN)$. In particular $\omega:=[\eta\,dy]\in H^1_{\rm mid}(\cN)$, and $\nabla_t^j\eta\in\cN_0$ for all $j\ge0$.
\end{lemma}

\begin{proof}
Since $\Lambda$ is $y$-adically complete and separated, and $\theta$ acts on $y^k\Lambda/y^{k+1}\Lambda$ as the invertible map $k+\mathrm{Res}$ for $k\ge1$, the map $\theta:y\Lambda\to y\Lambda$ is bijective. Given $\lambda\in\cN_0$, choose $F\in y\Lambda$ with $\theta F=y\lambda$. Then $\partial_yF=\lambda$, so $\lambda\,dy$ is exact in $\cN\otimes \Kt((y))\,dy$, and the image of $[\lambda\,dy]$ in $H^1_{\rm loc,0}(\cN)$ vanishes. Finally, $\eta\in\cN_0$, and $\nabla_t$ preserves $\cN_0$ by Lemma~\ref{bo:lem:N}(d).
\end{proof}

\subsection{The period}

\begin{lemma}\label{bo:lem:cycle}
Assume $n\ge2$. Let $\lambda\in\cN_0$ with $[\lambda\,dy]=0$ in $H^1(\cN)$. Then, for every real $t_0>R^2$ outside a finite set, the integral $\int_0^\infty s_u(\lambda)(y,t_0)\,dy$ converges absolutely and equals $0$.
\end{lemma}

\begin{proof}
Write $\lambda=\partial_yF$ with $F\in\cN$. Let $t_0>R^2$ be such that the coefficients of $\lambda$ and $F$ in the basis \eqref{bo:eq:basis} have no pole at $t_0$, and specialize the identity $\lambda=\partial_yF$ to $\cN_{t_0}$. Then $s(\lambda)=\partial_y\bigl(s(F)\bigr)$ for every local solution $s$ of $\cN_{t_0}$. By Lemma~\ref{bo:lem:N}(a),(b), Lemma~\ref{bo:lem:jordan} and Lemma~\ref{bo:lem:weight}, the solution $s_u$ of $\cN_{t_0}$ lies in $\im(T_0-1)=e\,(V\otimes W_m)$. Choose a local solution $s_w$ with $(T_0-1)s_w=s_u$, and let $\gamma_\varepsilon$ be the circle $|y|=\varepsilon$, traversed once counterclockwise from $\varepsilon$ to $\varepsilon$; analytic continuation of $s_w$ along $\gamma_\varepsilon$ gives $T_0s_w$.

Each element \eqref{bo:eq:basis} is sent by $s_u$ to a function bounded by $C(1+y)^{n+1}e^{-(\sqrt{t_0}-R)\sqrt y}$ for $y\ge1$, with $C$ depending on $t_0$. Indeed, $|K_0(z)|,|zK_1(z)|\le C'e^{-z}(1+z)$ for $z\ge\sqrt{t_0}$, and $|I_0(z)|,|zI_1(z)|\le(1+|z|)e^{|z|}$ for all $z$. Since $F$ has coefficients in $\C[y,y^{-1}]$, it follows that $s_u(F)(Y)\to0$ as $Y\to+\infty$, and that $\int_1^\infty|s_u(\lambda)|\,dy<\infty$. Near $y=0$, every local solution $s$ sends each element \eqref{bo:eq:basis} to a sum of products of $n+1$ functions of the form $\theta^ag$, with $g$ a Bessel solution and $a\in\{0,1\}$. Each such function is $O(1+|\log y|)$ on bounded sectors, so $s(\lambda)=O\bigl((1+|\log y|)^{n+1}\bigr)$ there, as $\lambda$ has coefficients in $\C[y]$. In particular $\int_0^1|s_u(\lambda)|\,dy<\infty$, and $\int_{\gamma_\varepsilon}s_w(\lambda)\,dy\to0$ as $\varepsilon\to0$.

For $0<\varepsilon<1$ we now compute
\[
\int_\varepsilon^\infty s_u(\lambda)\,dy+\int_{\gamma_\varepsilon}s_w(\lambda)\,dy
=-s_u(F)(\varepsilon)+\bigl((T_0-1)s_w\bigr)(F)(\varepsilon)=0 .
\]
In other words, $\lambda\,dy$ integrates to zero over the rapid decay cycle $[\varepsilon,\infty)\otimes s_u+\gamma_\varepsilon\otimes s_w$, whose boundary vanishes because $(T_0-1)s_w=s_u$. Letting $\varepsilon\to0$ gives $\int_0^\infty s_u(\lambda)\,dy=0$.
\end{proof}

\begin{proposition}\label{bo:prop:upper}
For $n\ge2$ and $m\in(\C^\times)^n$ we have $\ord\varpi_m\le o(m)$.
\end{proposition}

\begin{proof}
By Lemma~\ref{bo:lem:localexact} and Proposition~\ref{bo:prop:dimmid}, the $o(m)+1$ classes $\nabla_t^j\omega=[\nabla_t^j\eta\,dy]$, $0\le j\le o(m)$, lie in the $\Kt$-vector space $H^1_{\rm mid}(\cN)$ of dimension $o(m)$. Hence there are $\alpha_0,\dots,\alpha_{o(m)}\in\Kt$, not all zero, such that $\lambda:=\sum_j\alpha_j\nabla_t^j\eta$ satisfies $[\lambda\,dy]=0$ in $H^1(\cN)$. By Lemma~\ref{bo:lem:localexact}, $\lambda\in\cN_0$. By Lemma~\ref{bo:lem:N}(d), $s_u(\nabla_t^j\eta)=\partial_t^ju$. For real $t>R^2$, differentiation under the integral sign in \eqref{bo:eq:repK} is justified by bounds that are locally uniform in $t$, and gives
\[
\partial_t^j\varpi_m(t)=\frac12\int_0^\infty\partial_t^ju(y,t)\,dy .
\]
Hence, by Lemma~\ref{bo:lem:cycle},
\[
\sum_j\alpha_j(t)\,\partial_t^j\varpi_m(t)=\frac12\int_0^\infty s_u(\lambda)(y,t)\,dy=0
\]
for all real $t>R^2$ outside a finite set. By analytic continuation, $\sum_j\alpha_j\partial_t^j\varpi_m=0$ identically.
\end{proof}

\section{The lower bound}\label{bo:sec:lower}

\subsection{Facts about Fourier transforms}
Let $A_y=\C[y]\langle\partial_y\rangle$ and $A_s=\C[s]\langle\partial_s\rangle$. The isomorphism $A_y\to A_s$, $y\mapsto-\partial_s$, $\partial_y\mapsto s$, turns an $A_y$-module $\mathcal K$ into an $A_s$-module $\Four(\mathcal K)$. We use the following standard facts about holonomic $\cD$-modules on $\A^1$.

\begin{enumerate}
\item[(F1)] $\Four$ is an exact equivalence of categories, since it comes from a ring isomorphism. It preserves holonomicity and simplicity, and $\Four(\cO_{\A^1})\cong\delta_0$, $\Four(\delta_0)\cong\cO_{\A^1}$.
\item[(F2)] For an irreducible connection $L$ on a dense open set $j:U\hookrightarrow\A^1$, the intermediate extension $j_{!*}L$ is simple. Every holonomic module whose restriction to $U$ has composition factors $L_1,\dots,L_r$ has composition factors $j_{!*}L_1,\dots,j_{!*}L_r$ together with modules supported on $\A^1\setminus U$ \cite{HTT08}.
\item[(F3)] If $c$ is a regular singular point of $L$ with unipotent local monodromy $T$, the vanishing cycles of the \emph{intermediate extension} satisfy
\[
\phi_c(j_{!*}L)\cong\im(T-1),\qquad
\mu_c:=\dim\phi_c=\rk L-\dim\ker(T-1).
\]
Here is the local linear algebra, fixing the extension convention. A regular holonomic germ is described by nearby and vanishing cycle spaces $\Psi,\Phi$ and maps $\mathrm{can}:\Psi\to\Phi$, $\mathrm{var}:\Phi\to\Psi$ with $T_\Psi-1=\mathrm{var}\,\mathrm{can}$ and $T_\Phi-1=\mathrm{can}\,\mathrm{var}$; see \cite[Remark~5.7, p.~326]{Sab08}. Having no quotient supported at $c$ makes $\mathrm{can}$ surjective, and having no submodule supported there makes $\mathrm{var}$ injective. These are precisely the intermediate-extension conditions. Thus $\mathrm{var}$ identifies $\Phi$ with $\im(T_\Psi-1)$ and identifies their monodromies. In particular, a unipotent block of size $b$ contributes a block of size $b-1$, with a size-$1$ block contributing nothing.
\item[(F4)] (Rank of the Fourier transform.) Let $\mathcal K$ be holonomic, with regular singularities at finite distance. If $\rk_q$ is the rank of the slope-$q$ part of its formal completion at infinity, then
\[
\rk\Four(\mathcal K)
 =\sum_{c\in\A^1}\dim\phi_c(\mathcal K)
  +\sum_{q>1}(q-1)\rk_q.
\]
This is obtained by taking ranks in the formal stationary-phase theorem \cite[\S1, theorem ``Formal stationary phase'']{GL04}: the finite regular contribution is the vanishing-cycle space by \cite[Remark~5.7]{Sab08}, and the contribution from slope $q>1$ has rank $(q-1)\rk_q$ by \cite[Prop.~3.14(iii)]{BE04}, or \cite[Remarks~5.11(1)]{Sab08}. Only the case with all slopes at infinity at most $1$ is needed below.
\item[(F5)] (Stationary phase, with extension convention.) If $0$ is the only finite singularity of $\mathcal K$, its Fourier transform at $s=\infty$ is the direct sum of the local transform of the \emph{holonomic germ} $\mathcal K_0$ and the transform of the slope-$>1$ part at infinity; see \cite[\S1, formal stationary-phase theorem]{GL04}. For a regular holonomic germ at $0$, the first summand is the regular connection corresponding to $(\phi_0(\mathcal K),T)$, by \cite[Remark~5.7, p.~326]{Sab08}. Thus it is the vanishing-cycle monodromy, not the whole nearby-cycle monodromy, that is retained. The singularities of $\Four(\mathcal K)$ at finite distance are among $s=0$ and the constants $s=c$ occurring as slope-$1$ exponential factors $E^{cy}$ at the original infinity. The inverse local transform sends a positive slope $q<1$ to $q/(1-q)$; a slope-$0$ contribution remains regular or disappears. These location and slope statements follow from the stationary-phase diagram in \cite[Introduction, p.~310]{Sab08}, its Theorem~5.1 and Remarks~5.3(2), together with the inverse equivalence \cite[Prop.~3.10]{BE04}.
\end{enumerate}

\subsection{The chain of transforms}
For an entire function $h$ of exponential order less than $1$ (for example $|h(y)|\le Ce^{c\sqrt{|y|}}$), let $\widehat h(s)=\int_0^\infty e^{-sy}h(y)\,dy$ for $\Re s>0$. Integration by parts gives
\begin{equation}\label{bo:eq:laplace}
\widehat{yh}=-\partial_s\widehat h,\qquad\widehat{h'}=s\widehat h-h(0).
\end{equation}

\begin{lemma}\label{bo:lem:Lmodule}
Let $\mathcal K\subseteq\cO(\C)$ be the $A_y$-submodule generated by finitely many D-finite entire functions satisfying $|h(y)|\le C_he^{c|y|}$ for some $c\ge0$; in particular $\mathcal K$ is holonomic. Then every $h\in\mathcal K$ satisfies such a bound with $c$ replaced by any $c'>c$, so all $\widehat h$ are holomorphic for $\Re s>c$. Moreover, $\cL=\{\widehat h:h\in\mathcal K\}+\C[s]$ is an $A_s$-module, $\C[s]\cong\cO_{\A^1}$ is a submodule, and $\cL/\C[s]\cong\Four(\mathcal K)$. Consequently, the composition factors of $\cL$ are those of $\Four(\mathcal K)$ together with one copy of $\cO_{\A^1}$.
\end{lemma}

\begin{proof}
The bound for all $h\in\mathcal K$ follows from Cauchy's estimates for derivatives and is preserved by multiplication by polynomials, at the cost of replacing $c$ by $c'>c$. By \eqref{bo:eq:laplace}, $\cL$ is stable under $s$ and $\partial_s$, and $h\mapsto\widehat h\bmod\C[s]$ is a surjective map $\mathcal K\to\cL/\C[s]$ that is semilinear with respect to $A_y\cong A_s$. It is injective: $\widehat h(s)\to0$ as $s\to+\infty$, so if $\widehat h$ is a polynomial then $\widehat h=0$, and then $h=0$.
\end{proof}

Fix $n\ge1$ and $m\in(\C^\times)^n$, and let $G=G_m$, $H$ be as in Lemma~\ref{bo:lem:reps}. Let $M=(\cD\cdot G)|_{\Gm}$ (Corollary~\ref{bo:cor:M}) and $\mathcal K_1=j_{!*}M$.

\begin{lemma}\label{bo:lem:K1}
The module $\Four(\mathcal K_1)$ is simple, of generic rank $R'=\sum_{d\ge1}d\,a_d$. It is a connection on $\Gm$. At $s=\infty$ it is regular singular with unipotent monodromy having one Jordan block of size $d$ for each summand $V_d$, $d\ge1$. At $s=0$ all its slopes are $\le1$.
\end{lemma}

\begin{proof}
Simplicity follows from (F1), (F2) and Corollary~\ref{bo:cor:M}(a). The only singular point of $\mathcal K_1$ in $\A^1$ is $0$, and $\widehat{(\mathcal K_1)}_\infty$ has slopes in $\{0,1/2\}$ by Corollary~\ref{bo:cor:M}(c). By (F4), (F3) and Lemma~\ref{bo:lem:jordan}, the generic rank is $\mu_0=\dim W_m-\sum_da_d=\sum_dda_d$. By (F5) there are no singular points other than $0$ and $\infty$. The germ at $\infty$ is regular with the monodromy of $T_0$ on $\im(T_0-1)$, which has one Jordan block of size $d$ for each Jordan block of size $d+1\ge2$ of $T_0$. The slopes at $0$ are $0/(1-0)=0$ and $\tfrac12/(1-\tfrac12)=1$.
\end{proof}

Let $\iota:\Gm\to\Gm$, $v\mapsto s=1/v$, and set $N_v=\iota^*\bigl(\Four(\mathcal K_1)|_{\Gm}\bigr)$. Then $N_v$ is an irreducible connection of rank $R'$ on $\Gm$. At $v=0$ it is regular singular with the Jordan structure of Lemma~\ref{bo:lem:K1}, and at $v=\infty$ all slopes are $\le1$.

\begin{lemma}\label{bo:lem:Mhat}
The module $\widehat{\cM}:=\Four(j_{!*}N_v)$ is simple of generic rank $o(m)$. Its formal completion at $\tau=\infty$ is regular singular with unipotent monodromy having one Jordan block of size $d-1$ for each summand $V_d$ with $d\ge2$.
\end{lemma}

\begin{proof}
Simplicity follows from (F1) and (F2). The only singular point of $j_{!*}N_v$ in $\A^1_v$ is $0$, and there are no slopes $>1$ at $\infty$. By (F4) and (F3), the generic rank is
\[
\mu_0=R'-\#\{\text{Jordan blocks of }N_v\text{ at }0\}=\sum_{d\ge1}d\,a_d-\sum_{d\ge1}a_d=o(m).
\]
By (F5), the germ at $\infty$ is regular, with the monodromy of $N_v$ at $0$ restricted to the image of $T-1$, which has blocks of sizes $d-1$ for $d\ge2$.
\end{proof}

\begin{proposition}\label{bo:prop:lower}
For $n\ge1$ and $m\in(\C^\times)^n$, the module $\cD_t\cdot\varpi_m$ has $\widehat\cM$ (rescaled by $\tau=t/4$) as a composition factor. All its other composition factors are isomorphic to $\cO_{\A^1}$ or $\delta_0$. In particular $\ord\varpi_m\ge o(m)$.
\end{proposition}

\begin{proof}
\emph{Step 1.} The module $\cD_y\cdot G\subseteq\cO(\C)$ consists of entire functions of growth $O(e^{c\sqrt{|y|}})$. By (F2) its composition factors are $\mathcal K_1$ and copies of $\delta_0$. By Lemma~\ref{bo:lem:Lmodule}, $\cD_s\cdot\widehat G\subseteq\cL_1$, where the composition factors of $\cL_1$ are $\Four(\mathcal K_1)$ (once) and copies of $\cO_{\A^1}$.

\emph{Step 2.} Since $v^{-1}$ is a unit on $\Gm$, the restriction of $\cD_v\cdot H$ to $\Gm$ equals the module generated by $\widehat G(1/v)$, which is $\iota^*\bigl((\cD_s\cdot\widehat G)|_{\Gm}\bigr)$. Its composition factors are therefore among $N_v$ (at most once) and trivial rank one connections. Since $H$ is entire of exponential type, $\cD_v\cdot H\subseteq\cO(\C)$ satisfies the hypothesis of Lemma~\ref{bo:lem:Lmodule}. By (F2), its composition factors on $\A^1$ are among $j_{!*}N_v$ (at most once), $j_{!*}\cO_{\Gm}=\cO_{\A^1}$ and $\delta_0$.

\emph{Step 3.} By Lemma~\ref{bo:lem:Lmodule} and \eqref{bo:eq:repL}, $\cD\cdot\varpi_m$ (in the variable $\tau=t/4$) is contained in a module $\cL_2$ whose composition factors are among $\widehat\cM$ (at most once), $\Four(\cO_{\A^1})=\delta_0$, $\Four(\delta_0)=\cO_{\A^1}$, and $\cO_{\A^1}$.

\emph{Step 4.} Suppose $\widehat\cM$ is not a composition factor of $\cD\cdot\varpi_m$. Consider $E=\C[t,t^{-1}]\cdot(\cD\cdot\varpi_m)$, a module of single-valued analytic functions on the punctured neighbourhood $\{|t|>R^2\}$ of $t=\infty$. By Step 3 it has a filtration $0=E_0\subset E_1\subset\dots\subset E_r=E$ with $E_i/E_{i-1}\cong\cO_{\Gm}$. Choose $e_i\in E_i$ mapping to the horizontal generator of $E_i/E_{i-1}$. Then $\partial_te_i\in\sum_{j<i}\C[t,t^{-1}]\,e_j$, and $e_1,\dots,e_r$ generate $E$ over $\C[t,t^{-1}]$. By induction on $i$, each $e_i$ is a Laurent polynomial: $\partial_te_i$ is a Laurent polynomial, so $e_i$ is a Laurent polynomial plus a multiple of $\log t$, and the multiple is zero because $e_i$ is single-valued. Hence $\varpi_m\in E$ would be a Laurent polynomial. But $G_m$ is not a polynomial: each factor $I_0(\mu\sqrt y)$ is a transcendental entire function of order $1/2$ in $y$ and has infinitely many zeros. Hence infinitely many $c_k(m)$ are nonzero, a contradiction. So $\widehat\cM$ is a composition factor, and $\ord\varpi_m$, the generic rank of $\cD\cdot\varpi_m$, is at least $\rk\widehat\cM=o(m)$.
\end{proof}

\section{Proof of Theorem A}\label{bo:sec:proof}

Let $n\ge2$ and $m\in(\C^\times)^n$. Propositions~\ref{bo:prop:upper} and \ref{bo:prop:lower} give $\ord\varpi_m=o(m)$, which is (i). By Proposition~\ref{bo:prop:lower}, the generic rank of $\cD\cdot\varpi_m$ is $o(m)=\rk\widehat\cM$, so no composition factor $\cO_{\A^1}$ occurs. On a dense open set $U$ of the $t$-line containing a punctured neighbourhood of $\infty$, $\cD\cdot\varpi_m|_U\cong\widehat\cM|_U$. Hence $\cD\cdot\varpi_m$ is irreducible over $\C(t)$, which is (ii), and its local monodromy at $\infty$ is that of $\widehat\cM$, given by Lemma~\ref{bo:lem:Mhat}, which is (iii). Corollary~B follows from Lemma~\ref{bo:lem:count}(c), and Corollary~C from $W_m=V_n$. \qed

\begin{remark}\label{bo:rem:cyclic}
The proof shows more. The cyclic $\cD$-module generated by $\omega$ in $H^1_{\rm mid}(\cN)$ has rank between $\ord\varpi_m$ and $o(m)$, so it is all of $H^1_{\rm mid}(\cN)$. After tensoring with $\C(t)$, that module surjects onto $\C(t)\otimes\cD\cdot\varpi_m\cong\C(t)\otimes\widehat\cM$, and both have dimension $o(m)$. Hence the Gauss--Manin system $H^1_{\rm mid}(\cN)$ is itself irreducible of rank $o(m)$. In particular, the minimal Picard--Fuchs operator of the cohomology class $\omega$ also has order $o(m)$.
\end{remark}

\section{Reported computations and their status}\label{bo:sec:numerics}

The supplied manuscript describes the following computations. They are retained as a record of its numerical claims, not as newly verified computations. The source programs, precise random mass samples, coefficient cutoffs, and complete annihilation certificates were not supplied with the manuscript, and these searches were not independently rerun for this revision. In particular, failure to find a relation within a degree bound is not a proof of unrestricted minimality, and a relation fitted to finitely many coefficients is not by itself a certified differential equation. None of the results in this section is used in the proof of Theorem~A.

The first reported search works modulo the prime $67108859$ and considers
\[
\sum_{r\le\rho,\,\ell\le D}c_{r\ell}z^\ell\partial_z^r
\qquad\text{applied to}\qquad
z^{-1}\varpi_m(1/z)=\sum_kc_k(m)z^k,
\]
for random integer squared masses with prescribed multiplicities.
\begin{itemize}
\item \emph{Generic masses.} For $n=2,3,4,5$, candidate operators of orders $1,2,6,12$ and respective degrees $2,7,27,121$ are reported. The bounded searches reportedly find no candidate of order at most $5$ and degree at most $120$ for $n=4$, or of order at most $11$ and degree at most $160$ for $n=5$.
\item \emph{Other configurations in Table~\ref{bo:tab:conf}.} For each configuration the source reports a candidate of order $o(m)$, and no candidate of order $o(m)-1$ and degree at most $90$.
\item \emph{Local data.} For configurations $[1,1,1]$, $[1,1,1,1]$, $[4]$, $[2,2]$, $[2,1,1]$, $[3,1]$, $[1,1,1,1,1]$, $[5]$, $[2,2,1]$, and $[3,1,1]$, the reported local computations give respectively $1$, $4$, $1$, $2$, $3$, $2$, $5$, $1$, $3$, and $3$ power-series solutions. These numbers match the numbers of blocks predicted in characteristic zero. The reported indicial roots of $z^{-1}\varpi$ are $0^3,1,2,3$ for generic $n=4$ and $0^4,1^2,2^2,3^2,4^2$ for generic $n=5$. Finite-field data do not independently certify the characteristic-zero Jordan structure.
\item \emph{Second implementation.} The source describes a separate program modulo $2^{61}-1$, covering $24$ cases with $2\le n\le5$ and every multiplicity pattern in that range. Its cases reportedly include $(1,2,3)$, $(1,2,3,4)$, $(1,2,3,6)$, $(1,2,3,5)$, $(1,1,2,3)$, and $(1,i,2,3)$. It reports candidates of order $o(m)$, no candidates of order $o(m)-1$ with degree in $z$ at most $100$ in the $\theta_z$ representation, and logarithmic solution counts matching Theorem~A(iii). The independence and outputs of this implementation have not been checked here.
\item \emph{Published configurations with $n=6$.} For the configurations $[1^6]$, $[1^42]$, $[1^33]$, $[1^22^2]$, $[1^24]$, $[123]$, $[15]$, $[24]$, $[3^2]$, and $[6]$, \cite[\S4.5]{LV23} gives orders $29,23,17,18,12,13,8,9,9,5$, respectively. The representation-theoretic formula gives the same values.
\item \emph{Middle-cohomology calculation.} The source also reports exact rational calculations at $t_0=61/5$ for eleven multiplicity patterns with $n\le4$, with a pole-order truncation reported to be stable. The reported dimensions are $N$ for $H^1$ and $o(m)$ for the image of $\cN_0$, with $\nabla_t^j\omega$, $0\le j<o(m)$, spanning that image. A numerical evaluation using $400$ series terms is reported to give cancellation to about $60$ digits for $n\ge2$, and no such cancellation for $n=1$. The truncation-stability certificate and numerical error bounds have not been independently checked.
\end{itemize}
These reports are consistent with the theorem. The proof above, rather than the bounded searches, establishes exact minimality and monodromy. The $n=7$ calculation in \cite[\S4.6]{LV23} is likewise finite-field evidence at selected masses, not a proof of generic characteristic-zero minimality.

\section{Remarks}\label{bo:sec:remarks}

\begin{remark}[The case $n=1$]\label{bo:rem:n1}
For $n=1$, $\varpi_m(t)=(t-m_1^2)^{-1}$ has order $1$, whereas $o(m)=0$: here $W_m=V_1$ has no summand with $d\ge2$. The lower bound still holds, but the upper bound fails. Indeed, by Lemma~\ref{bo:lem:weight} the solution $K\otimes I$ is not in $\im(T_0-1)$, so the ray $[0,\infty)$ cannot be closed to a rapid decay cycle, and the boundary term at $y=0$ produces the pole at $t=m_1^2$.
\end{remark}

\begin{remark}[Singularities]
The singular points of $\widehat\cM$ in $\C^\times$ are among the thresholds and pseudo-thresholds $t=\bigl(\sum_i\varepsilon_im_i\bigr)^2$, $\varepsilon\in\{\pm1\}^n$. These arise in the chain of transforms from the exponential factors $E^{c^2v/4}$ of $N_v$ at $\infty$ (Corollary~\ref{bo:cor:M}(c)). This is the familiar Landau analysis, recovered here from the critical values of the phase. Part~\ref{part:structure}  shows that, for generic masses, every threshold is a genuine singularity, as is $t=0$ for $n\ge3$.
\end{remark}

\begin{remark}[The full Feynman integral]
Theorem~A concerns the maximal cut $\varpi_m$. The Feynman integral itself satisfies an inhomogeneous Picard--Fuchs equation $L\,I=S$ \cite[(1.6)]{LV23}. In our setting, Lemma~\ref{bo:lem:Lmodule} shows how boundary terms produce composition factors $\cO_{\A^1}$ and $\delta_0$, which is the source of such inhomogeneities. We do not study $S$ here; Part~\ref{part:structure}  shows that $S\ne0$, so that the full integral has order $o(m)+1$.
\end{remark}

\begin{remark}[Relation with Calabi--Yau operators]
By Corollary~C, for equal masses the minimal operator has maximal unipotent monodromy at $\infty$, as expected for the Calabi--Yau operators of \cite{BKV15,BFKNS21}. For generic masses and $n\ge4$, the monodromy at $\infty$ has a block of size $n-1$, coming from $V_n$, together with smaller blocks. For $n=4$ the block sizes $3,1,1,1$ are what one expects at a maximal (type III) degeneration of a family of K3 surfaces whose transcendental part has rank $6$; we do not pursue this here.
\end{remark}


\clearpage
\part{Monodromy and affine extensions}\label{part:structure}


\section{Singularities, symmetry, and the full integral}

\subsection{The Lairez--Vanhove conjecture}
For $m=(m_1,\dots,m_n)\in(\C^\times)^n$, the maximal cut of the $(n-1)$-loop banana (sunset) integral in two dimensions is, up to normalization, the torus period
\begin{equation}\label{bs:eq:varpi}
\varpi_m(t)=\sum_{k\ge0}c_k(m)\,t^{-k-1},\qquad c_k(m)=\sum_{|a|=k}\binom{k}{a_1,\dots,a_n}^{2}m_1^{2a_1}\cdots m_n^{2a_n},
\end{equation}
which converges for $|t|>R^2$, $R=\sum_i|m_i|$ \cite{LV23,BFKNS21}. We call the numbers
\[
\mu_\varepsilon=(\varepsilon_1m_1+\dots+\varepsilon_nm_n)^2,\qquad\varepsilon\in\{\pm1\}^n ,
\]
the \emph{thresholds}; in the physics literature, all but $(m_1+\dots+m_n)^2$ are usually called pseudo-thresholds. Since $\mu_\varepsilon=\mu_{-\varepsilon}$, generic masses have $2^{n-1}$ distinct thresholds. For positive masses, $R^2=(m_1+\dots+m_n)^2$ is the \emph{normal threshold}. Lairez and Vanhove computed minimal Picard--Fuchs operators for $n\le7$ and conjectured the following \cite[Conj.~4.0.1]{LV23}. For generic masses:
\begin{enumerate}
\item[(i)] the minimal operator annihilating $\varpi_m$ has order $o_n=2^n-\binom{n+1}{\lfloor(n+1)/2\rfloor}$;
\item[(ii)] its leading coefficient has the form $t^{\alpha(n)}\prod_{\mu}(t-\mu)\,\tilde q(t)$, where $\mu$ runs over the distinct thresholds, and the polynomial $\tilde q$ accounts for the apparent singularities. The non-apparent singularities are $t=0$ and the thresholds.
\end{enumerate}
The same count $o_n$ was proposed in \cite{BFKNS21}, and the formula is also used in \cite[\S3.1]{Van26}.

Part~\ref{part:order}  proved (i), in the following stronger form. Let $\mu_1^2,\dots,\mu_k^2$ be the distinct values among $m_1^2,\dots,m_n^2$, with multiplicities $n_1,\dots,n_k$, let $V=\C^2$ be the standard representation of $\sltwo$ and $V_d=\Sym^dV$, and decompose $W_m=V_{n_1}\otimes\dots\otimes V_{n_k}=\bigoplus_dV_d^{\oplus a_d}$.

\begin{theoremA}
For $n\ge2$ and $m\in(\C^\times)^n$, the minimal operator $L_m$ of $\varpi_m$ has order $o(m)=\sum_{d\ge2}(d-1)a_d$, it is irreducible, and its local monodromy at $t=\infty$ is unipotent with one Jordan block of size $d-1$ for each summand $V_d$ with $d\ge2$. For pairwise distinct squares, $o(m)=o_n$.
\end{theoremA}

The present paper is about (ii) and about the global structure of $L_m$.

\subsection{Results}
We say that $m$ is \emph{generic} if
\[
\text{(G)}\qquad\text{the $2^{n-1}$ numbers $\mu_\varepsilon$, $\varepsilon\in\{\pm1\}^n/\pm1$, are pairwise distinct and nonzero.}
\]
Condition (G) implies that the $m_i^2$ are pairwise distinct. Put
\[
\nu_n=2^{n-1}-\binom{n}{\lfloor n/2\rfloor}.
\]
Then $\nu_2=0$ and $\nu_n\ge1$ for $n\ge3$. Moreover $o_n-2\nu_n=2\binom{n}{\lfloor n/2\rfloor}-\binom{n+1}{\lfloor (n+1)/2\rfloor}$, which equals the Catalan number $\Cat(n/2)$ if $n$ is even and $0$ if $n$ is odd.

\begin{theoremB}
Let $n\ge2$, let $m$ satisfy {\rm(G)}, and let $L_m$ be the minimal operator of $\varpi_m$.
\begin{enumerate}
\item[(i)] The local system of solutions of $L_m$ extends to $\PP^1\setminus(\{0,\infty\}\cup\{\mu_\varepsilon\})$. Its local monodromy is nontrivial at every threshold, and it is nontrivial at $t=0$ if and only if $n\ge3$. Consequently, the non-apparent singular points of $L_m$ in $\C$ are exactly the $2^{n-1}$ thresholds, together with $t=0$ when $n\ge3$. Every other zero of the leading coefficient of $L_m$ is an apparent singularity.
\item[(ii)] At every threshold the local monodromy is a pseudo-reflection. It is a reflection, with eigenvalues $1,\dots,1,-1$, if $n$ is even, and a transvection, unipotent with $\rk(T-1)=1$, if $n$ is odd.
\item[(iii)] At $t=0$ the local monodromy $T_0$ is unipotent, with $(T_0-1)^2=0$ and $\rk(T_0-1)=\nu_n$. Thus it has $\nu_n$ Jordan blocks of size $2$, and $\Cat(n/2)$ blocks of size $1$ if $n$ is even (none if $n$ is odd).
\item[(iv)] If $m_1,\dots,m_n>0$, then near the normal threshold $R^2$
\[
\varpi_m(t)=A(t)+(t-R^2)^{\frac{n-3}2}\,h(t)\cdot\begin{cases}1,& n\text{ even},\\ \log(t-R^2),& n\text{ odd},\end{cases}
\]
with $A,h$ analytic at $R^2$ and $h(R^2)=\kappa_n\ne0$ given in \eqref{bs:eq:kappa}.
\end{enumerate}
\end{theoremB}

Theorem~B proves the description of the singular locus in (ii) of the conjecture; for $n=2$ the point $t=0$ is not singular. The product formula in (ii) asserts in addition that each threshold is a \emph{simple} zero of the leading coefficient, and it involves the multiplicity $\alpha(n)$ of $t=0$. We prove neither of these two statements. Remark~\ref{bs:rem:simple} gives a sufficient condition for simplicity at a threshold. It holds at the normal threshold for $n=3$. In the source-reported computations (Section~\ref{bs:sec:numerics}), which use one sample of generic masses for each $n\le5$ and work modulo a prime, every threshold is a simple zero. These computations also reproduce the degrees $2$, $17$, $98$ of $\tilde q$ found in \cite[\S\S4.2--4.4]{LV23} for $n=3,4,5$.

For the Galois group we allow mass configurations. Let $\mu_1,\dots,\mu_k$ and $n_1,\dots,n_k$ be as above, and let $\Lambda_m$ be the set of $\lambda=(\lambda_1,\dots,\lambda_k)$ with $\lambda_j\in\{n_j,n_j-2,\dots,-n_j\}$ and $\lambda\ne0$. Consider
\[
\text{(G$'$)}\qquad\text{the numbers }\Bigl(\sum_j\lambda_j\mu_j\Bigr)^2,\ \lambda\in\Lambda_m/\pm1,\ \text{are pairwise distinct and nonzero.}
\]
Condition (G$'$) is (G) when all $n_j=1$, and it holds automatically when all masses are equal.

\begin{theoremC}\label{bs:thm:biggroup}
Let $n\ge3$ and let $m$ satisfy {\rm(G$'$)}. The differential Galois group of $L_m$ over $\C(t)$ is $\Sp(o(m),\C)$ if $n$ is odd, and $\Ort(o(m),\C)$ if $n$ is even.
\end{theoremC}

For $n=2$ and generic masses the group is $\Ort(1)=\{\pm1\}$, since $\varpi_m=((t-\mu_+)(t-\mu_-))^{-1/2}$ with $\mu_\pm=(m_1\pm m_2)^2$.

\begin{corollaryD}
Under the hypotheses of Theorem~C, put $o=o(m)$. If $n$ is odd, the functions $\varpi_m^{(j)}$, $0\le j<o$, are algebraically independent over $\C(t)$. If $n$ is even, the ideal of polynomial relations among them over $\C(t)$ is generated by a single irreducible quadratic relation, which comes from the invariant bilinear form.
\end{corollaryD}

The number of closed walks of length $2k$ on the $d$-dimensional diamond lattice is $c_k(1,\dots,1)$ with $n=d+1$ (Section~\ref{bs:sec:galois}), so its lattice Green function is an equal-mass banana function \cite[(40)--(41)]{Gut10}. Hassani, Maillard and Zenine \cite{HMZ25} conjecture that the minimal number $N_{\rm var}$ of variables in a representation of a power series as the diagonal of a rational function is $n_{\log}+2$, where $n_{\log}$ is the highest power of the logarithm in the solutions at the origin \cite[\S1.2]{HMZ25}. They also conjecture that the differential Galois group of the minimal operator is orthogonal for $N_{\rm var}$ even and symplectic for $N_{\rm var}$ odd \cite[\S1.5]{HMZ25}, and they check the first conjecture for diamond lattices of dimension $d=3,4,5$, where $N_{\rm var}=d+1$.

\begin{corollaryE}
Let $d\ge2$ and let $P_d(x)=\sum_kW_{2k}x^k$, where $W_{2k}$ is the number of closed walks of length $2k$ on the $d$-dimensional diamond lattice. The differential Galois group of the minimal operator of $P_d$ over $\C(x)$ is $\Ort(d)$ if $d$ is odd and $\Sp(d)$ if $d$ is even. The highest power of the logarithm among the solutions at $x=0$ is $n_{\log}=d-1$. Thus the group is orthogonal exactly when $n_{\log}+2=d+1$ is even. This is the dichotomy predicted in \cite{HMZ25} through their conjectural relation $N_{\rm var}=n_{\log}+2$.
\end{corollaryE}

Finally, consider the full Feynman integral. For $m_1,\dots,m_n>0$ and real $t<R^2$ put
\begin{equation}\label{bs:eq:I}
I_m(t)=\frac12\int_0^\infty I_0\bigl(\sqrt{ty}\bigr)\prod_{i=1}^nK_0\bigl(m_i\sqrt y\bigr)\,dy ,
\end{equation}
where $I_0(\sqrt{ty})=J_0(\sqrt{|t|y})$ for $t<0$. In position space one finds that $(2\pi)^{1-n}I_m(-p^2)$ is the Euclidean banana integral in two dimensions with external momentum $p$ (Section~\ref{bs:sec:full}).

\begin{theoremF}\label{bs:thm:inhomogeneous}
Let $n\ge2$ and $m_1,\dots,m_n>0$, and let $L$ be any operator of order $o(m)$ annihilating $\varpi_m$. Then $LI_m=S_m$ as an identity of meromorphic functions on $(-\infty,R^2)$, away from poles of the coefficients, for a nonzero rational function $S_m\in\C(t)$. Consequently the minimal operator of $I_m$ has order $o(m)+1$; for generic masses this is $2^n-\binom{n+1}{\lfloor(n+1)/2\rfloor}+1$.
\end{theoremF}

The inhomogeneous equation $LI=S$ appears in \cite[(1.6)]{LV23}, where $S$ is not studied for banana integrals, and Part~\ref{part:order} did not study it either. In \cite{BFKNS21} the full integral is described as a special solution of such an equation, and for equal masses the inhomogeneity is given explicitly and is nonzero \cite[(3.6)]{BFKNS21}. Theorem~F shows that $S$ never vanishes, for arbitrary positive masses, so that the full integral always needs exactly one more order than its maximal cut. For $n=2$ and $L=\Delta\partial_t+\frac12\Delta'$, $\Delta=(t-\mu_+)(t-\mu_-)$, one has $S_m=-1$ (Example~\ref{bs:ex:n2}). For $n=3$ the source reports numerical evidence that $S_m$ is a nonconstant rational function whose coefficients involve logarithms of mass ratios (Section~\ref{bs:sec:numerics}).

\begin{theoremG}[Affine extension and full-integral jets]\label{bs:thm:affineintro}
Let $n\ge2$, let all masses be positive, and put $o=o(m)$. Write $L=L_m$, $S=LI_m\in\C(t)^\times$, and $A=(\partial_t-S'/S)L$, the minimal homogeneous operator of $I_m$. Let $G_m$ be the differential Galois group of $L$ on its solution space $V$.
\begin{enumerate}
\item The differential Galois group of $A$ is the full affine group
\[
V\rtimes G_m=\left\{\begin{pmatrix}g&b\\0&1\end{pmatrix}:g\in G_m,\ b\in V\right\}.
\]
Under {\rm(G$'$)} and $n\ge3$, this is $\mathbb G_a^o\rtimes\Sp(o)$ for odd $n$ and $\mathbb G_a^o\rtimes\Ort(o)$ for even $n$.
\item If $P\subset E$ are compatible Picard--Vessiot fields for $L$ and $A$, respectively, then
\[
E=P(I_m,I_m',\ldots,I_m^{(o-1)})
\]
is a purely transcendental extension of $P$ of degree $o$. In particular these $o$ full-integral jets are algebraically independent even after adjoining all solutions of the maximal-cut equation and their derivatives.
\item At $t=\infty$, $A$ is regular singular and its unipotent Jordan partition is obtained from that of $L$ by replacing its unique block of size $n-1$ with a block of size $n$, leaving every other block unchanged.
\end{enumerate}
\end{theoremG}

Explicit affine monodromy matrices for the four-loop equal-mass full integral already appear in \cite[\S3.3, (3.98)]{BDFKN22}. The affine-extension principle used here is standard differential Galois theory; we give an elementary proof and explicitly relate it to Bertrand's treatment of nonsplit inhomogeneous equations \cite[\S1.1, (N.3), and Thm.~1.1]{Ber01}. The application uses the all-mass irreducibility from Part~\ref{part:order} and the nonzero additional logarithm in Theorem~F. We present it as an additional consequence of the present framework, without claiming that the general principle or its banana application resolves a separately identified published conjecture. The stronger assertions about the full integral do not require {\rm(G$'$)}; that condition is used only to name $G_m$ as a classical group.

\subsection{Strategy}
Part~\ref{part:order} expresses $\varpi_m$ as a Laplace transform of an inversion of a Laplace transform of the product of Bessel functions $G_m(y)=\prod_iI_0(m_i\sqrt y)$. This yields a simple holonomic module $\Mhat$ of rank $o(m)$ that agrees with the $\cD$-module of $\varpi_m$ away from $t=0$. The singularities of $\Mhat$ come from critical values of phases. The phases are the exponential factors $e^{\pm c\sqrt y}$ of the Bessel product at $y=\infty$, with $c=\sum\pm m_i$. Stationary phase at finite distance \cite{Ari08,GL04,BE04} turns each pair $\pm c$ into a one-dimensional space of vanishing cycles of $\Mhat$ at $t=c^2$. The regular part that is left over produces the Jordan blocks at $t=0$. A homothety argument shows that the monodromy on these vanishing cycles depends only on $n$, and a Laplace-method computation at the normal threshold identifies it. For the Galois group, Poincaré duality on middle de Rham cohomology gives an invariant form of the right parity, and a theorem of Kazhdan--Margulis, Gabber and Beukers--Heckman \cite[Thm.~1.5]{Kat90} identifies the group. For the full integral, the cohomological relation of Part~\ref{part:order} applied to a non-closed integration path produces a rational boundary term, and a large-momentum asymptotic shows that this term is not zero.

\subsection*{Conventions}
We use the conventions of Part~\ref{part:order}. In particular $\theta=y\,d/dy$, $\Four$ is the Fourier transform of $\cD$-modules on $\A^1$ given by $y\mapsto-\partial_s$, $\partial_y\mapsto s$, and $E^{\varphi}$ denotes the rank one formal connection with solution $e^{\varphi}$. For a holonomic module $K$ on a curve and a point $x$, $\phi_x(K)$ denotes the formal vanishing cycles of $K$ at $x$, in the sense of Arinkin \cite[\S3.1, Cor.~4.2]{Ari08}. It is a connection on the punctured formal disc at $x$, and it determines the formal microlocalization of $K$ at $x$ \cite[Prop.~3.3]{Ari08}. On irregular summands, and on regular summands without the eigenvalue $1$, it agrees with the restriction of $K$ to the punctured formal disc; on regular unipotent summands it is the space of vanishing cycles with the operator induced by the monodromy \cite[\S3.3]{Ari08}. We call the latter operator the monodromy of $\phi_x(K)$; see (F8). A pseudo-reflection is an element $g\in\GL(V)$ with $\rk(g-1)=1$; it is a reflection if it is semisimple with eigenvalue $-1$, and a transvection if it is unipotent.

\section{Recollections}\label{bs:sec:recall}

\subsection{Inputs from the banana-order argument}\label{bs:sec:partI}
Throughout, $n\ge2$. We use the following results of Part~\ref{part:order}, with cross-references to that part.
\begin{enumerate}
\item[(P1)] (Lemma~\ref{bo:lem:reps} and its proof.) We have
\[
G_m(y)=\prod_iI_0(m_i\sqrt y)=\sum_k\frac{c_k(m)(y/4)^k}{(k!)^2},\qquad |G_m(y)|\le e^{R\sqrt{|y|}}.
\] The function $H(v)=\sum_kc_k(m)v^k/(4^kk!)$ is entire of exponential type at most $R^2/4$, $H(v)=v^{-1}\widehat G(1/v)$ with $\widehat G(s)=\int_0^\infty e^{-sy}G_m(y)\,dy$, and for $t>R^2$
\[
\varpi_m(t)=\frac14\int_0^\infty e^{-tv/4}H(v)\,dv=\frac12\int_0^\infty K_0\bigl(\sqrt{ty}\bigr)G_m(y)\,dy .
\]
\item[(P2)] (Corollary~\ref{bo:cor:M}, Lemmas~\ref{bo:lem:K1} and~\ref{bo:lem:Mhat}.) $M=(\cD\cdot G_m)|_{\Gm}\cong\bigotimes_j\Sym^{n_j}B_{\mu_j}$ is an irreducible connection of rank $N=\dim W_m$, where $B_\mu$ is the Bessel connection of $\theta^2-\frac{\mu^2}{4}y$. Let $\cK_1=j_{!*}M$. Then $\Four(\cK_1)$ is simple of rank $R'=\sum_{d\ge1}d\,a_d$ and smooth on $\Gm$. Let $N_v=\iota^*(\Four(\cK_1)|_{\Gm})$ with $\iota(v)=1/v$. Then $\Mhat=\Four(j_{!*}N_v)$, a module on the $\tau$-line, is simple of rank $o(m)$. It is regular singular at $\tau=\infty$ with unipotent monodromy having one Jordan block of size $d-1$ for each $V_d$, $d\ge2$. We view $\Mhat$ as a module on the $t$-line via $\tau=t/4$.
\item[(P3)] (Proposition~\ref{bo:prop:lower} and Section~\ref{bo:sec:proof}.) The composition factors of the $\C[t]\langle\partial_t\rangle$-module $\cD\cdot\varpi_m$ are $\Mhat$, once, and copies of $\delta_0$. Hence $\cD\cdot\varpi_m|_U\cong\Mhat|_U$ for every open set $U\subseteq\A^1$ with $0\notin U$.
\item[(P4)] (Section~\ref{bo:sec:upper}.) Over $\Kt=\C(t)$ there is a connection $\cN=B_t\otimes\bigotimes_j\Sym^{n_j}B_{\mu_j}$ on $\Gm$ of rank $2N$, with the basis~\eqref{bo:eq:basis} over $\Kt[y,y^{-1}]$, its $\Kt[y]$-span $\cN_0$, an element $\eta\in\cN_0$, an injective solution map $s_u$ with $s_u(\eta)=K_0(\sqrt{ty})G_m(y)$, and a derivation $\nabla_t$ with $s_u\circ\nabla_t=\partial_t\circ s_u$ that commutes with $\partial_y$ and preserves $\cN_0$ (Lemma~\ref{bo:lem:N}). The middle cohomology $H^1_{\rm mid}(\cN)$ has dimension $o(m)$ and is preserved by $\nabla_t$, and $\omega=[\eta\,dy]\in H^1_{\rm mid}(\cN)$ (Proposition~\ref{bo:prop:dimmid}, Lemma~\ref{bo:lem:localexact}). A relation $\sum_{j\le o(m)}\alpha_j\nabla_t^j\omega=0$ in $H^1(\cN)$, with $\alpha_j\in\Kt$ not all zero, yields the operator $\sum_j\alpha_j\partial_t^j$, which annihilates $\varpi_m$ (Lemma~\ref{bo:lem:cycle}, Proposition~\ref{bo:prop:upper}).
\item[(P5)] (Remark~\ref{bo:rem:cyclic}.) $H^1_{\rm mid}(\cN)$ is generated by $\omega$ as a $\Kt\langle\nabla_t\rangle$-module, and $\omega\mapsto\varpi_m$ induces an isomorphism of $\Kt$-connections $H^1_{\rm mid}(\cN)\cong\Kt\otimes\cD\cdot\varpi_m$. In particular $H^1_{\rm mid}(\cN)$ is irreducible of rank $o(m)$.
\end{enumerate}
We also use the facts (F1)--(F5) on Fourier transforms listed in Part~\ref{part:order}, Section~\ref{bo:sec:lower}.

\subsection{Fourier transforms near finite points}\label{bs:sec:finite}
We need the following standard facts on local Fourier transforms \cite{BE04,GL04,Sab08,Ari08}. Let $K$ be a holonomic $\cD$-module on $\A^1_v$ such that the formal completion $\widehat K_\infty$ at $v=\infty$ has slopes $\le1$. By the Levelt--Turrittin theorem we may write
\[
\widehat K_\infty\cong\bigoplus_{a\in\C}E^{av}\otimes K^{(a)},
\]
where each $K^{(a)}$ has slopes $<1$, and only finitely many $K^{(a)}$ are nonzero.
\begin{enumerate}
\item[(F6)] (Stationary phase at finite distance \cite[Cor.~2.3(1)]{Ari08}, \cite[Prop.~2.1]{GL04}.) For $x\in\A^1_\tau$, $\phi_x(\Four(K))\cong\Four^{(\infty,0)}(K^{(x)})$, where $\Four^{(\infty,0)}$ is the local Fourier transform from $\infty$ to $0$, transported to $x$ by translation. In particular $\phi_x(\Four(K))=0$ if $K^{(x)}=0$.
\item[(F7)] (Properties of $\Four^{(\infty,0)}$ \cite{BE04,Sab08,GS10}.) Let $E$ be a formal germ at $\infty$ with slopes $<1$.
\begin{enumerate}
\item[(a)] $\rk\Four^{(\infty,0)}(E)=\rk E-\irr E$ \cite[Prop.~3.14]{BE04}. A germ of pure slope $q<1$ goes to a germ of pure slope $q/(1-q)$ \cite[Thm.~3.3]{GS10}; in particular regular germs go to regular germs.
\item[(b)] Let $c\in\C^\times$ and let $U$ be regular of rank one. For the rank two germ $A=[y=u^2]_+(E^{cu}\otimes U)$, the explicit formula \cite[(5.4)]{Sab08} gives
\[
\Four^{(\infty,0)}(A)\cong E^{c^2/(4s)}\otimes\mathcal R\qquad\text{with $\mathcal R$ regular of rank one.}
\]
\item[(c)] $\Four^{(\infty,0)}$ is compatible with homotheties. If $h_\sigma(v)=\sigma v$ with $\sigma\in\C^\times$, then
\[
\Four^{(\infty,0)}(h_\sigma^*E)\cong h_{1/\sigma}^*\Four^{(\infty,0)}(E).
\]
This follows from the identity $\Four(h_\sigma^*K)\cong h_{1/\sigma}^*\Four(K)$ of global Fourier transforms and the locality in (F6).
\item[(d)] $\Four^{(\infty,0)}$ of the trivial rank one germ is the rank one germ of vanishing cycles of $\delta_0$, with trivial monodromy. Indeed $\Four(\cO_{\A^1})=\delta_0$, and (F6) applies with $K^{(0)}$ trivial.
\item[(e)] On regular germs, $\Four^{(0,\infty)}$ is the identity after renaming the variable \cite[Rem.~5.2]{Sab08}, and $\Four^{(\infty,0)}$ is inverse to $\Four^{(0,\infty)}$ up to the pullback by $s\mapsto-s$ \cite[Rem.~5.3(4)]{Sab08}; compare \cite[Prop.~3.10]{BE04} on the subcategory without horizontal sections. Hence $\Four^{(\infty,0)}$ sends a regular germ to a regular germ of the same rank whose monodromy is conjugate to the original one, both computed in the local coordinates $1/v$ and $s$.
\end{enumerate}
\item[(F8)] (Vanishing cycles of minimal extensions \cite[\S3.3]{Ari08}.) Let $K$ be a holonomic module that, near $x$, has no nonzero submodule or quotient supported at $x$. If $\phi_x(K)$ is regular, then $K$ is regular singular at $x$. If $L$ denotes the local system of $K$ on a punctured disc around $x$, with local monodromy $T$, then $\phi_x(K)$ is the direct sum of the generalized eigenspaces of $T$ for eigenvalues $\ne1$ and of $\im(T-1)$ on the generalized eigenspace for the eigenvalue $1$, with the action of $T$. In particular $\dim\phi_x(K)=\rk(T-1)$, and $T$ acts trivially on $\phi_x(K)$ if and only if $T$ is unipotent with $(T-1)^2=0$. This extends (F3) of Part~\ref{part:order}.
\end{enumerate}
In (F7)(b), the exponential factor $c^2/(4s)$ is the value of the phase $c\sqrt y-sy$ at its critical point $\sqrt y=c/(2s)$. For $n=1$ one checks it directly: the Laplace transform of $I_0(\mu\sqrt y)$ is $s^{-1}e^{\mu^2/(4s)}$.

\section{Formal structure along the chain}\label{bs:sec:formal}

Fix $n\ge2$ and a mass configuration satisfying (G$'$); under (G) all $n_j=1$. Let $U_n$ be the rank one regular formal connection on $\C((u^{-1}))$ with solution $u^{-n/2}$, and for $c\in\C^\times$ put
\[
A_c(n)=[y=u^2]_+\bigl(E^{cu}\otimes U_n\bigr),
\]
an irreducible formal connection on $\C((y^{-1}))$ of rank two and pure slope $1/2$. It depends only on $c^2$ and $n$. For $\lambda\in\Lambda_m$ put $c(\lambda)=\sum_j\lambda_j\mu_j$, and let $\Lambda_m^+$ be a set of representatives of $\Lambda_m/\pm1$.

\begin{lemma}\label{bs:lem:Minf}
The formal completion of $M$ at $y=\infty$ is
\[
\widehat M_\infty\cong\bigoplus_{\lambda\in\Lambda_m^+}A_{c(\lambda)}(n)\ \oplus\ A_0,
\]
where $A_0$ is zero unless all $n_j$ are even, in which case it is regular of rank one with solution $y^{-n/4}$. Under {\rm(G$'$)} the $c(\lambda)^2$, $\lambda\in\Lambda_m^+$, are pairwise distinct and nonzero.
\end{lemma}

\begin{proof}
By the Hankel expansions of $I_0$ and $K_0$, the Bessel connection $B_\mu$ has the formal solutions $e^{\pm\mu u}u^{-1/2}(1+O(u^{-1}))$ at $u=\sqrt y=\infty$. Hence $\widehat{(B_\mu)}_\infty\cong A_\mu(1)$. The solutions of $\Sym^{n_j}B_{\mu_j}$ are spanned by the products of $n_j$ such solutions, and the weight $\lambda_j\in\{n_j,n_j-2,\dots,-n_j\}$ labels the exponential factor $e^{\lambda_j\mu_ju}$, each occurring once. Tensoring, $M$ has, for each tuple $\lambda=(\lambda_j)_j$ (including $\lambda=0$), exactly one formal solution $e^{c(\lambda)u}u^{-n/2}(1+O(u^{-1}))$. Under (G$'$), $c(\lambda)=c(\lambda')$ forces $\lambda=\lambda'$, so these exponential factors are distinct and $\widehat M_\infty$ splits accordingly. The pair $\{\lambda,-\lambda\}$ with $\lambda\ne0$ gives $A_{c(\lambda)}(n)$. The tuple $\lambda=0$ exists only if all $n_j$ are even, and gives a regular rank one summand with solution $u^{-n/2}=y^{-n/4}$.
\end{proof}

\begin{lemma}\label{bs:lem:F0}
Let $\cK_1=j_{!*}M$. The formal completion of the connection $\Four(\cK_1)|_{\Gm}$ at $s=0$ is
\[
\bigoplus_{\lambda\in\Lambda_m^+}E^{c(\lambda)^2/(4s)}\otimes \mathcal R_{c(\lambda)}\ \oplus\ \mathrm{Reg},
\]
where each $\mathcal R_c$ is regular of rank one, and $\mathrm{Reg}$ is regular of rank $R'-|\Lambda_m^+|$. The monodromy of $\mathcal R_c$ depends only on $n$, not on $c$ or $m$. Moreover $\phi_0(j_{!*}\mathrm{Reg})\cong\Four^{(\infty,0)}(A_0)$. In particular, if $A_0=0$, then $\mathrm{Reg}$ is trivial, and under {\rm(G)} $\mathrm{Reg}\cong\cO^{\nu_n}$.
\end{lemma}

\begin{proof}
All slopes of $\widehat{(\cK_1)}_\infty$ are $<1$, so (F6) at $x=0$ gives $\phi_0(\Four(\cK_1))\cong\Four^{(\infty,0)}(\widehat{(\cK_1)}_\infty)$. By (F7)(a),(b), $\Four^{(\infty,0)}(A_c(n))\cong E^{c^2/(4s)}\otimes \mathcal R_c$ with $\mathcal R_c$ regular of rank one. By (F7)(c), with $h_\sigma^*A_c(n)\cong A_{\sigma^{1/2}c}(n)$, we get $\mathcal R_{\sigma^{1/2}c}\cong h_{1/\sigma}^*\mathcal R_c$. Homotheties do not change the monodromy of a regular germ, so the monodromy of $\mathcal R_c$ is the same for all $c\ne0$. Since $A_c(n)$ depends only on $(c^2,n)$, this monodromy depends only on $n$. The exponential factors $c(\lambda)^2/(4s)$ are pairwise distinct, so these summands do not interact.

The module $\Four(\cK_1)$ is simple, so near $s=0$ it has no nonzero submodule or quotient supported at $0$. Its formal germ at $0$ is therefore the sum of its irregular part, which is the meromorphic germ $\bigoplus_\lambda E^{c(\lambda)^2/(4s)}\otimes \mathcal R_{c(\lambda)}$, and of $j_{!*}\mathrm{Reg}$, where $\mathrm{Reg}$ is the regular part of the meromorphic germ. The formal vanishing cycles split in the same way. Comparing regular parts with the decomposition of $\Four^{(\infty,0)}(\widehat{(\cK_1)}_\infty)$ above gives $\phi_0(j_{!*}\mathrm{Reg})\cong\Four^{(\infty,0)}(A_0)$. If $A_0=0$, then $\phi_0(j_{!*}\mathrm{Reg})=0$, so by (F8) $\mathrm{Reg}$ has trivial monodromy and is trivial. Under (G), $|\Lambda_m^+|=2^{n-1}$, and $R'=\sum_dd\,a_d=\dim W_m-\sum_da_d=2^n-\binom n{\lfloor n/2\rfloor}$. Here $\sum_da_d=\binom n{\lfloor n/2\rfloor}$ is the number of irreducible summands of $V^{\otimes n}$ (Part~\ref{part:order}, proof of Lemma~\ref{bo:lem:count}). Hence $\rk\mathrm{Reg}=2^{n-1}-\binom n{\lfloor n/2\rfloor}=\nu_n$.
\end{proof}

Since $\iota$ exchanges $s=0$ and $v=\infty$, we obtain:

\begin{lemma}\label{bs:lem:Nv}
The formal completion of $N_v$ at $v=\infty$ is $\bigoplus_{\lambda\in\Lambda_m^+}E^{c(\lambda)^2v/4}\otimes \mathcal R'_{c(\lambda)}\oplus\iota^*\mathrm{Reg}$, where $\mathcal R'_c=\iota^*\mathcal R_c$ is regular of rank one with monodromy depending only on $n$. Under {\rm(G)}, $\iota^*\mathrm{Reg}\cong\cO^{\nu_n}$.
\end{lemma}

\section{Local monodromy and the singular locus}\label{bs:sec:local}

For $c\in\C^\times$ put $\tau_c=c^2/4$; in the variable $t=4\tau$ these are the thresholds $c^2$.

\begin{proposition}\label{bs:prop:phi}
Assume {\rm(G$'$)}.
\begin{enumerate}
\item[(a)] $\Mhat$ is smooth on $\A^1_\tau\setminus(\{0\}\cup\{\tau_{c(\lambda)}:\lambda\in\Lambda_m^+\})$.
\item[(b)] For each $\lambda\in\Lambda_m^+$, $\dim\phi_{\tau_{c(\lambda)}}(\Mhat)=1$, and the monodromy acts on this line by a scalar $\rho_n$ that depends only on $n$.
\item[(c)] Under {\rm(G)}, $\dim\phi_0(\Mhat)=\nu_n$, and the monodromy acts trivially on $\phi_0(\Mhat)$.
\item[(d)] $\Mhat$ is regular singular at every point of $\PP^1$. Consequently $\cD\cdot\varpi_m$ is regular holonomic, and $L_m$ is Fuchsian.
\end{enumerate}
\end{proposition}

\begin{proof}
Apply (F6) to $K=j_{!*}N_v$, whose formal completion at $\infty$ is given by Lemma~\ref{bs:lem:Nv}; its slopes are $\le1$. We have $K^{(x)}=\mathcal R'_{c(\lambda)}$ if $x=\tau_{c(\lambda)}$, $K^{(0)}=\iota^*\mathrm{Reg}$, and $K^{(x)}=0$ otherwise. Since $\Mhat$ is simple of positive rank, it has no nonzero submodule or quotient supported at a point, so $\phi_x(\Mhat)=0$ means that $\Mhat$ is smooth at $x$. This gives (a). For (b), (F6) gives $\phi_{\tau_c}(\Mhat)\cong\Four^{(\infty,0)}(\mathcal R'_c)$, which by (F7)(a) has rank one. Its monodromy depends only on the formal germ $\mathcal R'_c$, hence only on $n$ by Lemma~\ref{bs:lem:Nv}. For (c), $\phi_0(\Mhat)\cong\Four^{(\infty,0)}(\cO^{\nu_n})$ has rank $\nu_n$ and trivial monodromy by (F7)(d). For (d), all $K^{(x)}$ are regular, so all $\phi_x(\Mhat)$ are regular by (F7)(a), and $\Mhat$ is regular singular at every finite point by (F8). At $\infty$ it is regular by (P2). The composition factors of $\cD\cdot\varpi_m$ are $\Mhat$ and copies of $\delta_0$ by (P3), all regular holonomic.
\end{proof}

\begin{lemma}\label{bs:lem:holo}
Let $p\in\C^\times$. Every solution of $L_m$ on a punctured disc around $p$ that is invariant under the local monodromy at $p$ extends holomorphically to $p$. In particular, if the local monodromy at $p$ is trivial, then all solutions of $L_m$ are holomorphic at $p$.
\end{lemma}

\begin{proof}
Let $g$ be such a solution. It is single-valued, and it is meromorphic at $p$ because $L_m$ is regular singular there (Proposition~\ref{bs:prop:phi}(d)). The function $g$ is annihilated by every $Q\in\C[t]\langle\partial_t\rangle$ with $Q\varpi_m=0$. Indeed, over $\C(t)$ the annihilator of $\varpi_m$ is generated by $L_m$, so $rQ=BL_m$ for some operator $B$ and some nonzero $r\in\C[t]$, and then $rQg=BL_mg=0$. Choose an affine open set $U\ni p$ with $0\notin U$, let $\mathrm{Mer}_p$ be the space of germs of meromorphic functions at $p$, and let $\cO_p\subset\mathrm{Mer}_p$ be the holomorphic germs; both are modules over $\cO(U)\langle\partial_t\rangle$. Then $P\varpi_m\mapsto Pg$ is a well-defined $\cO(U)\langle\partial_t\rangle$-linear map from $\cD\cdot\varpi_m|_U\cong\Mhat|_U$ (by (P3)) to $\mathrm{Mer}_p$. The module $\Mhat|_U$ is simple, being the restriction of a simple holonomic module to an open set on which it is nonzero. The image of this map in $\mathrm{Mer}_p/\cO_p$ is a quotient of $\Mhat|_U$ all of whose elements are killed by a power of $t-p$. Since $\Mhat|_U$ has positive rank, this image is zero, and so $g\in\cO_p$.
\end{proof}

\begin{proposition}\label{bs:prop:normal}
Let $n\ge2$, $m_1,\dots,m_n>0$ and $R=\sum_im_i$. As $v\to+\infty$,
\[
H(v)=C_1\,e^{R^2v/4}\,v^{-(n-1)/2}\bigl(1+O(v^{-1})\bigr),\qquad C_1=\sqrt\pi\,R\,(\pi R)^{-n/2}(m_1\cdots m_n)^{-1/2}.
\]
Consequently there is a function $A$, analytic at $R^2$, such that, as $t\to R^2$ from above, with $x=t-R^2$,
\[
\varpi_m(t)=\begin{cases}A(t)+\kappa_n\,x^{\frac{n-3}2}+O\bigl(x^{\frac{n-1}2}\bigr),& n\text{ even},\\[3pt] A(t)+\kappa_n\,x^{\frac{n-3}2}\log x+O\bigl(x^{\frac{n-1}2}|\log x|\bigr),& n\text{ odd},\end{cases}
\]
where
\begin{equation}\label{bs:eq:kappa}
\kappa_n=\frac{C_1}{4}\cdot4^{-\frac{n-3}2}\cdot\begin{cases}\Gamma\bigl(\frac{3-n}2\bigr),& n\text{ even},\\[2pt] (-1)^{\frac{n-1}2}\big/\bigl(\frac{n-3}2\bigr)!,& n\text{ odd}.\end{cases}
\end{equation}
\end{proposition}

\begin{proof}
For real $z\to+\infty$, $I_0(z)=e^z(2\pi z)^{-1/2}(1+O(z^{-1}))$. Hence
\[
G_m(y)=\kappa\,y^{-n/4}e^{R\sqrt y}\bigl(1+O(y^{-1/2})\bigr),\qquad\kappa=\prod_i(2\pi m_i)^{-1/2}.
\]
In $\widehat G(s)=\int_0^\infty e^{-sy}G_m(y)\,dy$ substitute $y=u^2$. The phase $-su^2+Ru$ has its maximum at $u_*=R/(2s)$, with value $R^2/(4s)$ and second derivative $-2s$. Laplace's method, in which the odd terms of the Gaussian expansion integrate to zero, gives
\[
\widehat G(s)=2\kappa\sqrt\pi\,(R/2)^{1-n/2}\,e^{R^2/(4s)}\,s^{(n-3)/2}\bigl(1+O(s)\bigr)\qquad(s\to0^+).
\]
The error coming from $O(y^{-1/2})$ is also $O(s)$, because $y\asymp s^{-2}$ near the maximum, and the contribution of $u$ away from $u_*$ is exponentially smaller. Since $H(v)=v^{-1}\widehat G(1/v)$, this gives the asymptotics of $H$, with $C_1=2\kappa\sqrt\pi(R/2)^{1-n/2}$, which simplifies to the stated value.

Put $\varepsilon=x/4$, $k=(n-1)/2$, and fix $v_0>0$ with $H(v)=C_1e^{R^2v/4}v^{-k}(1+r(v))$ and $|r(v)|\le C/v$ for $v\ge v_0$. Then
\[
\varpi_m(t)=\frac14\int_0^{v_0}e^{-tv/4}H(v)\,dv+\frac{C_1}4\int_{v_0}^\infty e^{-\varepsilon v}v^{-k}\,dv+\frac{C_1}4\int_{v_0}^\infty e^{-\varepsilon v}v^{-k}r(v)\,dv .
\]
The first term is entire in $t$. The second term equals $\frac{C_1}4\varepsilon^{k-1}\Gamma(1-k,\varepsilon v_0)$. If $k\notin\Z$, this is $\frac{C_1}4\Gamma(1-k)\varepsilon^{k-1}$ plus an entire function of $\varepsilon$. If $k\in\Z_{\ge1}$, it is $\frac{C_1}4\frac{(-1)^k}{(k-1)!}\varepsilon^{k-1}\log\varepsilon$ plus an entire function of $\varepsilon$, by the expansion of $\Gamma(1-k,z)$ in terms of the exponential integral. For the third term, let $j_0$ be the largest integer $<k$ and write $e^{-z}=\sum_{j\le j_0}(-z)^j/j!+\rho(z)$ with $|\rho(z)|\le C'\min(z^{j_0},z^{j_0+1})$ for $z\ge0$. The terms with $j\le j_0$ give a polynomial in $\varepsilon$, because $\int_{v_0}^\infty v^{j-k-1}\,dv<\infty$. Splitting the remaining integral at $v=1/\varepsilon$ bounds it by $O(\varepsilon^k)$ if $k\notin\Z$ and by $O(\varepsilon^k|\log\varepsilon|)$ if $k\in\Z$. Substituting $\varepsilon=x/4$, and absorbing $\log4$ into the analytic part, gives the claim with \eqref{bs:eq:kappa}.
\end{proof}

For $n=2$ the proposition gives $\varpi_m\sim\frac12(m_1m_2)^{-1/2}(t-R^2)^{-1/2}$, in agreement with $\varpi_m=((t-(m_1+m_2)^2)(t-(m_1-m_2)^2))^{-1/2}$.

\begin{proof}[Proof of Theorem~B]
(i) and (iii). By (P3), on every open set not containing $0$ the local system of $L_m$ is that of $\Mhat$. This also holds on a punctured disc around $0$. By Proposition~\ref{bs:prop:phi}(a) the local system extends over every point of $\C^\times$ other than the thresholds. By Proposition~\ref{bs:prop:phi}(b) and (F8), the local monodromy at each threshold satisfies $\rk(T-1)=1$, so it is nontrivial. By Proposition~\ref{bs:prop:phi}(c),(d) and (F8), $T_0$ is unipotent with $(T_0-1)^2=0$ and $\rk(T_0-1)=\nu_n$; this is (iii). It is nontrivial if and only if $\nu_n>0$, that is, $n\ge3$.

Let $p$ be a zero of the leading coefficient of $L_m$ that is neither a threshold nor $0$. The local monodromy at $p$ is trivial, so by Lemma~\ref{bs:lem:holo} all solutions are holomorphic at $p$, and $p$ is an apparent singularity. For $n=2$, $L_m=\Delta\partial_t+\frac12\Delta'$ with $\Delta=(t-\mu_+)(t-\mu_-)$, and $\Delta(0)=\mu_+\mu_-\ne0$ by (G), so $t=0$ is not a singular point.

(ii) and (iv). By Proposition~\ref{bs:prop:phi}(b) and (F8), the local monodromy $T$ at each threshold is a pseudo-reflection: if $\rho_n\ne1$ it is semisimple with eigenvalues $1,\dots,1,\rho_n$, and if $\rho_n=1$ it is a transvection. Since $\rho_n$ depends only on $n$, it suffices to compute it for one choice of masses. Take $m_1,\dots,m_n>0$ satisfying (G) (such masses exist for every $n$), and the normal threshold $R^2$, which corresponds to $\lambda=(1,\dots,1)$. Write $x=t-R^2$. By Lemma~\ref{bs:lem:holo}, the monodromy-invariant solutions near $R^2$ are holomorphic at $R^2$. Hence, on a slit disc around $R^2$,
\[
\varpi_m=\begin{cases}g+x^\beta g_1,& \rho_n\ne1,\\ g+g_1\log x,&\rho_n=1,\end{cases}
\]
where $g$ and $g_1$ are meromorphic at $R^2$, and $e^{2\pi i\beta}=\rho_n$ with $\beta\notin\Z$. If $\rho_n\ne1$, then $g$ is invariant, hence holomorphic. If $\rho_n=1$, then $g_1=(2\pi i)^{-1}(T-1)\varpi_m$ lies in $\im(T-1)$, hence is invariant and holomorphic. We compare with Proposition~\ref{bs:prop:normal}, using uniqueness of asymptotic expansions as $x\to0^+$ in the scale $x^e(\log x)^j$ with $e\in\C$, $j\in\N$. (A finite sum of convergent series in this scale that is $O(x^r|\log x|^N)$ has no terms with $\Re e<r$. Indeed, for a fixed real part the sum is $x^a$ times a polynomial in $\log x$ whose coefficients are almost periodic functions of $\log x$, and such an expression tends to $0$ only if it vanishes.)

Let $n$ be even. If $\rho_n=1$, then $(g-A)+g_1\log x=\kappa_nx^{(n-3)/2}+O(x^{(n-1)/2})$, where all exponents on the left are integers. Since $(n-3)/2\notin\Z$ and $\kappa_n\ne0$, this is impossible. Hence $\rho_n\ne1$ and $(g-A)+x^\beta g_1=\kappa_nx^{(n-3)/2}+O(x^{(n-1)/2})$. The exponents of $g-A$ are integers and those of $x^\beta g_1$ lie in $\beta+\Z$, so $\beta\equiv(n-3)/2\pmod\Z$ and $\rho_n=e^{\pi i(n-3)}=-1$. Moreover $x^\beta g_1-\kappa_nx^{(n-3)/2}$ has no terms of exponent $<(n-1)/2$, so $x^\beta g_1=x^{(n-3)/2}h$ with $h$ holomorphic and $h(R^2)=\kappa_n$. This is (iv), with $g$ in place of $A$.

Let $n$ be odd. If $\rho_n\ne1$, then $(g-A)+x^\beta g_1$ has no logarithmic terms, but it equals $\kappa_nx^{(n-3)/2}\log x$ up to an error $O(x^{(n-1)/2}|\log x|)$, which is impossible. Hence $\rho_n=1$, and $(g-A)+(g_1-\kappa_nx^{(n-3)/2})\log x=O(x^{(n-1)/2}|\log x|)$. Thus $g_1-\kappa_nx^{(n-3)/2}$ has no terms of exponent $<(n-1)/2$, so $g_1=x^{(n-3)/2}h$ with $h$ holomorphic and $h(R^2)=\kappa_n$, and $g-A$ has no terms of negative exponent, so $g$ is holomorphic. This is (iv), with $g$ in place of $A$.

In both cases $\rho_n=(-1)^{n+1}$, which proves (ii).
\end{proof}

\begin{remark}[Configurations]\label{bs:rem:config}
Under (G$'$), Proposition~\ref{bs:prop:phi}(a),(b),(d), Lemma~\ref{bs:lem:holo} and the proof of Theorem~B(ii) apply without change. Thus every nonzero value $c(\lambda)^2$, $\lambda\in\Lambda_m$, is a non-apparent singular point, at which the local monodromy is a reflection if $n$ is even and a transvection if $n$ is odd. At $t=0$ the answer depends on $A_0$. If some $n_j$ is odd, then $A_0=0$, and Theorem~B(iii) holds with $\nu_n$ replaced by $R'-|\Lambda_m^+|$. If all $n_j$ are even, then $\phi_0(j_{!*}\mathrm{Reg})\cong\Four^{(\infty,0)}(A_0)$ is a line on which the monodromy acts by $(-1)^{n/2}$, because $A_0$ has the solution $y^{-n/4}$. By (F8), the monodromy of $\mathrm{Reg}$ is then a reflection if $n\equiv2\pmod4$ and a transvection if $n\equiv0\pmod4$. By (F7)(e) and (F6), the monodromy on $\phi_0(\Mhat)\cong\Four^{(\infty,0)}(\iota^*\mathrm{Reg})$ is of the same type, and (F8) then shows the following. If $n\equiv2\pmod4$, $T_0$ has the eigenvalue $-1$ with multiplicity one, and its unipotent part $T_{0,u}$ satisfies $(T_{0,u}-1)^2=0$. If $n\equiv0\pmod4$, $T_0$ is unipotent with exactly one Jordan block of size $3$, and all other blocks have size at most $2$. For equal masses this is consistent with Table~\ref{bs:tab:equal}: the indicial root $3/2$ at $t=0$ for $n=6$, the triple root for $n=4$, where $T_0$ is a single block of size $3$, and the roots $1^2,2^3,3^2$ for $n=8$, where the blocks have sizes $3,2,2$.
\end{remark}

\section{Self-duality and the differential Galois group}\label{bs:sec:galois}

\subsection{The invariant form}
Let $\varepsilon_\mu$ be the class of $1$ in the Bessel connection $B_\mu$. Then $B_\mu$ carries the alternating form with $\langle\varepsilon_\mu,\theta\varepsilon_\mu\rangle=1$. It is horizontal: since $\theta^2\varepsilon_\mu=\frac{\mu^2}4y\,\varepsilon_\mu$, we get $\theta\langle\varepsilon_\mu,\theta\varepsilon_\mu\rangle=\langle\theta\varepsilon_\mu,\theta\varepsilon_\mu\rangle+\frac{\mu^2}4y\langle\varepsilon_\mu,\varepsilon_\mu\rangle=0$, and the other pairs are checked in the same way. Dually, the $\theta$-Wronskian $f\,\theta g-g\,\theta f$ of two solutions is constant. On the factor $B_t$ of $\cN$ the form is also horizontal for $\nabla_t$, because $\nabla_t\varepsilon=t^{-1}\theta\varepsilon$ and $\nabla_t\theta\varepsilon=\frac y4\varepsilon$ (Part~\ref{part:order}, Lemma~\ref{bo:lem:N}(d)), while $\nabla_t$ acts trivially on the other factors. The form on $\Sym^dV$ induced by an alternating form on $V=\C^2$ has parity $(-1)^d$, where $+1$ means symmetric and $-1$ alternating. Taking symmetric and tensor powers, we obtain on $\cN=B_t\otimes\bigotimes_j\Sym^{n_j}B_{\mu_j}$ a nondegenerate $\Kt[y,y^{-1}]$-bilinear form, horizontal for $\partial_y$ and for $\nabla_t$, of parity $(-1)^{1+n}$.

\begin{proposition}\label{bs:prop:dual}
$H^1_{\rm mid}(\cN)$ carries a nondegenerate $\Kt$-bilinear form that is horizontal for the Gauss--Manin connection and has parity $(-1)^n$. Consequently, the differential Galois group of $L_m$ is contained in $\Ort(o(m))$ if $n$ is even and in $\Sp(o(m))$ if $n$ is odd.
\end{proposition}

\begin{proof}
The form identifies $\cN$ with its dual $\cN^\vee$, compatibly with $\partial_y$ and $\nabla_t$. Poincaré duality for the de Rham cohomology of connections on curves gives a perfect pairing $H^1_c(\Gm,\cN)\times H^1(\Gm,\cN^\vee)\to\Kt$; compare \cite[\S2, Cor.~2.15]{FSY22} for the same construction for symmetric powers of Kloosterman connections. Explicitly, a class in $H^1_c$ is represented by a $1$-form $\alpha$ with values in $\cN$ together with formal primitives $f_x$ of $\alpha$ at $x\in\{0,\infty\}$, and it pairs with $[\beta]$ to $\sum_{x}\Res_x\langle f_x,\beta\rangle$. The pairing induces a nondegenerate pairing on the image $H^1_{\rm mid}(\cN)$ of $H^1_c\to H^1$, which coincides with the kernel of $H^1\to H^1_{\rm loc,0}$ used in Part~\ref{part:order}, Section~\ref{bo:sec:upper}. For $a,b\in H^1_{\rm mid}$, represented by $\alpha_a,\alpha_b$ with formal primitives $f_{a,x},f_{b,x}$, we have $\langle a,b\rangle=\sum_x\Res_x\langle f_{a,x},\alpha_b\rangle$. Since $d\langle f_{a,x},f_{b,x}\rangle=\langle\alpha_a,f_{b,x}\rangle+\langle f_{a,x},\alpha_b\rangle$ has no residue, $\langle a,b\rangle=-(-1)^{1+n}\langle b,a\rangle$. So the parity is $(-1)^n$. Since the form on $\cN$ is $\nabla_t$-horizontal and residues commute with $\partial_t$, the pairing is horizontal. By (P5), $H^1_{\rm mid}(\cN)\cong\Kt\otimes\cD\cdot\varpi_m$ as $\Kt$-connections, so the solution space of $L_m$ carries an invariant nondegenerate form of parity $(-1)^n$.
\end{proof}

\subsection{Proof of Theorem C}
Let $\mathcal V$ be the solution space of $L_m$ near a base point, $\Gamma\subseteq\GL(\mathcal V)$ the monodromy group, and $G$ the differential Galois group.

\begin{lemma}\label{bs:lem:regular}
$G$ is the Zariski closure of $\Gamma$. The group $\Gamma$ is generated by $T_\infty$ and by local monodromies $T_c$ at the thresholds $c(\lambda)^2$, $\lambda\in\Lambda_m^+$, each of which is a pseudo-reflection.
\end{lemma}

\begin{proof}
By Proposition~\ref{bs:prop:phi}(d), $L_m$ is Fuchsian, so $G$ is the Zariski closure of $\Gamma$ by Schlesinger's theorem \cite{vdPS03}. By Theorem~B(i) and Remark~\ref{bs:rem:config}, the local system extends to $\PP^1$ minus $\{0,\infty\}$ and the thresholds $c(\lambda)^2$. The fundamental group of this complement is generated by loops around all but one of the punctures; we omit $0$. The local monodromies at the thresholds are pseudo-reflections by Theorem~B(ii) and Remark~\ref{bs:rem:config}.
\end{proof}

\begin{lemma}\label{bs:lem:lie}
Let $n\ge3$. Then $\mathcal V$ is irreducible under the identity component $G^0$.
\end{lemma}

\begin{proof}
By Theorem~A, $\mathcal V$ is irreducible under $G$, and $T_\infty$ is unipotent with a Jordan block of size $n-1\ge2$. A unipotent element of $G$ lies in $G^0$, since the Zariski closure of the group it generates is connected. We first note that $G^0$ is reductive: its unipotent radical $U$ is characteristic in $G^0$, hence normal in $G$. The nonzero fixed space $\mathcal V^U$ is $G$-invariant, so irreducibility gives $\mathcal V^U=\mathcal V$. Since the representation is faithful, $U=1$. Thus restriction to $G^0$ is completely reducible, as needed below.

If $n$ is odd, all $T_c$ are transvections. Then all generators of $\Gamma$ are unipotent, so $G=G^0$, and the claim follows.

Let $n$ be even, so that the $T_c$ are reflections. Decompose $\mathcal V|_{G^0}=\mathcal V_1\oplus\dots\oplus\mathcal V_r$ into isotypic components; $G$ permutes them transitively. Suppose $r\ge2$. If a reflection $T_c$ maps some $\mathcal V_i$ to $\mathcal V_{i'}\ne\mathcal V_i$, then $v\mapsto T_cv-v$ is injective on $\mathcal V_i$, so $\dim\mathcal V_i\le\rk(T_c-1)=1$. Then all $\mathcal V_i$ are lines on which $G^0$ acts by characters, so $G^0$ is diagonalizable, which contradicts $1\ne T_\infty\in G^0$. Hence every $T_c$ preserves every $\mathcal V_i$, and so does $T_\infty\in G^0$. By Lemma~\ref{bs:lem:regular} the group $\Gamma$, and hence $G$, preserves every $\mathcal V_i$, which contradicts transitivity. So $r=1$, and $\mathcal V|_{G^0}\cong W\otimes\C^e$ with $W$ irreducible.

Suppose $e\ge2$. If $\dim W=1$, then $G^0$ acts by scalars, which is again impossible. If $\dim W\ge2$, every $g\in G$ normalizes $G^0$, so $W$ twisted by conjugation with $g$ is isomorphic to $W$. Hence there is $X\in\GL(W)$ such that $(X^{-1}\otimes1)g$ commutes with $G^0$, and by Schur's lemma $g=X\otimes Y$ for some $Y\in\GL(\C^e)$. A reflection $X\otimes Y$ is semisimple, so $X$ and $Y$ are semisimple, with eigenvalues $a_1,\dots$ and $b_1,\dots$, say. Among the eigenvalues $a_ib_j$ exactly one equals $-1$ and the others equal $1$. If $a_1b_1=-1$, then $a_1b_1=(a_1b_2)(a_2b_1)/(a_2b_2)=1$, a contradiction. Hence $e=1$.
\end{proof}

\begin{proof}[Proof of Theorem~C]
By Proposition~\ref{bs:prop:dual} and Lemma~\ref{bs:lem:lie}, $G^0$ is a connected subgroup of $\Ort(\mathcal V)$ (for $n$ even) or of $\Sp(\mathcal V)$ (for $n$ odd) acting irreducibly on $\mathcal V$. Hence it is reductive. By Schur's lemma its center acts by scalars, and the only scalars in $\Ort(\mathcal V)$ and $\Sp(\mathcal V)$ are $\pm1$. So $\mathfrak g=\Lie(G^0)$ is semisimple, it acts irreducibly on $\mathcal V$, and it is normalized by the pseudo-reflections $T_c\in G$. By the theorem of Kazhdan--Margulis, Gabber and Beukers--Heckman \cite[Thm.~1.5]{Kat90}, \cite{BH89}, $\mathfrak g$ is one of $\mathfrak{sl}(\mathcal V)$, $\mathfrak{so}(\mathcal V)$ and $\mathfrak{sp}(\mathcal V)$ (for some nondegenerate form in the last two cases).

Since $\mathcal V$ is an irreducible $\mathfrak g$-module, Schur's lemma shows that the space of $\mathfrak g$-invariant bilinear forms on $\mathcal V$ has dimension at most one. Let $n$ be odd. Then $\mathfrak g$ preserves the alternating form of Proposition~\ref{bs:prop:dual}, so $\mathfrak g\ne\mathfrak{so}(\mathcal V)$, which preserves a symmetric form. If $\dim\mathcal V\ge3$, then also $\mathfrak g\ne\mathfrak{sl}(\mathcal V)$, which preserves no nonzero bilinear form. Hence $\mathfrak g=\mathfrak{sp}(\mathcal V)$, which also holds if $\dim\mathcal V=2$. So $G^0=\Sp(\mathcal V)$ and $G=\Sp(\mathcal V)$. Let $n$ be even. Then $\dim\mathcal V=o(m)\ge n-1\ge3$, and the same argument with the symmetric form gives $\mathfrak g=\mathfrak{so}(\mathcal V)$. So $G^0=\SO(\mathcal V)$, and $G=\Ort(\mathcal V)$ because $G\subseteq\Ort(\mathcal V)$ contains a reflection.
\end{proof}

\subsection{Algebraic relations}
\begin{proof}[Proof of Corollary~D]
Let $E$ be a Picard--Vessiot extension for $L_m$ containing $\varpi_m$, and put $F=\C(t)(\varpi_m,\dots,\varpi_m^{(o-1)})$. By the Galois correspondence, $F=E^H$ with $H$ the stabilizer of $\varpi_m\in\mathcal V$. Hence $\trdeg_{\C(t)}F=\dim G-\dim H=\dim(G\cdot\varpi_m)$ \cite{vdPS03}. For $G=\Sp(\mathcal V)$ the orbit of a nonzero vector is $\mathcal V\setminus\{0\}$, so $\trdeg F=o$. For $G=\Ort(\mathcal V)$ and $o\ge3$ the orbit is open in the quadric $\{v:B(v,v)=B(\varpi_m,\varpi_m)\}$, where $B$ is the invariant form, so $\trdeg F=o-1$.

In the cyclic basis $e_i=\nabla_t^i\omega$, $0\le i<o$, of $H^1_{\rm mid}(\cN)\cong\Kt\otimes\cD\cdot\varpi_m$, the form of Proposition~\ref{bs:prop:dual} has a symmetric Gram matrix $g\in\GL_o(\C(t))$. A solution $f\in\mathcal V$ corresponds to the horizontal linear form $e_i\mapsto f^{(i)}$, and the dual form on these linear forms has Gram matrix $g^{-1}$. Hence $B(\varpi_m,\varpi_m)=\vec\varpi^{\,T}g^{-1}\vec\varpi$ with $\vec\varpi=(\varpi_m^{(i)})_{i<o}$, a quadratic relation with constant value. The ideal of relations is a prime ideal of height one in a polynomial ring over $\C(t)$, hence principal, generated by an irreducible $Q$ that divides this quadric. If $Q$ had degree one, there would be $b,a_i\in\C(t)$, not all $a_i$ zero, with $\sum_ia_i\varpi_m^{(i)}=b$. For $b=0$ this contradicts the minimality of $L_m$. For $b\ne0$ the line $\C(t)\,b\subset\Kt\otimes\cD\cdot\varpi_m$ would be a submodule of rank one, which contradicts irreducibility, since $o\ge3$. So $Q$ is the quadric itself, which is therefore irreducible.
\end{proof}

\subsection{Diamond lattices}
\begin{proof}[Proof of Corollary~E]
The $d$-dimensional diamond lattice is bipartite. Its edges from a vertex of the first class point in the directions $e_0,\dots,e_d$, where $e_0+\dots+e_d=0$ and any $d$ of the $e_i$ are linearly independent, and its edges from a vertex of the second class point in the directions $-e_0,\dots,-e_d$. A closed walk of length $2k$ from a vertex of the first class is a sequence $(i_1,j_1,\dots,i_k,j_k)$ with $\sum_s(e_{i_s}-e_{j_s})=0$. Since the only linear relation among the $e_i$ is $\sum e_i=0$, this holds if and only if $(i_s)$ and $(j_s)$ have the same multiset of values. Hence $W_{2k}=\sum_{|a|=k}\binom{k}{a}^2=c_k(1,\dots,1)$ with $n=d+1$, and $P_d(x)=x^{-1}\varpi_{(1,\dots,1)}(1/x)$. By (P1) this is also $\int_0^\infty u\,I_0(\sqrt x\,u)^{d+1}K_0(u)\,du$, which is the representation \cite[(41)]{Gut10}, due to Glasser and Montaldi, with $x=z^2$. The substitution $t=1/x$ preserves $\C(t)=\C(x)$, and multiplying by $x^{-1}$ does not change the Picard--Vessiot extension. So the differential Galois group of $P_d$ is that of $L_m$ for equal masses $m_i=1$, $n=d+1$. Equal masses satisfy (G$'$), and $o(m)=d$ by Theorem~A. Theorem~C gives $\Ort(d)$ for $d$ odd ($n$ even) and $\Sp(d)$ for $d$ even ($n$ odd). The local monodromy of $\varpi$ at $t=\infty$, that is $x=0$, is a single unipotent Jordan block of size $d$ (Part~\ref{part:order}, Corollary~\ref{bo:cor:equal}), so $n_{\log}=d-1$.
\end{proof}

In the variable $z$ with $x=z^2$, which is often used for lattice Green functions, the Galois group becomes a subgroup of finite index with the same identity component, $\SO(d)$ or $\Sp(d)$.

\section{The full Feynman integral}\label{bs:sec:full}

In position space, the two-dimensional Euclidean propagator is $K_0(m|x|)/(2\pi)$, so the banana integral with external momentum $p$ is
\[
\int_{\mathbb R^2}e^{ip\cdot x}\prod_i\frac{K_0(m_i|x|)}{2\pi}\,d^2x=(2\pi)^{1-n}\int_0^\infty rJ_0(|p|r)\prod_iK_0(m_ir)\,dr=(2\pi)^{1-n}I_m(-|p|^2),
\]
with $y=r^2$. For real $t<R^2$ the integrand of \eqref{bs:eq:I} decays like $e^{(\sqrt{\max(t,0)}-R)\sqrt y}$, and $I_m$ is real-analytic on $(-\infty,R^2)$.

Fix a real $t<R^2$. Let $s_F$ denote the solution of $\cN$ attached to $I_0(\sqrt{ty})\otimes\bigotimes_jK_0(\mu_j\sqrt y)^{n_j}$, in the notation of Part~\ref{part:order}, Section~\ref{bo:sec:upper}. Since $I_0(\sqrt{ty})$, like $K_0(\sqrt{ty})$, depends only on $ty$ and solves $\theta^2f=\frac{ty}4f$, the computation in the proof of Part~\ref{part:order}, Lemma~\ref{bo:lem:N}(d) gives $s_F\circ\nabla_t=\partial_t\circ s_F$. In particular $I_m=\frac12\int_0^\infty s_F(\eta)\,dy$.

\begin{lemma}\label{bs:lem:S}
Let $\xi=\sum_j\alpha_j\nabla_t^j\eta=\partial_yF$, with $F\in\cN$, be a relation as in (P4), and let $L=\sum_j\alpha_j\partial_t^j$. Then $LI_m=S$ on $(-\infty,R^2)$ minus a finite set, where
\[
S(t)=-\frac12\lim_{\varepsilon\to0^+}s_F(F)(\varepsilon,t)
\]
exists and is a rational function of $t$.
\end{lemma}

\begin{proof}
Differentiation under the integral sign, justified by bounds locally uniform in $t$, gives $LI_m=\frac12\int_0^\infty s_F(\xi)\,dy$. Near $y=0$, $s_F(\xi)=O((1+|\log y|)^{n+1})$ because $\xi\in\cN_0$ (Part~\ref{part:order}, proof of Lemma~\ref{bo:lem:cycle}). As $y\to\infty$, $s_F(F)\to0$ because $F$ has Laurent polynomial coefficients and $s_F$ of each basis element decays exponentially. Hence $\int_\varepsilon^\infty s_F(\xi)\,dy=-s_F(F)(\varepsilon)$, and the limit exists. To see that it is rational, expand. The functions $I_0(\sqrt{ty})$ and $\theta I_0(\sqrt{ty})$ are power series in $y$ whose coefficients are polynomials in $t$. The functions $K_0(\mu\sqrt y)$ and $\theta K_0(\mu\sqrt y)$ are sums $\sum_ky^k(a_k+b_k\log y)$ with constant $a_k,b_k$. The coefficients of $F$ in the basis~\eqref{bo:eq:basis} of Part~\ref{part:order} lie in $\Kt[y,y^{-1}]$. Hence $s_F(F)(\varepsilon)=\sum_{k,\ell}r_{k,\ell}(t)\,\varepsilon^k(\log\varepsilon)^\ell$, with $k$ bounded below, $0\le\ell\le n$, and $r_{k,\ell}\in\C(t)$, the series converging for small $\varepsilon$. Since the limit exists for all $t$ in an interval, $r_{k,\ell}=0$ for $k<0$, and for $k=0$, $\ell>0$. Hence $S=-\frac12r_{0,0}\in\C(t)$.
\end{proof}

\begin{lemma}\label{bs:lem:asym}
As $t\to-\infty$,
\[
I_m(t)=\frac{n}{2^{n-1}}\,\frac{(\log(-t))^{n-1}}{-t}\bigl(1+O(1/\log(-t))\bigr).
\]
\end{lemma}

\begin{proof}
Put $a=-t>0$ and $G^K(y)=\prod_iK_0(m_i\sqrt y)$. By the Mellin--Parseval formula \cite[Ch.~4]{BH86},
\[
\int_0^\infty a^{w-1}I_m(-a)\,da=\tfrac12\,\mathcal M[G^K](1-w)\,\mathcal M[J_0(\sqrt\cdot)](w),\qquad\mathcal M[J_0(\sqrt\cdot)](w)=\frac{4^w\Gamma(w)}{\Gamma(1-w)},
\]
in a strip $0<\Re w<\delta$. Near $y=0$, $G^K(y)=\sum_{\ell=0}^na_\ell(\log y)^\ell+O(y|\log y|^n)$ with $a_n=(-\frac12)^n$. So $\mathcal M[G^K](1-w)$ continues meromorphically to $\Re w<2$, with a single pole at $w=1$, of order $n+1$ and leading term $a_n(-1)^nn!\,(1-w)^{-n-1}$. Here is a uniform vertical-line estimate that justifies the contour shift after continuation. Put $\Theta=y\partial_y$ and $Q(y)=\Theta^{n+1}G^K(y)$. The convergent logarithmic expansion at zero gives $Q(y)=O(y(1+|\log y|)^n)$, uniformly in every closed subsector of $|\arg y|<\pi$. At infinity $Q$ has exponential decay in $\sqrt{|y|}$ in these subsectors. Thus $\mathcal M[Q](s)$ is holomorphic for $\Re s>-1$. Integration by parts, initially for $\Re s>0$, gives
\[
\mathcal M[G^K](s)=(-s)^{-n-1}\mathcal M[Q](s).
\]
For any fixed $0<\vartheta<\pi$, rotate the Mellin integration ray to argument $\vartheta\operatorname{sgn}(\Im s)$. The two endpoint arcs vanish for $\Re s>-1$, and this gives $\mathcal M[Q](\sigma+i\tau)=O(e^{-\vartheta|\tau|})$, uniformly for $\sigma$ in a compact subinterval of $(-1,\infty)$. Consequently the meromorphic continuation of $\mathcal M[G^K]$ has the same exponential decay, up to a polynomial factor, throughout the closed strips needed below. Stirling's formula shows that $\Gamma(w)/\Gamma(1-w)$ has only polynomial growth on vertical lines. These estimates justify both Mellin inversion and shifting its contour from $0<\Re w<1/4$ to $\Re w=2-\delta'$, for any fixed $0<\delta'<1$. Since $1/\Gamma(1-w)=(1-w)(1+O(1-w))$, the product has a pole of order $n$ at $w=1$, with leading term $2a_n(-1)^nn!\,(1-w)^{-n}$. Shifting the contour across $w=1$ gives $I_m(-a)=a^{-1}P(\log a)+O(a^{-2+\delta'})$, with $P$ a polynomial of degree $n-1$ and leading coefficient $2a_n(-1)^nn!/(n-1)!=n/2^{n-1}$.
\end{proof}

\begin{proof}[Proof of Theorem~F]
In Lemma~\ref{bs:lem:S} we have $\alpha_{o(m)}\ne0$, since otherwise $\varpi_m$ would be annihilated by an operator of order $<o(m)$, contradicting Theorem~A. Hence every operator of order $o(m)$ annihilating $\varpi_m$ is a $\C(t)^\times$-multiple of the operator of Lemma~\ref{bs:lem:S}, and it suffices to treat that one. By Lemma~\ref{bs:lem:S}, $LI_m=S$ with $S\in\C(t)$. Suppose $S=0$. Then $I_m$ continues to a solution of $L$ near $t=\infty$. By Theorem~A every such solution has the form $\sum_{\ell=0}^{n-2}(\log t)^\ell f_\ell(t)$ with $f_\ell$ meromorphic at $\infty$, because the Jordan blocks there have size at most $n-1$. Along $t\to-\infty$ such a function is a finite sum of terms $c\,|t|^e(\log|t|)^\ell$ with $e\in\Z$ and $\ell\le n-2$, plus smaller terms, so it cannot be asymptotic to $c\,(\log(-t))^{n-1}/t$ with $c\ne0$. This contradicts Lemma~\ref{bs:lem:asym}. Hence $S\ne0$.

The operator $(\partial_t-S'/S)\circ L$ annihilates $I_m$, so $\ord I_m\le o(m)+1$. Conversely, $\cD\cdot I_m$ contains $\cD\cdot S$, of generic rank one, and $\cD\cdot I_m/\cD\cdot S$ is generated by the class of $I_m$, which is killed by $L$. Over $\C(t)$ this quotient is therefore a quotient of $\C(t)\otimes\cD/\cD L\cong\Kt\otimes\cD\cdot\varpi_m$, which is irreducible. It is nonzero, since otherwise $I_m\in\C(t)\cdot S=\C(t)$, contradicting Lemma~\ref{bs:lem:asym}. Hence the generic rank of $\cD\cdot I_m$ is $o(m)+1$.
\end{proof}

\subsection{The affine differential Galois group}
We first record the standard irreducible case of the unipotent-radical principle for inhomogeneous equations \cite[\S1.1, (N.3), and Thm.~1.1]{Ber01}.

\begin{lemma}[Irreducible affine extension]\label{bs:lem:affine}
Let $K$ be a differential field of characteristic zero with algebraically closed constant field $C$. Let $L\in K\langle\partial\rangle$ be irreducible of order $r\ge1$, let $s\in K^\times$, and suppose $Lf=s$ has no solution in $K$. Let $G$ act on the solution space $V$ of $L$, and put $A=(\partial-s'/s)L$. Then the differential Galois group of $A$ is $V\rtimes G$. If $P\subset E$ are compatible Picard--Vessiot fields for $L$ and $A$, and $f\in E$ satisfies $Lf=s$, then
\[
E=P(f,f',\ldots,f^{(r-1)})
\]
is purely transcendental of degree $r$ over $P$.
\end{lemma}

\begin{proof}
The solution space $W$ of $A$ fits into the exact sequence
\[
0\longrightarrow V\longrightarrow W
 \xrightarrow{\ f\mapsto Lf/s\ } C\longrightarrow0.
\]
Surjectivity follows from the dimension count $\dim W=r+1$ and $\dim V=r$. If $H$ is the Galois group of $A$, restriction to $P$ is a surjective morphism $H\to G$. Choosing a lift of $1\in C$ identifies $H$ with a closed subgroup of $V\rtimes G$, with matrices $\left(\begin{smallmatrix}g&b\\0&1\end{smallmatrix}\right)$. The kernel $U$ of $H\to G$ is a closed subgroup of the additive vector group $V$, hence a vector subspace in characteristic zero. Conjugation makes it $G$-invariant. Irreducibility of $L$ gives $U=0$ or $U=V$.

The faithful irreducible representation of $G$ implies that $G$ is reductive: its connected unipotent radical is normal in $G$, has a nonzero fixed vector, and hence fixes all of $V$. If $U=0$, the morphism $H\to G$ is an isomorphism in characteristic zero, so $H$ is reductive as well. Finite-dimensional rational representations of such a group are completely reducible (also when it is disconnected, since its component group is finite). Thus the displayed exact sequence splits as an $H$-representation. An $H$-fixed lift $h\in W$ of $1$ then belongs to the fixed field $E^H=K$ and satisfies $Lh=s$, contrary to the hypothesis. Therefore $U=V$. Surjectivity to $G$ now implies $H=V\rtimes G$: multiply any lift of $g$ by an appropriate translation to obtain $(g,0)$.

A basis of $V$, together with $f$, is a basis of $W$. The relation $Lf=s$ expresses $f^{(r)}$ and all higher derivatives using the first $r$ jets. Hence the Picard--Vessiot field is
\[
E=P(f,\ldots,f^{(r-1)}).
\] Its relative Galois group is $U=V$, of dimension $r$. The Picard--Vessiot dimension formula gives $\trdeg_P E=r$. Since the displayed field is generated by precisely $r$ elements over $P$, those elements are algebraically independent, proving the final assertion.
\end{proof}

\begin{proof}[Proof of Theorem~G(i)--(ii)]
Theorem~A gives the irreducibility of $L$ for every positive mass configuration. Theorem~F gives $LI_m=S\ne0$. There is no rational solution $h\in\C(t)$ of $Lh=S$: otherwise $I_m-h$ would be a homogeneous solution of $L$. At infinity it would have logarithmic degree at most $n-2$ by Theorem~A. Adding a rational function cannot produce the nonzero term of logarithmic degree $n-1$ in Lemma~\ref{bs:lem:asym}. Thus Lemma~\ref{bs:lem:affine} applies. The identifications with the symplectic and orthogonal groups under {\rm(G$'$)} are exactly Theorem~C.
\end{proof}

\begin{corollary}[All polynomial relations through the next derivative]\label{bs:cor:fullrelations}
For all positive masses, write $L=\sum_{j=0}^o a_j(t)\partial_t^j$ with $a_o\ne0$. The ideal of polynomial relations over $\C(t)$, and also over $P$, among $I_m,I_m',\ldots,I_m^{(o)}$ is generated by the single affine linear polynomial
\[
\sum_{j=0}^o a_j(t)X_j-S(t).
\]
Under {\rm(G$'$)} and $n\ge3$, consider the $2o$ functions
\[
\varpi_m,\ldots,\varpi_m^{(o-1)},\quad I_m,\ldots,I_m^{(o-1)}.
\]
They are algebraically independent over $\C(t)$ if $n$ is odd. If $n$ is even, their ideal of relations is generated by the quadratic relation involving only the maximal-cut jets from Corollary~D.
\end{corollary}

\begin{proof}
Eliminate $X_o$ using its nonzero coefficient $a_o$. A further relation would give a polynomial relation among $I_m,\ldots,I_m^{(o-1)}$ over $P$, contrary to Theorem~G(ii). For the second assertion, expand any proposed relation as a polynomial in the full-integral jets, with coefficients that are polynomials in the maximal-cut jets. Independence over $P$ forces every such coefficient to vanish. Apply Corollary~D to those coefficients.
\end{proof}

\subsection{The complete Jordan partition at infinity}
\begin{proof}[Proof of Theorem~G(iii)]
The equation $A$ is an extension of the regular singular equation $L$ by a trivial rank-one equation. Regular singularity is closed under extensions \cite[Ch.~3]{vdPS03}. In this instance it can also be seen directly: variation of constants expresses an inhomogeneous solution as a regular singular fundamental matrix times integrals of its inverse times a meromorphic column. On bounded sectors the integrands, their primitives, and hence the solutions have at most power-logarithmic growth. Thus $A$ is regular singular at infinity. Its solution-space exact sequence is
\[
0\longrightarrow V\longrightarrow W\xrightarrow{L/S}\C\longrightarrow0.
\]
Since monodromy on $V$ and on the quotient is unipotent, the monodromy $T$ on $W$ is unipotent. Put $N=\log T$ and $r=n-1$. Theorem~A says that $N|_V$ has a unique largest block of size $r$: the highest weight $n$ occurs once in $W_m=\bigotimes_j\Sym^{n_j}\C^2$, and every other highest weight is smaller. The asymptotic in Lemma~\ref{bs:lem:asym} and regular singularity imply $N^r I_m\ne0$; otherwise the Frobenius expansion of $I_m$ would have logarithmic degree at most $r-1$. On the other hand $NI_m\in V$ and $N^rV=0$, so $N^{r+1}I_m=0$.

Choose an invariant decomposition $V=V_r\oplus V_{<r}$, where $V_r$ is the unique length-$r$ Jordan block and all blocks of $V_{<r}$ have smaller length. Put $a=NI_m$. Since $N^{r-1}a\ne0$, its projection to $V_r$ is cyclic. Therefore
\[
U=\operatorname{span}_\C\{a,Na,\ldots,N^{r-1}a\}
\]
projects isomorphically onto $V_r$ and is the graph of an $\C[N]$-linear map $V_r\to V_{<r}$. In particular $V=U\oplus V_{<r}$. As $I_m$ projects to $1$ in $W/V$, this gives
\[
W=\operatorname{span}_\C\{I_m,NI_m,\ldots,N^rI_m\}\oplus V_{<r}.
\]
The first summand is a single block of size $r+1=n$, and the other blocks are unchanged. This also proves that the extension is already nonsplit in the local category at infinity.
\end{proof}

\begin{example}[$n=2$]\label{bs:ex:n2}
Let $m_1,m_2>0$, $\mu_\pm=(m_1\pm m_2)^2$ and $\Delta=(t-\mu_+)(t-\mu_-)$. Then $\varpi_m=\Delta^{-1/2}$ is annihilated by $L=\Delta\partial_t+\frac12\Delta'$. For $t=-p^2<0$ the classical evaluation of the bubble integral gives
\[
I_m(t)=\int_0^\infty rJ_0(pr)K_0(m_1r)K_0(m_2r)\,dr=\Delta^{-1/2}\log\frac{m_1^2+m_2^2-t+\sqrt\Delta}{2m_1m_2},
\]
and both sides are analytic on $(-\infty,\mu_-)$.
With $u=m_1^2+m_2^2-t$ one has $\Delta=u^2-4m_1^2m_2^2$, and the logarithm has $t$-derivative $-\Delta^{-1/2}$. Hence $LI_m=-1$. The source reports checking the closed form against numerical quadrature of \eqref{bs:eq:I} for $m=(1,3)$ to $30$ digits.
\end{example}

\section{Source-reported computations}\label{bs:sec:numerics}
The numerical data in this section are retained from the supplied second
banana manuscript. Their code and full certificates were not supplied and
these runs were not reproduced for this compilation. They are not premises
of any proof. Finite-field rank searches or numerical fits by themselves
would not certify a characteristic-zero minimal operator or an exact
inhomogeneous identity.

The source reports that all modular computations use the prime $p=2^{61}-1$ and python-flint. For each case it reports computing a candidate minimal operator of $f(z)=\sum_kc_k(m)z^k=t\varpi_m(t)$, $z=1/t$, as $\sum_{r,d}c_{r,d}z^d\partial_z^r$, of the order predicted by Theorem~A and of minimal degree in $z$. It then reports analysing its leading coefficient $q(z)$, in which a threshold $\mu$ appears as the root $z=1/\mu$.

\begin{table}[ht]
\centering
\caption{Generic masses $m_i$ (not squared). The column $\gamma$ gives the indicial root at each threshold other than $0,1,\dots,\ord-2$. At $t=0$ we list the indicial roots of $f$, with multiplicity, and the number of logarithm-free solutions, which equals the number of Jordan blocks of $T_0$. The last column gives the degree of the remaining factor of $q$, with the value of \cite{LV23} in parentheses.}\label{bs:tab:generic}
\begin{tabular}{@{}cclcccc@{}}
\toprule
$n$ & masses & order, $\deg_z$ & thresholds & $\gamma$ & $t=0$: roots; log-free & apparent\\
\midrule
$3$ & $10,38,56$ & $2$, $7$ & $4$ simple & $0$ & $1^2$; $1$ & $2$ ($2$)\\
$4$ & $5,7,56,57$ & $6$, $27$ & $8$ simple & $1/2$ & $1^2,2^2,3,4$; $4$ & $17$ ($17$)\\
$5$ & $10,17,25,36,39$ & $12$, $121$ & $16$ simple & $1$ & $1^2,\dots,6^2$; $6$ & $98$ ($98$)\\
\bottomrule
\end{tabular}
\end{table}

\begin{itemize}
\item \emph{Generic masses} (Table~\ref{bs:tab:generic}). Every threshold is a simple root of $q$, with indicial roots $0,1,\dots,o_n-2$ and $\gamma=(n-3)/2$, the exponent of Theorem~B(iv). At $t=0$ the number of logarithm-free solutions is $o_n-\nu_n=1,4,6$. Together with $o_n=2,6,12$ this is consistent with exactly $\nu_n=1,2,6$ blocks of size $2$, as in Theorem~B(iii). The degrees $2,17,98$ of the remaining factor of $q$ coincide with those found in \cite{LV23}.
\item \emph{Equal masses} $m_i=1$, $3\le n\le8$ (Table~\ref{bs:tab:equal}). The nonzero thresholds $(n-2j)^2$ are simple roots with extra indicial root $(n-3)/2$, as in Remark~\ref{bs:rem:config}.
\item \emph{The inhomogeneous term.} Let $L_t$ be the operator in $t$ defined by $L_tF=L_z(tF)$, where $L_z$ is the primitive integral operator in $z$ and $z=1/t$; then $L_t\varpi_m=0$. The source reports computing $I_m$ and its $t$-derivatives by numerical quadrature at $30$ digits, using $\partial_t^kI_0(\sqrt{ty})=2^{-k}(y/t)^{k/2}I_k(\sqrt{ty})$. For $m=(1,3)$ it reports finding $L_tI_m=t$, to all digits, at the six points $t=-7$, $-2$, $-\frac13$, $\frac12$, $\frac{16}3$, $\frac{32}3$. This agrees with Example~\ref{bs:ex:n2}, since here $L_t=-tL$. For $m=(1,2,4)$ it reports evaluating $t^7L_tI_m$ at $16$ points and fitting a polynomial of degree $8$, with residual below $10^{-27}$. Its coefficients of $t^0,\dots,t^3$ vanish to working precision, and PSLQ at $22$ digits identifies the others, giving
\[
L_tI_m=\tfrac32t+(-42+15\ell)+\frac{189-345\ell}{t}+\frac{1470+6405\ell}{t^2}+\frac{\frac{3675}2-20475\ell}{t^3},\qquad\ell=\log2.
\]
This fitted expression is consistent with Theorem~F, but the numerical fit alone is not an exact certification of its coefficients.
\item \emph{Mass configurations for $n=6$.} The orders computed in \cite[\S4.5]{LV23} for the ten configurations $[1^6],[1^42],[1^33],[1^22^2],[1^24],[123],[15],[24],[3^2],[6]$ are $29,23,17,18,12,13,8,9,9,5$. These are exactly the values $o(m)$ of Theorem~A.
\end{itemize}

\begin{table}[ht]
\centering
\caption{Equal masses $m_i=1$: order $n-1$, nonzero thresholds $(n-2j)^2$, the extra indicial root $\gamma$ at each of them, and the indicial roots of $f=t\varpi$ at $t=0$.}\label{bs:tab:equal}
\begin{tabular}{@{}cccll@{}}
\toprule
$n$ & order & thresholds & $\gamma$ & roots at $t=0$\\
\midrule
$3$ & $2$ & $1,9$ & $0$ & $1^2$\\
$4$ & $3$ & $4,16$ & $1/2$ & $1^3$\\
$5$ & $4$ & $1,9,25$ & $1$ & $1^2,2^2$\\
$6$ & $5$ & $4,16,36$ & $3/2$ & $1^2,2^2,\frac32$\\
$7$ & $6$ & $1,9,25,49$ & $2$ & $1^2,2^2,3^2$\\
$8$ & $7$ & $4,16,36,64$ & $5/2$ & $1^2,2^3,3^2$\\
\bottomrule
\end{tabular}
\end{table}

\section{Remarks}\label{bs:sec:remarks}

\begin{proposition}[Exact local multiplicity from the invariant jets]\label{bs:prop:jetdefect}
Let $L$ be a rational-coefficient Fuchsian scalar operator of order $r\ge1$ at $t=\mu$. Suppose its local monodromy is a semisimple pseudo-reflection with exceptional eigenvalue different from $1$, and all monodromy-invariant solutions are holomorphic at $\mu$. Let
\[
0\le\alpha_0<\cdots<\alpha_{r-2}
\]
be the vanishing sequence of that $(r-1)$-dimensional space, and put
\[
s=\min\bigl(\mathbb Z_{\ge0}\setminus\{\alpha_0,\ldots,\alpha_{r-2}\}\bigr).
\]
For $r=1$ the sequence is empty and $s=0$. The multiplicity of $\mu$ in the leading coefficient of the primitive polynomial-coefficient normalization of $L$ is exactly $r-s$. In particular it is one if and only if $\alpha_j=j$ for every $0\le j\le r-2$.
\end{proposition}

\begin{proof}
Write $x=t-\mu$ and divide $L$ by its leading coefficient:
\[
L=\partial_x^r+\sum_{j=0}^{r-1}a_j(x)\partial_x^j.
\]
Fuchsianity gives $\operatorname{poleord}_0(a_j)\le r-j$. The primitive polynomial leading coefficient therefore has multiplicity equal to the maximum of these pole orders and zero. Frobenius theory and semisimplicity give a basis consisting of holomorphic solutions with valuations $\alpha_j$ and one solution $x^\beta u(x)$, where $u(0)\ne0$ is holomorphic and $\beta\notin\mathbb Z$. The indicial polynomial consequently has precisely the $r$ distinct roots $\alpha_0,\ldots,\alpha_{r-2},\beta$.

For $0\le j<s$, choose a holomorphic solution of valuation $j$. Their matrix of derivatives of orders $0,\ldots,s-1$ is invertible at zero. The equations $Lf=0$ for these $s$ solutions express $a_0,\ldots,a_{s-1}$ as holomorphic linear combinations of $1,a_s,\ldots,a_{r-1}$. Every coefficient on the latter list has pole order at most $r-s$, so this bounds all coefficients by $r-s$. When $s=0$ this is just the Fuchsian bound.

Set $A_j=\lim_{x\to0}x^{r-j}a_j(x)$ and use falling powers $(X)_{\underline j}=X(X-1)\cdots(X-j+1)$. The indicial polynomial is
\[
P(X)=(X)_{\underline r}+\sum_{j=0}^{r-1}A_j(X)_{\underline j}.
\]
The roots $0,\ldots,s-1$ force $A_0=\cdots=A_{s-1}=0$ by triangular evaluation. Since $s$ is not among the roots, $P(s)=s!A_s\ne0$. Thus $a_s$ has pole order exactly $r-s$, proving equality and the final characterization.
\end{proof}

Applied to every threshold of an even-$n$ banana equation under {\rm(G$'$)}, Proposition~\ref{bs:prop:jetdefect} turns the unresolved leading-coefficient multiplicity into an exact finite-jet question. It does not assert that the requisite jets are independent. Local monodromy and the exceptional exponent alone do not force this: the order-two equations with solution spaces $\langle1,x^{1/2}\rangle$ and $\langle x,x^{1/2}\rangle$ have the same reflection monodromy and exceptional exponent, but their primitive operators are $2x\partial_x^2+\partial_x$ and $2x^2\partial_x^2-x\partial_x+1$, with leading multiplicities one and two. A proof of the full Lairez--Vanhove leading-coefficient assertion therefore still needs additional control of the cyclic vector and its invariant jets. Neither the threshold simplicity statement for all $n$ nor the general multiplicity $\alpha(n)$ at zero is resolved here.

\begin{remark}[Simplicity of the thresholds]\label{bs:rem:simple}
Let $\mu$ be a threshold, $r=\ord L_m$, and let $\mathcal V^T$ be the space of solutions near $\mu$ that are invariant under the local monodromy. By Theorem~B(ii) and Lemma~\ref{bs:lem:holo}, $\dim\mathcal V^T=r-1$ and the elements of $\mathcal V^T$ are holomorphic at $\mu$. By the leading coefficient of $L_m$ we mean that of its primitive form in $\C[t]\langle\partial_t\rangle$. Suppose that their $(r-2)$-jets at $\mu$ are linearly independent. Then the monic operator $K$ of order $r-1$ with solution space $\mathcal V^T$ has coefficients holomorphic at $\mu$, since they are quotients of determinants by the Wronskian $w$ of a basis of $\mathcal V^T$, and $w(\mu)\ne0$. Write $L_m=a\,(\partial_t-b)K$ with $a\in\C(t)$. Then $b=W'/W-w'/w$, where $W$ is the Wronskian of all solutions. Since $W$ is multiplied by $\det T$ under monodromy and $L_m$ is Fuchsian, $W'/W$ has at most a simple pole at $\mu$, and so does $b$. Hence the coefficients of $a^{-1}L_m$ have at most simple poles at $\mu$, and $\mu$, being a non-apparent singular point, is a simple zero of the leading coefficient of $L_m$. For $n$ even the converse holds as well: at a simple zero the indicial roots are $0,1,\dots,r-2$ and one further root, which is not an integer because $T$ has the eigenvalue $-1$, so the holomorphic solutions have the exponents $0,\dots,r-2$. For $n=3$ the condition says that the invariant solution does not vanish at $\mu$, and at the normal threshold this follows from Theorem~B(iv), since $(T-1)\varpi_m=2\pi i\,\tilde h$ with $\tilde h(R^2)=\kappa_3\ne0$. We do not know a general argument for the other thresholds or for $n\ge4$.
\end{remark}

\begin{remark}[Picard--Lefschetz]
The maximal cut is a period of a family of Calabi--Yau varieties of dimension $n-2$ \cite{BKV15,BFKNS21,LV23}. If the family acquires an ordinary double point at a threshold, the Picard--Lefschetz formula predicts a reflection when $n-2$ is even and a transvection when $n-2$ is odd. Theorem~B(ii) establishes this parity for the minimal operator without using the geometry.
\end{remark}

\begin{remark}[Quadratic relations]
For $n$ even the quadratic relation of Corollary~D is a bilinear relation of the type studied for Bessel moments and Feynman integrals in \cite{FSY20,DPSS24}. Corollary~D says that, for the maximal cut and its derivatives of order below the minimal order, it is the only algebraic relation over $\C(t)$, and that for $n$ odd there is none.
\end{remark}

\begin{remark}[Exceptional groups]
For equal masses and $n=8$ the operator has order $7$ and is orthogonal. A priori its Galois group could have been of type $G_2$, as happens for some operators of order $7$. The threshold reflections exclude this, since the normalizer of $G_2$ in $\GL_7$ contains no reflections, and Theorem~C gives $\Ort(7)$.
\end{remark}


\clearpage
\part{Recurrence and arithmetic consequences}\label{part:transfers}
% Independently checked finite-dimensional representation theory.
% Its identification with differential order uses the banana order theorem;
% this proposition does not independently prove that analytic identification.
\section{Multiplicity partitions and strict increases under splitting}

Write $V_d=\operatorname{Sym}^d\mathbb C^2$. For a partition
$\boldsymbol n=(n_1,\ldots,n_k)$ of $n\ge2$, put
\[
 W_{\boldsymbol n}=\bigotimes_{j=1}^k V_{n_j}
   =\bigoplus_{d\ge0}a_dV_d,
 \qquad \mathfrak o(\boldsymbol n)=\sum_{d\ge2}(d-1)a_d.
\]
Define $P_{\boldsymbol n}(z)=\prod_j(1+z+\cdots+z^{n_j})
=\sum_h c_hz^h$ and $D=\prod_j(n_j+1)$.

\begin{proposition}\label{mc:prop:partition-order}
The representation-theoretic order functional is
\[
 \mathfrak o(\boldsymbol n)=
 \begin{cases}
 D-2c_{(n-1)/2},&n\text{ odd},\\
 D-c_{n/2}-c_{n/2-1},&n\text{ even}.
 \end{cases}
\]
Splitting any part $a+b$ into two positive parts $a,b$ never decreases
$\mathfrak o$. For $n\ge4$ every such split strictly increases it, by at
least $n-3$; for $n=2,3$ every split leaves it unchanged. Consequently,
when $n\ge4$ the unique minimizing partition is $(n)$ and the unique
maximizing partition is $(1,\ldots,1)$, with values
\[
 n-1\quad\hbox{and}\quad
 2^n-\binom{n+1}{\lfloor(n+1)/2\rfloor},
\]
respectively. For $n=2,3$ all partitions have value $n-1$.
\end{proposition}

\begin{proof}
The character of $W_{\boldsymbol n}$ is $z^{-n}P_{\boldsymbol n}(z^2)$.
Every irreducible summand has the parity of $n$. If $n$ is odd, each
summand has exactly one vector of weight $1$, so
$\sum_d a_d=c_{(n-1)/2}$ and $a_0=0$. If $n$ is even, each summand has
exactly one vector of weight $0$, and
$\sum_d a_d=c_{n/2}$, $a_0=c_{n/2}-c_{n/2-1}$. In both cases
\[
 D-\mathfrak o=2\sum_d a_d-a_0,
\]
which proves the coefficient formula.

Let $W'$ be the tensor product of the factors not being split. The
Clebsch--Gordan decomposition gives the actual direct-sum identity
\[
 V_a\otimes V_b\otimes W'
 \cong (V_{a+b}\otimes W')\ \oplus\!
 \bigoplus_{h=1}^{\min(a,b)}V_{a+b-2h}\otimes W'.
\]
The functional $\sum_{d\ge2}(d-1)a_d$ is additive and nonnegative on
representations, proving monotonicity. The $h=1$ summand has highest
weight $n-2$, of multiplicity one. For $n\ge4$ its contribution is at
least $(n-2)-1=n-3>0$. For $n=2,3$ all added summands have highest
weight at most $n-2\le1$ and contribute zero. Every partition is obtained
from $(n)$ by successive splits and refines to $(1,\ldots,1)$; strictness
therefore proves both uniqueness assertions. For $(n)$ the value is
$n-1$. For $(1,\ldots,1)$ use $P(z)=(1+z)^n$ and Pascal's identity in
the coefficient formula.
\end{proof}

\begin{remark}
When the banana order theorem identifies the minimal differential order
with $\mathfrak o(\boldsymbol n)$, Proposition~\ref{mc:prop:partition-order}
has the following precise interpretation. Separating one repeated squared
mass into two distinct squared-mass classes strictly raises the order for
$n\ge4$. The extremal classes are equality of all squared masses and
pairwise distinct squared masses. No assertion about differential order
follows from this representation-theoretic proposition without that
analytic identification.
\end{remark}

\section{Two different minimal sizes of banana recurrences}
\label{tr:rt-section}

For nonzero complex masses $m_1,\ldots,m_n$, put
\[
 c_k(m)=\sum_{a_1+\cdots+a_n=k}
       \binom{k}{a_1,\ldots,a_n}^{\!2}\prod_i m_i^{2a_i},
 \qquad F_m(z)=\sum_{k\ge0}c_k(m)z^k.
\]
Thus $F_m(z)=z^{-1}\varpi_m(1/z)$.  Write the distinct squared
masses as $\mu_1^2,\ldots,\mu_s^2$, with multiplicities
$n_1,\ldots,n_s$, so $\sum_jn_j=n$.  Define
\begin{equation}
\label{tr:rt-tau}
 Q=\prod_{j=1}^s(n_j+1),\qquad
 \delta=\begin{cases}1,&\text{every }n_j\text{ is even},\\0,&\text{otherwise},\end{cases}
 \qquad T=\frac{Q-\delta}{2}.
\end{equation}
Let $\mathcal T_m$ be the set of nonzero squared signed sums
\[
 \left(\sum_{j=1}^s\lambda_j\mu_j\right)^2,
 \qquad \lambda_j\in\{n_j,n_j-2,\ldots,-n_j\},\quad \lambda\ne0.
\]
The condition called \textup{(G$'$)} in Part~\ref{part:structure} means that these
numbers are nonzero and distinct after identifying $\lambda$ with
$-\lambda$.  Under this condition $|\mathcal T_m|=T$.

A polynomial recurrence is an identity, valid for all sufficiently
large integers $k$,
\begin{equation}
\label{tr:rt-recurrence}
 \sum_{j=0}^r p_j(k)c_{k+j}(m)=0,
 \qquad p_j\in\mathbb C[k],\quad p_0p_r\ne0.
\end{equation}
Its \emph{shift width} is $r$, whereas its \emph{coefficient degree}
is $D=\max_j\deg p_j$.  Rational-coefficient recurrences have the
same possible shift widths after clearing denominators.

The analytic input below is precisely the result established in
Parts~\ref{part:order} and~\ref{part:structure}: the minimal operator of $F_m$ is irreducible and
regular singular, has order
\begin{equation}
\label{tr:rt-ode-order}
 o(m)=\sum_{d\ge2}(d-1)a_d,
 \qquad \bigotimes_{j=1}^s\operatorname{Sym}^{n_j}\mathbb C^2
           =\bigoplus_d(\operatorname{Sym}^d\mathbb C^2)^{\oplus a_d},
\end{equation}
and, under \textup{(G$'$)}, has nonidentity local monodromy at each
$z=1/\mu$, $\mu\in\mathcal T_m$.  In Part~\ref{part:structure} these threshold
statements for multiplicities follow from Proposition~\ref{bs:prop:phi}
and Remark~\ref{bs:rem:config}; no definite-critical-point claim for
banana saddle points is needed here.  The upper bound proved next
is independent of this analytic input.

\begin{theorem}[Separate width and degree optima]
\label{tr:rt-main}
Let $n\ge2$.  For arbitrary nonzero masses, $c_k(m)$ has a
rational-coefficient recurrence of shift width at most $T$.
If \textup{(G$'$)} holds, then
\begin{equation}
\label{tr:rt-two-optima}
 \min r=T,\qquad \min D=o(m).
\end{equation}
The two minima range separately over all nonzero polynomial
recurrences \eqref{tr:rt-recurrence}.  The same minima hold if
coefficients are required to belong to
$K[k]$, where $K=\mathbb Q(m_1^2,\ldots,m_n^2)$.
In particular, for generic distinct squared masses
\[
 \min r=2^{n-1},\qquad
 \min D=2^n-\binom{n+1}{\lfloor(n+1)/2\rfloor},
\]
whereas for equal masses
\[
 \min r=\left\lfloor\frac{n+1}{2}\right\rfloor,
 \qquad \min D=n-1.
\]
These assertions do not claim that a single recurrence attains both
minima.
\end{theorem}

For example, when $n=2$, put $A=m_1^2+m_2^2$ and
$B=(m_1^2-m_2^2)^2$.  The generating function is
$(1-2Az+Bz^2)^{-1/2}$, and its coefficients satisfy
\[
 (k+1)c_{k+1}-A(2k+1)c_k+Bk\,c_{k-1}=0\qquad(k\ge1).
\]
The width is two and the coefficient degree is one for generic
masses; when the squared masses coincide, the width drops to one.


\subsection{An explicit parity system gives the upper bound}

Define the entire functions by their convergent series
\[
 I_0(x)=\sum_{a\ge0}\frac{(x/2)^{2a}}{(a!)^2},\qquad
 I_1(x)=\sum_{a\ge0}\frac{(x/2)^{2a+1}}{a!(a+1)!}.
\]
Direct differentiation gives
$I_0'(x)=I_1(x)$ and $I_1'(x)=I_0(x)-I_1(x)/x$.
For the index set
$\mathcal A=\prod_{j=1}^s\{0,\ldots,n_j\}$, put
\[
 u_\alpha(x)=\prod_{j=1}^s
       I_0(\mu_jx)^{n_j-\alpha_j}I_1(\mu_jx)^{\alpha_j},
 \qquad |\alpha|=\sum_j\alpha_j.
\]
The product rule proves the exact system
\begin{equation}
\label{tr:rt-bessel-system}
 (x\partial_x+|\alpha|)u_\alpha
 =x\sum_j\mu_j\bigl((n_j-\alpha_j)u_{\alpha+e_j}
                       +\alpha_j u_{\alpha-e_j}\bigr),
\end{equation}
with terms outside $\mathcal A$ omitted.  Let $E$ and $O$ be the
subsets of even and odd $|\alpha|$, respectively.  Their sizes are
$(Q+\delta)/2$ and $(Q-\delta)/2=T$, since
\[
 |E|-|O|=\prod_j\sum_{a=0}^{n_j}(-1)^a=\delta.
\]
Let $A$ be the $E\times O$ matrix and $B$ the $O\times E$ matrix
of the right side of \eqref{tr:rt-bessel-system}, with the factor $x$
removed.  Their entries are constants in
$K'=K(\mu_1,\ldots,\mu_s)$; they need not be transposes.

Parity gives vector expansions
\[
 (u_\alpha)_{\alpha\in E}=\sum_{k\ge0}e_kx^{2k},\qquad
 (u_\beta)_{\beta\in O}=\sum_{k\ge0}b_kx^{2k+1}.
\]
Set
\[
 D_E(k)=\operatorname{diag}_{\alpha\in E}(2k+|\alpha|),
 \qquad D_O(k)=\operatorname{diag}_{\beta\in O}(2k+1+|\beta|).
\]
Comparison of coefficients in \eqref{tr:rt-bessel-system} gives
\begin{equation}
\label{tr:rt-coefficient-system}
 D_E(k+1)e_{k+1}=A b_k,\qquad D_O(k)b_k=B e_k.
\end{equation}
All diagonal matrices inverted below are invertible for $k\ge0$.
Consequently the smaller, odd sector closes:
\[
 b_{k+1}=D_O(k+1)^{-1}B D_E(k+1)^{-1}A b_k.
\]
For $w_k=4^k(k!)^2b_k$, this becomes
\begin{equation}
\label{tr:rt-rational-system}
 w_{k+1}=R(k)w_k,\qquad
 R(k)=4(k+1)^2D_O(k+1)^{-1}B D_E(k+1)^{-1}A
       \in\operatorname{Mat}_T(K'(k)).
\end{equation}
Let $A_0$ be the row of $A$ indexed by $\alpha=0$.
Since
$u_0(x)=\sum_{k\ge0}c_k(m)x^{2k}/(4^k(k!)^2)$,
the empty coordinate of \eqref{tr:rt-coefficient-system} reads
\begin{equation}
\label{tr:rt-observation}
 c_{k+1}(m)=2(k+1)A_0w_k.
\end{equation}

\begin{lemma}[A finite rational system gives a scalar recurrence]
\label{tr:rt-system-lemma}
If $w_{k+1}=R(k)w_k$ in dimension $T$ and
$a_k=\ell(k)w_k$, with $R$ and $\ell$ rational over a field,
then $a_k$ satisfies a rational-coefficient recurrence of width at
most $T$.
\end{lemma}
\begin{proof}
For $0\le j\le T$, express $a_{k+j}=\ell_j(k)w_k$, where
$\ell_0=\ell$ and
$\ell_j(k)=\ell(k+j)R(k+j-1)\cdots R(k)$ for $j>0$.
The $T+1$ rows $\ell_j(k)$ lie in a $T$-dimensional vector space
over the rational-function field, so they are linearly dependent.
Their relation gives the claimed recurrence outside finitely many
denominator zeros.  Removing zero endpoints and shifting the index
cannot increase its width.  No invertibility or cyclicity assumption
on $R(k)$ is needed.
\end{proof}

Applying the lemma to \eqref{tr:rt-rational-system}--\eqref{tr:rt-observation}
proves the upper bound $r\le T$.  The extension $K'/K$ is finite.
Expanding the recurrence coefficients in a $K$-basis of $K'$ gives
at least one nonzero recurrence over $K(k)$, because each $c_k(m)$
belongs to $K$.  Clearing denominators completes the upper-bound
proof over the stated coefficient field.

\subsection{Singularities force width; irreducibility forces degree}

\begin{lemma}[Recurrence singularity budget]
\label{tr:rt-budget}
Suppose a convergent generating function $F(z)=\sum_{k\ge0}a_kz^k$
has a recurrence \eqref{tr:rt-recurrence} of width $r$ and coefficient
degree $D$.  With $\theta=z\partial_z$, it satisfies
\begin{equation}
\label{tr:rt-operator}
 LF=P(z),\qquad
 L=\sum_{j=0}^r z^{r-j}p_j(\theta-j),\qquad P\in\mathbb C[z].
\end{equation}
If $a_j=[k^D]p_j(k)$, the highest ordinary-derivative coefficient
of $L$ is
\begin{equation}
\label{tr:rt-leading-h}
 z^DH(z),\qquad H(z)=\sum_{j=0}^r a_jz^{r-j}\ne0,
 \qquad \deg H\le r.
\end{equation}
Every finite nonzero singularity of every continued branch of $F$
is a root of $H$.
\end{lemma}
\begin{proof}
Multiply the recurrence by $z^{k+r}$ and sum over indices where it
holds.  The omitted initial terms give the polynomial $P$.
The $j$th term is $z^{r-j}p_j(\theta-j)F$ up to a polynomial.
The coefficient of $\partial_z^D$ is \eqref{tr:rt-leading-h}; its
different powers of $z$ cannot cancel identically.  At a nonzero
point where $H$ is nonzero, \eqref{tr:rt-operator} is an ordinary
linear differential equation with analytic coefficients and forcing.
Every solution branch extends holomorphically there.  For $D=0$,
the same conclusion follows from $F=P/H$.
\end{proof}

This is the singularity-budget argument in the Franel
manuscript~\cite{KeioFranel}, reproduced here with its hypotheses explicit.
For positive masses $|c_k(m)|\le(\sum_i|m_i|)^{2k}$, and the same
bound follows by taking absolute values for complex masses, so
$F_m$ has positive radius of convergence.

Under \textup{(G$'$)}, every $1/\mu$ with $\mu\in\mathcal T_m$ is a
singularity of some continued branch of $F_m$.  To justify this
statement for the specific period, rather than merely for its
equation, use irreducibility and regular singularity of its minimal
operator.  Its monodromy representation is irreducible, and hence
the span of the monodromy orbit of any nonzero solution is the full
solution space.  If every continued branch of $F_m$ were
holomorphic at $1/\mu$, the local monodromy would fix that entire
span.  This contradicts the nonidentity threshold monodromy from
Part~\ref{part:structure}.  Lemma~\ref{tr:rt-budget} now gives $r\ge T$, proving the
first equality in \eqref{tr:rt-two-optima}.

\begin{lemma}[No nonzero rational image]
\label{tr:rt-no-rational-image}
Let $F$ be nonrational with irreducible minimal operator
$M\in\mathbb C(z)\langle\partial_z\rangle$ of order $o$.
If $LF\in\mathbb C(z)$ for a differential operator $L$ over
$\mathbb C(z)$, then $M$ right-divides $L$ and $LF=0$.
In particular, any nonzero such $L$ has order at least $o$.
\end{lemma}
\begin{proof}
The greatest common right divisor of $M$ and $L$ is either $M$
or $1$, since $M$ is irreducible.  In the latter case the Euclidean
algorithm gives operators $U,V$ with $UM+VL=1$.  Applying this
identity to $F$ expresses $F=V(LF)$ as a rational function, a
contradiction.  Therefore $M$ right-divides $L$ and $LF=0$.
\end{proof}

The threshold monodromy and the preceding orbit argument show
that $F_m$ is nonrational under \textup{(G$'$)}.  Applying
Lemma~\ref{tr:rt-no-rational-image} to \eqref{tr:rt-operator} gives
$D\ge o(m)$ and also shows that its polynomial forcing vanishes.
Conversely, write the minimal order-$o(m)$ operator using
$\partial_z^i=z^{-i}\theta(\theta-1)\cdots(\theta-i+1)$ and
clear denominators.  This gives an annihilator
$\sum_s z^s Q_s(\theta)$ with $\deg Q_s\le o(m)$.
Taking coefficients gives an eventual polynomial recurrence of
coefficient degree at most $o(m)$, proving the second equality.
It can also be obtained over $K[k]$: a relation with bounded width
and degree is a null vector for a matrix with entries in $K$;
extending the coefficient field does not change the rank of that
finite-dimensional system.  Although the system has infinitely
many rows, its row span is generated by finitely many rows.

\subsection{What minimum width says about a differential equation}

\begin{corollary}[Threshold polynomial and a higher-order equation]
\label{tr:rt-threshold-polynomial}
Assume \textup{(G$'$)}.  In every recurrence of minimum width $T$,
with coefficient degree $D$ and $a_j=[k^D]p_j(k)$, one has
\begin{equation}
\label{tr:rt-characteristic}
 \sum_{j=0}^T a_j\lambda^j
       =C\prod_{\mu\in\mathcal T_m}(\lambda-\mu),\qquad C\ne0.
\end{equation}
In particular both endpoint polynomials have degree $D$.
There exists a homogeneous polynomial differential equation for
$F_m$ whose highest derivative coefficient is
\begin{equation}
\label{tr:rt-clean-leading}
 C z^D\prod_{\mu\in\mathcal T_m}(1-\mu z).
\end{equation}
It is regular singular everywhere, has a simple leading-coefficient
zero at each reciprocal threshold, and has no other finite nonzero
singular point.  Its order $D$ is allowed to exceed $o(m)$.
\end{corollary}
\begin{proof}
The polynomial $H$ of Lemma~\ref{tr:rt-budget} has degree at most
$T$ and has all $T$ distinct nonzero numbers $1/\mu$ as roots.
Thus $H=C\prod_\mu(1-\mu z)$, and reversing coefficients gives
\eqref{tr:rt-characteristic}.  Its nonzero constant and leading
coefficients prove the endpoint assertion.
Choose a minimum-width recurrence and its operator $L$ in
\eqref{tr:rt-operator}.  Lemma~\ref{tr:rt-no-rational-image} gives $LF_m=0$,
and the leading coefficient is \eqref{tr:rt-clean-leading}.
Writing $L=\sum_{i=0}^D q_i(z)\partial_z^i$, the Euler-operator
form gives $z^i\mid q_i(z)$ and $\deg q_i\le T+i$.
These are the regular-singularity bounds at $0$ and $\infty$.
At each other leading-coefficient root the root is simple, so
the regular-singularity bounds hold there as well.
\end{proof}

\begin{corollary}[A strict obstruction to simultaneous optimality]
\label{tr:rt-tradeoff}
Assume \textup{(G$'$)}. If the primitive polynomial form of the minimal operator of $F_m$
has a finite nonzero apparent singular point, no recurrence can
simultaneously have width $T$ and coefficient degree $o(m)$.
More generally, if that leading coefficient has $A$ additional
distinct nonzero roots besides the reciprocal thresholds, every
recurrence of coefficient degree $o(m)$ has width at least $T+A$.
\end{corollary}
\begin{proof}
A degree-$o(m)$ recurrence gives an annihilator $L$ of order
$o(m)$, so $L$ is a rational multiple of the minimal operator $M$.
Normalize $M$ to polynomial coefficients with greatest common
divisor one.  Since $L$ also has polynomial coefficients, its
rational multiplier must be a polynomial: a denominator would
divide every coefficient of $M$.  Thus every nonzero root of the
leading coefficient of $M$ remains in the leading coefficient of
$L=z^DH$ at its highest derivative.  The degree bound
$\deg H\le r$ yields $r\ge T+A$.
\end{proof}

The clean higher-order equation in
Corollary~\ref{tr:rt-threshold-polynomial} therefore does not prove
the remaining simplicity assertion for the \emph{minimal}
Picard--Fuchs operator in the Lairez--Vanhove conjecture.  It is a
transfer result about recurrence width, coefficient degree, and
the cost of removing apparent points.  No separate assertion of
priority or of a new external conjecture closure is made.

% Self-contained transfer of the discrete-Abel integrality mechanism.
\section{An integral solution lattice for the two-mass coefficient recurrence}

\begin{theorem}\label{ar:thm:two-mass-lattice}
Let $a,b\in\mathbb Z$ with $a\ne b$, put $S=a+b$, $Q=(a-b)^2$, and set
\[
 U_n=\sum_{j=0}^n\binom nj^2a^jb^{n-j}\qquad(n\ge0).
\]
Consider all rational sequences satisfying, only for $n\ge1$,
\begin{equation}\label{ar:eq:two-mass-recurrence}
 (n+1)f_{n+1}=(2n+1)S f_n-nQ f_{n-1}.
\end{equation}
Their integral solution lattice is exactly $\mathbb ZU$; their
bounded-common-denominator subspace is exactly $\mathbb QU$. More
precisely, if $c=f_1-Sf_0\ne0$, then for every prime $p\nmid Q$ and every
$h\ge1$,
\[
 \min\{v_p(f_{p^h-1}),v_p(f_{p^h})\}\le v_p(c)-h.
\]
Thus boundedness at even one prime $p\nmid Q$ characterizes the line
$\mathbb QU$. For every finite set of primes $\mathcal P$, the solutions
in $\mathbb Z[\mathcal P^{-1}]$ are exactly
$\mathbb Z[\mathcal P^{-1}]U$.
\end{theorem}

\begin{proof}
A constant-term calculation, in the ring of formal power series, gives
\begin{align*}
 \sum_{n\ge0}U_nz^n
 &=\operatorname{CT}_x\frac1{1-z(S+ax+b/x)}\\
 &=\frac1{1-Sz}\sum_{j\ge0}\binom{2j}{j}
       \left(\frac{abz^2}{(1-Sz)^2}\right)^j
  =(1-2Sz+Qz^2)^{-1/2}.
\end{align*}
Indeed $U_n=\operatorname{CT}_x[(ax+b)(1+x^{-1})]^n$.
Writing this generating series as $F$, differentiation gives
$(1-2Sz+Qz^2)F'=(S-Qz)F$; coefficient extraction proves
\eqref{ar:eq:two-mass-recurrence} for $U$, with $U_0=1$, $U_1=S$.

For any other solution put $W_n=U_nf_{n+1}-U_{n+1}f_n$.
Subtracting the two recurrences yields
\[
 W_n=\frac{nQ}{n+1}W_{n-1}
     =\frac{Q^n}{n+1}(f_1-Sf_0).
\]
Since the $U_n$ are integers, at $n=p^h-1$ the valuation of this
determinant gives exactly the asserted inequality. If the sequence is
bounded in $\mathbb Q_p$ for one $p\nmid Q$, then $c=0$. The nonzero
coefficient $n+1$ ensures forward uniqueness, so $f=f_0U$.
An integral sequence is bounded at any prime $p\nmid Q$ and has
$f_0\in\mathbb Z$, proving the lattice assertion. Multiplication by a
common denominator proves the rational-line assertion. For localization,
choose a prime outside both $\mathcal P$ and the prime divisors of $Q$;
the same argument gives $f=f_0U$ with $f_0\in\mathbb Z[\mathcal P^{-1}]$.
All converses follow from integrality of $U$.
\end{proof}

\begin{remark}
For integer banana masses $m_1,m_2$ of distinct squares, take
$a=m_1^2$ and $b=m_2^2$. This is a transfer of the classical discrete
Abel identity and the prime-denominator mechanism, not a new order
formula for the banana differential equation. The exclusion $a=b$ is
essential: then $Q=0$, the recurrence imposed only for $n\ge1$ has a
free initial value $f_0$ invisible to all later terms. In particular the
sequence $(1,0,0,\ldots)$ is an integral solution and is not a multiple of
$U$. For $a=b\in\mathbb Z$, a second independent integral solution is
$f_0=0$, $f_n=\frac12\binom{2n}{n}a^{n-1}$ for $n\ge1$ (with its
polynomial interpretation at $a=0$); hence in that degenerate case the
integral initial-value lattice is all of $\mathbb Z^2$.
\end{remark}

\appendix
\section{Exact finite checks and reproducibility}\label{app:verification}
The proofs above do not use guessed recurrences or floating-point rank tests.
For this compilation, the following finite checks were run with exact integers
and rational numbers:
\begin{enumerate}
\item The central-coefficient formula for the differential order was compared
with independent repeated Clebsch--Gordan decomposition for all 506 partitions
of integers $2\le n\le14$. All 2059 elementary splits of those partitions
satisfied the stated strict-increase bound.
\item The three tableau recurrences and their Casoratians were checked against
counts obtained independently from the hook-length formula, through index 14.
\item The two-mass recurrence and its Casoratian were checked for all 56 ordered
pairs of distinct integers $a,b\in\{-3,-2,\ldots,4\}$, through index 24.
\item The Jacobi operator recursion reproduced all eight displayed coefficients
of the height-12 operator in exact rational arithmetic. This uses the finite
operator construction in Part~\ref{part:tableaux}, which supplies the complete
algorithm without an interpolation or guessing step.
\end{enumerate}
The following Python 3 program reproduces the first three checks using only
the standard library. The tests establish the reported finite identities;
they do not establish the all-parameter theorems, historical novelty, or a
formal verification of the differential-module arguments.
\begin{lstlisting}[language=Python]
"""Finite exact checks; none is used as an all-parameter proof."""
from fractions import Fraction as F
from math import factorial, comb, prod
from collections import Counter
from itertools import product

def partitions(n, largest=None):
    if n == 0:
        yield ()
        return
    for a in range(min(n, largest or n), 0, -1):
        for tail in partitions(n-a, a):
            yield (a,) + tail

def tensor_decomposition(parts):
    ans = Counter({0: 1})
    for a in parts:
        new = Counter()
        for d, multiplicity in ans.items():
            for h in range(min(d, a)+1):
                new[d+a-2*h] += multiplicity
        ans = new
    return ans

def rank_from_rep(parts):
    return sum((d-1)*v for d,v in tensor_decomposition(parts).items() if d>=2)

def rank_from_coeff(parts):
    p = [1]
    for a in parts:
        q = [0]*(len(p)+a)
        for i,v in enumerate(p):
            for j in range(a+1): q[i+j] += v
        p = q
    n = sum(parts)
    return prod(a+1 for a in parts) - (2*p[n//2] if n%2 else p[n//2]+p[n//2-1])

partition_count = split_count = 0
for n in range(2,15):
    for parts in partitions(n):
        partition_count += 1
        o = rank_from_rep(parts)
        assert o == rank_from_coeff(parts)
        for j,c in enumerate(parts):
            for a in range(1,c//2+1):
                split_count += 1
                new = parts[:j] + (a,c-a) + parts[j+1:]
                gap = rank_from_rep(new)-o
                assert gap >= n-3 if n>=4 else gap == 0

def hook_count(shape):
    hooks=[]
    for i,row in enumerate(shape):
        for j in range(row):
            hooks.append(row-j + sum(r>j for r in shape[i+1:]))
    return factorial(sum(shape))//prod(hooks)

def tableau_count(n,k):
    return sum(hook_count(p) for p in partitions(n) if len(p)<=k)

for k in (2,3,4):
    u=[tableau_count(n,k) for n in range(15)]
    f=[F(0),F(1)]
    for n in range(1,14):
        if k==2: a,b,c=n+2,2,4*n
        elif k==3: a,b,c=n+3,2*n+3,3*n
        else: a,b,c=(n+4)*(n+5),8*n+20,16*n*(n+1)
        assert a*u[n+1] == b*u[n]+c*u[n-1]
        f.append((b*f[n]+c*f[n-1])/a)
    for n in range(14):
        w=u[n]*f[n+1]-u[n+1]*f[n]
        if k==2: expected=F(2*(-4)**n,(n+1)*(n+2))
        elif k==3: expected=F(6*(-3)**n,(n+1)*(n+2)*(n+3))
        else: expected=F(2880*(-16)**n,(n+1)*(n+2)**2*(n+3)**2*(n+4)**2*(n+5))
        assert w == expected

pair_count=0
for a,b in product(range(-3,5),repeat=2):
    if a==b: continue
    pair_count+=1
    s,q=a+b,(a-b)**2
    u=[sum(comb(n,j)**2*a**j*b**(n-j) for j in range(n+1)) for n in range(25)]
    f=[F(0),F(1)]
    for n in range(1,24):
        assert (n+1)*u[n+1] == (2*n+1)*s*u[n]-n*q*u[n-1]
        f.append(((2*n+1)*s*f[n]-n*q*f[n-1])/(n+1))
    for n in range(24):
        assert u[n]*f[n+1]-u[n+1]*f[n] == F(q**n,n+1)

print('PASS: partitions',partition_count,'elementary splits',split_count)
print('PASS: hook-formula tableau recurrence/Casoratian checks, k=2,3,4; indices through 14')
print('PASS: two-mass recurrence/Casoratian checks for',pair_count,'integer pairs; indices through 24')
\end{lstlisting}

\clearpage
\begin{thebibliography}{99}
\bibitem{BFK95}
F.~Bergeron, L.~Favreau, and D.~Krob,
\emph{Conjectures on the enumeration of tableaux of bounded height},
Discrete Math.\ \textbf{139} (1995), 463--468.

\bibitem{BG00}
F.~Bergeron and F.~Gascon,
\emph{Counting Young tableaux of bounded height},
J.\ Integer Seq.\ \textbf{3} (2000), Article 00.1.7.

\bibitem{BCLSS07}
A.~Bostan, F.~Chyzak, G.~Lecerf, B.~Salvy, and \'E.~Schost,
\emph{Differential equations for algebraic functions},
in: Proceedings of ISSAC 2007, ACM, 2007; arXiv:cs/0703121.

\bibitem{BRS}
A.~Bostan, T.~Rivoal, and B.~Salvy,
\emph{Minimization of differential equations and algebraic values of $E$-functions},
arXiv:1906.05529.

\bibitem{DLMF}
F.~W.~J.~Olver, A.~B.~Olde Daalhuis, D.~W.~Lozier, B.~I.~Schneider,
R.~F.~Boisvert, C.~W.~Clark, B.~R.~Miller, B.~V.~Saunders,
H.~S.~Cohl, and M.~A.~McClain, eds.,
\emph{NIST Digital Library of Mathematical Functions},
\url{https://dlmf.nist.gov/10.40}, accessed 18 October 2026.

\bibitem{For10}
P.~J.~Forrester,
\emph{Log-gases and Random Matrices},
London Math.\ Soc.\ Monographs Ser.~34, Princeton University Press, 2010.

\bibitem{FK22}
P.~J.~Forrester and S.~Kumar,
\emph{Differential recurrences for the distribution of the trace of the
$\beta$-Jacobi ensemble},
Physica D \textbf{434} (2022), 133220.
\url{https://doi.org/10.1016/j.physd.2022.133220}.
Preprint arXiv:2011.00787v2 (1 December 2020), Corollary~2.5,
equations~(2.19)--(2.20).
\url{https://arxiv.org/abs/2011.00787v2}.

\bibitem{FW08}
P.~J.~Forrester and S.~O.~Warnaar,
\emph{The importance of the Selberg integral},
Bull.\ Amer.\ Math.\ Soc.\ (N.S.) \textbf{45} (2008), 489--534.

\bibitem{FW26}
P.~J.~Forrester and F.~Wei,
\emph{Higher-order linear differential equations for unitary matrix integrals: applications and generalisations},
SIGMA \textbf{22} (2026), 015; arXiv:2508.20797.

\bibitem{Ges90}
I.~M.~Gessel,
\emph{Symmetric functions and P-recursiveness},
J.\ Combin.\ Theory Ser.~A \textbf{53} (1990), 257--285.

\bibitem{Inc56}
E.~L.~Ince,
\emph{Ordinary Differential Equations},
Dover, New York, 1956.

\bibitem{KL05}
M.~Klazar and F.~Luca,
\emph{On integrality and periodicity of the Motzkin numbers},
Aequationes Math.\ \textbf{69} (2005), 68--75.
\url{https://doi.org/10.1007/s00010-004-2734-x}.

\bibitem{Mat25}
A.~Matveeva,
\emph{On the integrality of some P-recursive sequences},
arXiv:2511.02121v1 (2025), Sections~2--3, especially Theorem~3.4.
\url{https://arxiv.org/abs/2511.02121v1}.

\bibitem{Kol68}
E.~R.~Kolchin,
\emph{Algebraic groups and algebraic dependence},
Amer.\ J.\ Math.\ \textbf{90} (1968), 1151--1164.

\bibitem{Mis19}
M.~Mishna,
\emph{On standard Young tableaux of bounded height},
in: Recent Trends in Algebraic Combinatorics, Springer, Cham, 2019; arXiv:1805.08130.

\bibitem{vdPS03}
M.~van der Put and M.~F.~Singer,
\emph{Galois theory of linear differential equations},
Grundlehren Math. Wiss. 328, Springer, Berlin, 2003.

\bibitem{Was65}
W.~Wasow,
\emph{Asymptotic Expansions for Ordinary Differential Equations},
Interscience, New York, 1965.

\bibitem{Won89}
R.~Wong,
\emph{Asymptotic Approximations of Integrals},
Academic Press, Boston, 1989.

\bibitem{BE04}
S.~Bloch and H.~Esnault,
\emph{Local Fourier transforms and rigidity for $\mathcal D$-modules},
Asian J. Math. \textbf{8} (2004), no.~4, 587--605; arXiv:math/0312343.

\bibitem{BFKNS21}
K.~B\"onisch, F.~Fischbach, A.~Klemm, C.~Nega, and R.~Safari,
\emph{Analytic structure of all loop banana integrals},
J. High Energy Phys. \textbf{2021}, no.~5, 066; arXiv:2008.10574.

\bibitem{BKV15}
S.~Bloch, M.~Kerr, and P.~Vanhove,
\emph{A Feynman integral via higher normal functions},
Compos. Math. \textbf{151} (2015), no.~12, 2329--2375.

\bibitem{BS08}
J.~M.~Borwein and B.~Salvy,
\emph{A proof of a recursion for Bessel moments},
Experiment. Math. \textbf{17} (2008), 223--230.

\bibitem{BSWZ12}
J.~M.~Borwein, A.~Straub, J.~Wan, and W.~Zudilin,
\emph{Densities of short uniform random walks},
Canad. J. Math. \textbf{64} (2012), 961--990.

\bibitem{Duh25}
C.~Duhr,
\emph{Modular forms for three-loop banana integrals},
arXiv:2502.15325 (2025).

\bibitem{FSY22}
J.~Fres\'an, C.~Sabbah, and J.-D.~Yu,
\emph{Hodge theory of Kloosterman connections},
Duke Math. J. \textbf{171} (2022), no.~8, 1649--1747.

\bibitem{GL04}
R.~Garc\'{\i}a L\'opez,
\emph{Microlocalization and stationary phase},
Asian J. Math. \textbf{8} (2004), no.~4, 747--768; arXiv:math/0307366.

\bibitem{HTT08}
R.~Hotta, K.~Takeuchi, and T.~Tanisaki,
\emph{$D$-modules, perverse sheaves, and representation theory},
Progr. Math. 236, Birkh\"auser, Boston, 2008.

\bibitem{Kat90}
N.~M.~Katz,
\emph{Exponential sums and differential equations},
Ann. of Math. Stud. 124, Princeton Univ. Press, 1990.

\bibitem{LV23}
P.~Lairez and P.~Vanhove,
\emph{Algorithms for minimal Picard--Fuchs operators of Feynman integrals},
Lett. Math. Phys. \textbf{113} (2023), no.~2, Paper No.~37; arXiv:2209.10962.

\bibitem{Mal74}
B.~Malgrange,
\emph{Sur les points singuliers des \'equations diff\'erentielles},
Enseign. Math. (2) \textbf{20} (1974), 147--176.

\bibitem{Mal91}
B.~Malgrange,
\emph{\'Equations diff\'erentielles \`a coefficients polynomiaux},
Progr. Math. 96, Birkh\"auser, Boston, 1991.

\bibitem{Sab08}
C.~Sabbah,
\emph{An explicit stationary phase formula for the local formal Fourier--Laplace transform},
in: Singularities I, Contemp. Math. 474, Amer. Math. Soc., 2008, 309--330; arXiv:0706.3570. See especially Remark~5.7 on p.~326 for regular holonomic germs.

\bibitem{SZ23}
A.~Straub and W.~Zudilin,
\emph{Sums of powers of binomials, their Ap\'ery limits, and Franel's suspicions},
Int. Math. Res. Not. IMRN \textbf{2023}, no.~11, 9861--9879. arXiv:2112.09576.

\bibitem{Van14}
P.~Vanhove,
\emph{The physics and the mixed Hodge structure of Feynman integrals},
in: String-Math 2013, Proc. Sympos. Pure Math. 88, Amer. Math. Soc., 2014, 161--194.

\bibitem{Van26}
P.~Vanhove,
\emph{The multiloop sunset to all orders},
Ann. Henri Poincar\'e (2026), doi:10.1007/s00023-026-01734-8; arXiv:2603.03183.

\bibitem{Ari08}
D.~Arinkin,
\emph{Fourier transform and middle convolution for irregular $\mathcal D$-modules},
arXiv:0808.0699 (2008).

\bibitem{BDFKN22}
K.~B\"onisch, C.~Duhr, F.~Fischbach, A.~Klemm, and C.~Nega,
\emph{Feynman integrals in dimensional regularization and extensions of Calabi--Yau motives},
J. High Energy Phys. \textbf{2022}, no.~9, 156; arXiv:2108.05310.

\bibitem{Ber01}
D.~Bertrand,
\emph{Unipotent radicals of differential Galois group and integrals of solutions of inhomogeneous equations},
Math. Ann. \textbf{321} (2001), no.~3, 645--666; doi:10.1007/s002080100243.

\bibitem{BH86}
N.~Bleistein and R.~A.~Handelsman,
\emph{Asymptotic expansions of integrals},
Dover, New York, 1986.

\bibitem{BH89}
F.~Beukers and G.~Heckman,
\emph{Monodromy for the hypergeometric function ${}_nF_{n-1}$},
Invent. Math. \textbf{95} (1989), no.~2, 325--354.

\bibitem{DPSS24}
C.~Duhr, F.~Porkert, C.~Semper, and S.~F.~Stawinski,
\emph{Twisted Riemann bilinear relations and Feynman integrals},
J. High Energy Phys. \textbf{2025}, no.~3, 019; arXiv:2407.17175.

\bibitem{FSY20}
J.~Fres\'an, C.~Sabbah, and J.-D.~Yu,
\emph{Quadratic relations between Bessel moments},
Algebra Number Theory \textbf{17} (2023), no.~3, 541--602; arXiv:2006.02702.

\bibitem{GS10}
A.~Graham-Squire,
\emph{Calculation of local formal Fourier transforms},
arXiv:1011.3525 (2010).

\bibitem{Gut10}
A.~J.~Guttmann,
\emph{Lattice Green functions in all dimensions},
J. Phys. A \textbf{43} (2010), 305205; arXiv:1004.1435.

\bibitem{HMZ25}
S.~Hassani, J.-M.~Maillard, and N.~Zenine,
\emph{On the diagonals of rational functions: the minimal number of variables},
J. Phys. A \textbf{58} (2025), 255202; arXiv:2502.05543.

\bibitem{KeioFranel}
Keio,
\emph{Minimal Recurrence Orders and Integral Solution Lattices for Franel Sequences},
unpublished research manuscript supplied with this compilation, dated 5 September 2026.
Used as methodological motivation; the transferred results are proved here independently.

\end{thebibliography}
\end{document}
